
% new file names from date 20171119: 
% block5_171119.tex  

\renewcommand{\thisfile}{block5\_190328}  

\cleardoublepage 
\mysection{Perfect and Tilting DG Modules over NC DG Rings} 
\label{sec:perf-tilt-NC}
 
\AYcopyright

Perfect and tilting complexes over noncommutative rings are 
among the important applications of derived categories to ring theory. 
In this section of the book we study the more general concepts of 
perfect DG modules and tilting DG bimodules over noncommutative DG rings. 

Throughout the section we adhere to the following convention, that extends 
Convention \ref{conv:2490}. 

\begin{conv} \label{conv:3485}
There is a fixed a nonzero commutative base ring $\K$ (e.g.\ a field, or 
the ring of integers $\Z$). By default all linear objects and operations are 
$\K$-linear. This means that the DG rings are DG central $\K$-rings (Definition 
\ref{dfn:1071}), the DG modules are $\K$-linear, the DG bimodules 
are central over $\K$, and all homomorphisms between them are $\K$-linear. 
Likewise all linear categories and functors are $\K$-linear.
The symbol $\ot$ stands for $\ot_{\K}$. 
\end{conv}

All the results in this section specialize to rings: the category 
$\catt{Rng} \centover \K$ of {\em central $\K$-rings} (also known as 
associative unital $\K$-algebras) is a full subcategory of the category
$\catt{DGRng} \centover \K$ of DG central $\K$-rings. 

From subsection \ref{subsec:flat-dg-rng} onward we shall impose a flatness 
condition on our DG rings. 



\mysubsection{Algebraically Perfect DG Modules}
\label{subsec:perf-dgmods}
Here we define {\em algebraically perfect DG modules}, and prove several of 
their characterizations, in Theorems \ref{thm:3400}, \ref{thm:3415} and 
\ref{thm:3416}. See Remark \ref{rem:3491} regarding nomenclature. 

Semi-free filtrations of DG modules were introduced in Definition 
\ref{dfn:1575}. 

\begin{dfn} \label{dfn:3411}
Let $A$ be a DG ring and $P$ a DG $A$-module. 
\begin{enumerate}
\item Let $F = \bigl\{ F_j(P) \bigr\}_{j \geq -1}$ be a semi-free filtration of 
$P$. We say that $F$ has {\em finite extension length} if 
there is a number $j_1 \in \N$ such that $F_{j_1}(P) = P$.

\item The DG module $P$ is called {\em semi-free of finite extension length} 
if it admits some semi-free filtration $F$ that has finite extension length. 
\end{enumerate}
\end{dfn}

Recall that for a DG module $M$ and an integer $k$ we write  
$M[k] := \opn{T}^k(M)$, the $k$-th translation of $M$. See Definition 
\ref{dfn:2120}.

\begin{dfn} \label{dfn:3402}
Let $A$ be a DG ring and $P$ a DG $A$-module. 
\begin{enumerate}
\item We call $P$ a {\em finite free DG $A$-module} 
\index{Differential graded module! finite free}
if there is an isomorphism 
$P \cong \bigoplus_{i = 1}^r A[-k_i]$
in $\dcat{C}_{\mrm{str}}(A)$ for some $r \in \N$ and $k_i \in \Z$. 

\item A {\em finite semi-free filtration} 
\index{Filtration on a DG module! finite semi-free}
on $P$ is a semi-free filtration 
$F = \bigl\{ F_j(P) \bigr\}_{j \geq -1}$ that has these properties:
\begin{itemize}
\item Each $\opn{Gr}^F_j(P)$ is a finite free DG $A$-module.

\item The filtration $F$ has finite extension length (Definition 
\ref{dfn:3411}(1)). 
\end{itemize}

\item The DG module $P$ is called {\em finite semi-free} 
\index{Differential graded module! finite semi-free}
if it admits 
some finite semi-free filtration. 
\end{enumerate}
\end{dfn}

Saturated full triangulated subcategories, and those generated by a set of 
objects, were introduced in Definitions \ref{dfn:4915} and \ref{dfn:4916}.

\begin{dfn} \label{dfn:3401}
A DG $A$-module $L$ is called {\em algebraically perfect} 
\index{Differential graded module! algebraically perfect}
if $L$ belongs to the saturated full triangulated  subcategory of $\dcat{D}(A)$ 
generated by the DG module $A$. 
\end{dfn}

\begin{rem} \label{rem:3491}
The definition above is new. Earlier texts used the term {\em perfect}, and  
this referred, somewhat ambiguously, to several distinct properties of DG 
modules, or complexes, that are sometimes equivalent to each other. See Theorems 
\ref{thm:3400}, \ref{thm:3415} and \ref{thm:3416}. For commutative DG rings 
there is another notion -- that of a {\em geometrically perfect} DG module; see 
Remark \ref{rem:3425}. 
\end{rem}

The projective dimension of a DG module was introduced in Definition 
\ref{dfn:2128}.

\begin{prop} \label{prop:3415}
If $L \in \dcat{D}(A)$ is algebraically perfect, then it has finite projective 
dimension.
\end{prop}

\begin{proof}
The projective dimension of the DG module $A$ is zero. If $L$ has finite 
projective dimension then so do all its translates $L[k]$. 
If there's a distinguished triangle 
$L' \to L'' \to L \xar{\triangle}$
s.t.\ both $L'$ and $L''$ have finite projective dimension, then so does $L$. 
If $L$ is a direct summand of $L'$, and  $L'$ has finite projective 
dimension, then so does $L$.
Now use Proposition \ref{prop:4470}.
\end{proof}


Recall that the category $\dcat{D}(A)$ has arbitrary direct sums, and they are 
the direct sums in $\dcat{C}_{\mrm{str}}(A)$; see Theorem \ref{thm:3140} and 
Corollary \ref{cor:1665}. 

\begin{dfn} \label{dfn:3405}
A DG $A$-module $L$ is called a {\em compact object of $\dcat{D}(A)$}
\index{Differential graded module! compact}
if the functor  
\[ \opn{Hom}_{\dcat{D}(A)}(L, -) : \dcat{D}(A) \to \dcat{M}(\K) \]
commutes with infinite direct sums. 
\end{dfn}

Explanation: Let $\{ M_z \}_{z \in Z}$ be a collection of objects of 
$\dcat{D}(A)$, with direct sum 
$M := \bigoplus_{z \in Z} M_z$ and embeddings 
$e_z : M_z \to M$. The canonical homomorphism of $\K$-modules 
\begin{equation} \label{eqn:3480}
\begin{aligned}
& \Phi_L : 
\bigoplus_{z \in Z} \, \opn{Hom}_{\dcat{D}(A)}(L, M_z) \to
\opn{Hom}_{\dcat{D}(A)} (L, M) \ ,
\\[0.4em]
& \Phi_L := \bigl\{ \opn{Hom}(\opn{id}_L, e_z) \bigr\}_{z \in Z}
\end{aligned}
\end{equation}
is always injective. Compactness of $L$ says that $\Phi_L$ is bijective.  
Some texts use the name ``small'' instead of ``compact''; see Remark 
\ref{rem:3426}.  

\begin{prop} \label{prop:3545}
Let $A$ and $B$ be DG rings, $F : \dcat{D}(A) \to \dcat{D}(B)$ an 
equivalence of triangulated categories, and $L \in \dcat{D}(A)$. Then $L$ is a 
compact object of $\dcat{D}(A)$ if and only if $F(L)$ is a 
compact object of $\dcat{D}(B)$.  
\end{prop}

\begin{proof}
Suppose $\{ N_z \}_{z \in Z}$ is a collection of objects of $\dcat{D}(B)$. 
Because $F$ is an equivalence of categories, we can assume that 
$N_z = F(M_z)$ for some $M_z \in \dcat{D}(A)$.
The equivalence $F$ respects coproducts. Namely, if 
$M = \bigoplus_{z \in Z} M_z$ in $\dcat{D}(A)$, 
with embeddings $e_z : M_z \to M$, 
then the object $N := F(M) \in \dcat{D}(B)$, with the morphisms 
$F(e_z) : N_z \to N$,
is a coproduct of the collection of objects
$\{ N_z) \}_{z \in Z}$. 
There is a commutative diagram of $\K$-modules
\[  \UseTips  \xymatrix @C=10ex @R=6ex {
\bigoplus_{z \in Z} \opn{Hom}_{\dcat{D}(A)}(L, M_z) 
\ar[r]^(0.56){\Phi_L}
\ar[d]_{\bigoplus F}^{\cong}
&
\opn{Hom}_{\dcat{D}(A)}(L, M) 
\ar[d]^{F}_{\cong}
\\
\bigoplus_{z \in Z} \opn{Hom}_{\dcat{D}(B)} \bigl( F(L), N_z \bigr) 
\ar[r]^(0.56){\Phi_{F(L)}}
&
\opn{Hom}_{\dcat{D}(B)} \bigl( F(L), N \bigr) 
} \]
We see that if $\Phi_L$ is an isomorphism, then so is $\Phi_{F(L)}$.

For the reverse direction we do the same, but now for a quasi-inverse 
$G : \dcat{D}(B) \to \dcat{D}(A)$ of $F$, the object 
$F(L) \in \dcat{D}(B)$ and the object $L \cong G(F(L)) \in \dcat{D}(A)$.
\end{proof}

Here are a few lemmas about compact objects that will be needed for the proof 
of Theorem \ref{thm:3400}. 

\begin{lem} \label{lem:3405}
Let $L$ be a compact object of $\dcat{D}(A)$, let 
$\{ M_z \}_{z \in Z}$ be a collection of objects of $\dcat{D}(A)$, 
let $M := \bigoplus_{z \in Z} \, M_z$, 
and let $\phi : L \to M$
be a morphism in $\dcat{D}(A)$. Then there is a finite subset 
$Z_0 \sub Z$, such that writing 
$M_0 := \bigoplus_{z \in Z_0} \, M_z$, 
and denoting by $\ga_0 : M_0 \to M$ the inclusion, 
there exists a morphism 
$\phi_0 : L \to M_0$ in $\dcat{D}(A)$ satisfying 
$\phi = \opn{Q}(\ga_0) \circ \phi_0$. 
\end{lem}

\begin{proof}
We are given 
$\phi \in \opn{Hom}_{\dcat{D}(A)}(L, M)$.
Using the canonical isomorphism (\ref{eqn:3480})
that we get by the compactness of $L$, 
we can express $\phi$ as a finite sum, say  
$\phi = \Phi_L \bigl( \sum\nolimits_{z \in Z_0} \phi_z \bigr)$,
where $Z_0 \sub Z$ is a finite subset, and 
$\phi_z \in \opn{Hom}_{\dcat{D}(A)}(L, M_z)$.
The morphism 
\[ \phi_0 := \sum\nolimits_{z \in Z_0} \phi_z : L \to 
\bigoplus\nolimits_{z \in Z_0} \, M_z = M_0 \]
in $\dcat{D}(A)$ has the required property.
\end{proof}

\begin{rem} \label{rem:3405}
It might be helpful to say some words about the lemma. 
The inclusion $\ga_0 : M_0 \to M$ is a (split) monomorphism
in the abelian category $\dcat{C}_{\mrm{str}}(A)$. Suppose we choose a 
K-projective resolution $P \to L$. The morphisms $\phi : L \to M$
and $\phi_0 : L \to M_0$ in $\dcat{D}(A)$ are represented by 
homomorphisms 
$\til{\phi} : P \to M$ and $\til{\phi}_0 : P \to M_0$
in $\dcat{C}_{\mrm{str}}(A)$. 
There is no reason for the homomorphism $\til{\phi}$ to factor through $M_0$
in $\dcat{C}_{\mrm{str}}(A)$. 
All we can say is that {\em there is a homotopy}
$\til{\phi} \twoto \ga_0 \circ \til{\phi}_0$. 

Later, after proving Theorem \ref{thm:3400}, we will know that there is a 
finite semi-free DG $A$-module $Q$, with homomorphisms 
$P \xar{\al} Q \xar{\be} P$ in $\dcat{C}_{\mrm{str}}(A)$,
and with a homotopy 
$\be \circ \al \twoto \opn{id}_P$.  
The homomorphism 
$\til{\phi} \circ \be : Q \to M$
does factor through some finite direct sum $M_0$, and hence so does 
$\til{\phi} \circ \be \circ \al : P \to M$. 
The homotopy $\be \circ \al \twoto \opn{id}_P$
induces the homotopy $\til{\phi} \twoto \ga_0 \circ \til{\phi}_0$
that was mentioned above. 
\end{rem}

The next definition is taken from \cite{BoNe}. 

\begin{dfn} \label{dfn:3410}
Suppose 
$\bigl( \{ M_i \}_{i \in \N}, \{ \mu_i \}_{i \in \N} \bigr)$ 
is a direct system in $\dcat{D}(A)$. Let 
\[ \phi \, : \, \bigoplus\nolimits_{i \in \N} \, M_i \ \to \ 
\bigoplus\nolimits_{i \in \N} \, M_i \]
be the morphism in $\dcat{D}(A)$ defined by 
\[ \phi|_{M_i} := (\opn{id}, -\mu_i) : M_i \to M_i \oplus M_{i + 1} . \]
A {\em homotopy colimit} 
\index{Homotopy colimit}
of the direct system $\{ M_i \}_{i \in \N}$
is an object $M \in \dcat{D}(A)$ that is a cone on $\phi$, in the sense of 
Definition \ref{dfn:3600}. Namely $M$ sits in some distinguished triangle 
\[ \bigoplus\nolimits_{i \in \N} \, M_i \xar{\ \phi \ } 
\bigoplus\nolimits_{i \in \N} \, M_i 
\xar{\psi} M \xar{\ \triangle \ } \]
in $\dcat{D}(A)$.
\end{dfn}

Observe that a homotopy colimit exists, and it is unique up to a 
non\-unique isomorphism in $\dcat{D}(A)$. See Corollary \ref{cor:3600}.

\begin{lem} \label{lem:3610}
If 
$\bigl( \{ M_i \}_{i \in \N}, \{ \mu_i \}_{i \in \N} \bigr)$
is a direct system in $\dcat{C}_{\mrm{str}}(A)$, and we take 
$M := \lim_{i \to} M_i$,
the direct limit in $\dcat{C}_{\mrm{str}}(A)$, then $M$ is a 
homotopy colimit of the direct system 
$\bigl( \{ M_i \}_{i \in \N}, \{ \opn{Q}(\mu_i) \}_{i \in \N} \bigr)$ 
in $\dcat{D}(A)$
\end{lem}

\begin{proof}
Let 
\[ \til{\phi} \, : \, \bigoplus\nolimits_{i \in \N} \, M_i \ \to \ 
\bigoplus\nolimits_{i \in \N} \, M_i \]
be the morphism in $\dcat{C}_{\mrm{str}}(A)$ defined like in Definition 
\ref{dfn:3410} above; so the morphism $\phi$ in that definition equals 
$\opn{Q}(\til{\phi})$. 

An easy calculation shows that 
\[ 0 \to \bigoplus\nolimits_{i \in \N} \, M_i \xar{\ \til{\phi} \ } 
\bigoplus\nolimits_{i \in \N} \, M_i \xar{\ \til{\psi} \ } M \to 0 , \] 
where $\til{\psi}$ is the canonical homomorphism, is an exact sequence in 
$\dcat{C}_{\mrm{str}}(A)$. 
By Proposition \ref{prop:2165} there is a distinguished triangle 
\[ \bigoplus\nolimits_{i \in \N} \, M_i \xar{\ \opn{Q}(\til{\phi}) \ } 
\bigoplus\nolimits_{i \in \N} \, M_i \xar{\ \opn{Q}(\til{\psi}) \ } M 
\xar{ \, \triangle \, } \]
in $\dcat{D}(A)$. 
\end{proof}

\begin{lem} \label{lem:3410} 
Let $L$ be a compact object of $\dcat{D}(A)$, and 
let $\{ M_i \}_{i \in \N}$ be a direct system in $\dcat{D}(A)$, with homotopy 
colimit $M$. Then the morphism $\psi$ in Definition \tup{\ref{dfn:3410}}
induces an isomorphism
\[ \lim_{i \to} \, \opn{Hom}_{\dcat{D}(A)}(L, M_i) \iso 
\opn{Hom}_{\dcat{D}(A)}(L, M) \]
in $\dcat{M}(\K)$. 
\end{lem}

\begin{proof}
Applying the cohomological functor $\opn{Hom}_{\dcat{D}(A)}(L, -)$
to the distinguished triangle in Definition \ref{dfn:3410}, we get 
long exact sequence of $\K$-modules
\begin{equation}  \label{eqn:4790}
\begin{aligned}
&
\cdots \xar{\de^{k - 1}}
\opn{Hom}_{\dcat{D}(A)} \Bigl( L, 
\Bigl( \bigoplus\nolimits_{i \in \N} \, M_i \Bigr)[k] \Bigr)
\\
& \quad
\xar{ \opn{Hom}(L, \phi[k]) }
\opn{Hom}_{\dcat{D}(A)} \Bigl( L, 
\Bigl( \bigoplus\nolimits_{i \in \N} \, M_i \Bigr)[k] \Bigr)
\\
& \quad  
\xar{ \opn{Hom}(L, \psi[k]) }
\opn{Hom}_{\dcat{D}(A)} \bigl( L, M[k] \bigr) \xar{\de^{k}}  \cdots .
\end{aligned}
\end{equation}
Using the compactness of $L$ and the resulting isomorphism $\Phi_L$, we obtain 
from (\ref{eqn:4790}) the long exact sequence of $\K$-modules
\begin{equation}  \label{eqn:3411}
\begin{aligned}
&
\cdots \xar{\de^{k - 1}} 
\bigoplus\nolimits_{i \in \N} \, \opn{Hom}_{\dcat{D}(A)} 
\bigl( L, M_i[k] \bigr) 
\\
& \quad  
\xar{\phi^k} \bigoplus\nolimits_{i \in \N} \, 
\opn{Hom}_{\dcat{D}(A)} \bigl( L, M_i[k] \bigr) 
\\
& \quad  
\xar{\psi^k} \opn{Hom}_{\dcat{D}(A)} \bigl( L, M[k] \bigr)  \xar{\de^{k}}  
\cdots . 
\end{aligned}
\end{equation}
We know that for every $k$ the homomorphism $\phi^k$ above is 
injective. This implies that the connecting homomorphisms $\de^k$
in (\ref{eqn:3411}) are all zero. Hence, when we take $k = 0$, and we 
replace the two occurrences of  
``$\cdots$'' in (\ref{eqn:3411}) by ``$0$'', we get a short exact sequence.  
Finally, a calculation like in the proof of Lemma \ref{lem:3610} shows that 
\[ \opn{Coker}(\phi^0) \cong
\lim_{i \to} \, \opn{Hom}_{\dcat{D}(A)}(L, M_i) . \qedhere \]
\end{proof}

\begin{lem} \label{lem:3406} 
Let $L$ be a compact object of $\dcat{D}(A)$, let $M$ be a DG $A$-module, 
and let $\{ F_j(M) \}_{j \geq -1}$ be a filtration of $M$ in
$\dcat{C}_{\mrm{str}}(A)$ such that 
$\bigcup_j F_j(M) = M$. Given a morphism 
$\phi : L \to M$ in $\dcat{D}(A)$, there is an index $j_1 \in \N$, and a 
morphism $\phi_{j_1} : L \to F_{j_1}(M)$ in $\dcat{D}(A)$, 
such that 
$\phi = \opn{Q}(\ga_{j_1}) \circ \phi_{j_1}$,
where $\ga_{j_1} : F_{j_1}(M) \to M$ is the inclusion. 
\end{lem}

\begin{proof}
Here $M \cong \lim_{j \to} \, F_j(M)$ in $\dcat{C}_{\mrm{str}}(A)$. By Lemma 
\ref{lem:3610} we know that $M$ is a homotopy colimit of the direct system 
$\{ F_j(M) \}_{j \geq -1}$. According to Lemma \ref{lem:3410} we 
get a bijection 
\begin{equation} \label{eqn:3420}
\lim_{j \to} \, \opn{Hom}_{\dcat{D}(A)}(L, F_j(M)) \cong  
\opn{Hom}_{\dcat{D}(A)}(L, M) . 
\end{equation}
So there is some $j_1 \in \N$ and some 
$\phi_{j_1} \in \opn{Hom}_{\dcat{D}(A)}(L, F_{j_1}(M))$
that goes to $\phi$ under the bijection (\ref{eqn:3420}).
\end{proof}

The next lemma is adapted from \cite[Lemma tag=09R2]{SP}. 

\begin{lem} \label{lem:3412}
Let $L$ be a compact object of $\dcat{D}(A)$, let $P$ be a semi-free
DG $A$-module of finite extension length, and let $\phi : L \to P$ be a 
morphism 
in $\dcat{D}(A)$.
Then there is a finite semi-free DG $A$-module $P'$, with a morphism 
$\phi' : L \to P'$ in $\dcat{D}(A)$ and a monomorphism 
$\ga' : P' \to P$ in $\dcat{C}_{\mrm{str}}(A)$, such that 
$\phi = \opn{Q}(\ga') \circ \phi'$. 
\end{lem}

\begin{proof}
Let $\{ F_j(P) \}_{j \geq -1}$ be a semi-free filtration of $P$ of finite 
extension length; say $F_{j_1}(P) = P$. We will construct DG submodules
\[ P_0 \sub P_1 \sub \cdots \sub P_{j_1} \sub P_{j_1 + 1} \]
of $P$, such that properties (a)-(c) below hold for every 
$k \in [0, j_1 + 1]$. For each such $k$ and each $j \geq -1$ we define  
$F_j(P_k) := F_j(P) \cap P_k$.
\begin{itemize}
\item[(a)] The filtration $\{ F_j(P_k) \}_{j \geq -1}$ is a semi-free 
filtration on $P_k$.

\item[(b)] The free DG $A$-module $\opn{Gr}^F_j(P_k)$ is finite free for all 
$j \geq k$. 

\item[(c)] The morphism $\phi$ factors through $P_k$; namely there is a 
morphism $\phi_k : L \to P_k$ in $\dcat{D}(A)$ s.t.\ 
$\phi = \opn{Q}(\ga_k) \circ \phi_k$, where $\ga_k : P_k \to P$ is the 
inclusion. 
\end{itemize}

For $k = 0$, the DG module $P' := P_0$ is finite semi-free. The morphism 
$\phi' := \phi_0 : L \to P'$
in $\dcat{D}(A)$ and the inclusion 
$\ga' := \ga_0 : P' \to P$
in $\dcat{C}_{\mrm{str}}(A)$ 
will satisfy $\phi = \opn{Q}(\ga') \circ \phi'$. 

The construction of the DG modules $P_k$ and the morphisms $\phi_k$ is by 
descending induction on $k$, starting from $k = j + 1$. 
We start with $P_{j_1 +1 } := P$ and $\phi_{j_1 + 1} := \phi$. 
Properties (a)-(c) hold trivially here. 

Now to the inductive step. Suppose that $k \in [1, \ldots, j + 1]$, 
and we already found a DG submodule 
$P_k \sub P$ and a morphism $\phi_k : L \to P_k$ as required. 
We are going to produce a DG submodule 
$P_{k - 1} \sub P_k$ such that 
$\opn{Gr}^F_{j}(P_{k - 1}) = \opn{Gr}^F_{j}(P_k)$
for all $j \neq k - 1$, and 
$\opn{Gr}^F_{k - 1}(P_{k - 1})$ is a finite free DG submodule of 
$\opn{Gr}^F_{k - 1}(P_k)$. 
This means that we have to shrink $F_{k - 1}(P_k)$ suitably.

For every $j \geq 0$ the DG $A$-module $\opn{Gr}^F_{j}(P_k)$ is free, and 
we choose a basis for it, namely a collection 
$\{ \bar{p}_z \}_{z \in Z_j}$
of elements such that 
$\opn{Gr}^F_{j}(P_k) = \bigoplus_{z \in Z_j} A \cd \bar{p}_z$
in $\dcat{C}_{\mrm{str}}(A)$. 
Moreover, we choose these bases such that the indexing set 
$Z_j$ is finite for every $j \geq k$ (this is possible by property (b)), 
and $Z_j$ is empty for every $j \geq j_1 + 1$ (due to the fact that 
$F_{j_1}(P) = P$). 

By definition 
$\opn{Gr}^F_{j}(P_k) = F_{j}(P_k) / F_{j - 1}(P_k)$,
so we can lift the basis 
$\{ \bar{p}_z \}_{z \in Z_j}$ of $\opn{Gr}^F_{j}(P_k)$ to 
a collection $\{ p_z \}_{z \in Z_j}$ of homogeneous elements 
$p_z \in F_{j}(P_k)$. 
Let's define the indexing set $Z := \coprod_{j \geq k} Z_j$.
Then there are decompositions
\begin{equation} \label{eqn:5085}
P^{\natural}_k = F_{k - 2}(P_k)^{\natural} \oplus
\Bigl( \bigoplus\nolimits_{z \in Z_{k - 1}} \, A^{\natural} \cd p_z \Bigr) 
\oplus 
\Bigl( \bigoplus\nolimits_{z \in Z} \, A^{\natural} \cd p_z \Bigr) 
\end{equation}
and 
\[ F_{k - 1}(P_k)^{\natural} = F_{k - 2}(P_k)^{\natural} \oplus 
\Bigl( \bigoplus\nolimits_{z \in Z_{k - 1}} \, A^{\natural} \cd p_z \Bigr)   \]
in $\dcat{G}_{\mrm{str}}(A^{\natural})$, and each $A^{\natural} \cd p_z$ is a 
free graded $A^{\natural}$-module of rank $1$. 

Because $Z$ is a finite set, there is a finite subset $Y \sub Z_{k - 1}$ such 
that 
\[ \d(p_z) \in F_{k - 2}(P_k) \oplus
\Bigl( \bigoplus\nolimits_{y \in Y} \, A \cd p_y \Bigr) \oplus 
\Bigl( \bigoplus\nolimits_{z' \in Z} \, A \cd p_{z'} \Bigr) \]
for every $z \in Z$. 
For every $z \in Z_{k - 1}$ the element 
$\bar{p}_z \in \opn{Gr}^F_{k - 1}(P_k)$ satisfies 
$\d(\bar{p}_z) = 0$, and therefore the element 
$p_z \in F_{k - 1}(P_k)$ satisfies $\d(p_z) \in F_{k - 2}(P_k)$.
In particular this is true for $p_y$ with $y \in Y$. 
And $F_{k - 2}(P_k)$ is a DG submodule of $P_k$.
Combining all these observations we see that the graded submodule
\begin{equation} \label{eqn:5086}
Q^{\natural} := F_{k - 2}(P_k)^{\natural} \oplus
\Bigl( \bigoplus\nolimits_{y \in Y} \, A^{\natural} \cd p_y \Bigr) \oplus 
\Bigl( \bigoplus\nolimits_{z \in Z} \, A^{\natural} \cd p_z \Bigr) 
\end{equation}
of $P_k^{\natural}$ actually underlies a DG submodule $Q \sub P_k$. 
(Compare this to the decomposition (\ref{eqn:5085}) of $P_k^{\natural}$.)

Consider the quotient DG $A$-module $N := P_k / Q$,
with the canonical projection $\al: P_k \to N$, that's an epimorphism in 
$\dcat{C}_{\mrm{str}}(A)$. 
Define the set $X := Z_{k - 1} - Y$
and the elements 
$n_x := \al(p_x) \in N$. 
Then $N^{\natural}$ is a free graded $A^{\natural}$-module, with basis 
indexed by $X$, i.e.\ 
$N^{\natural} = \bigoplus\nolimits_{x \in X} A^{\natural} \cd n_x$.
Furthermore, because 
$\d(p_x) \in F_{k - 2}(P_k) \sub Q$
for all $x \in X$, we have $\d(n_x) = 0$;
so actually $N$ is a free DG $A$-module with basis $\{ n_x \}_{x \in 
X}$. 

Next consider the morphism
$\th : L \to N$ in $\dcat{D}(A)$ with formula 
$\th := \opn{Q}(\al) \circ \phi_k$.
By Lemma \ref{lem:3405} there is a finite subset $X_0 
\sub X$ such that $\th$ factors through the DG submodule
$N_0 := \bigoplus\nolimits_{x \in X_0} A \cd n_x$. 
I.e.\ there's a morphism $\th_0 : L \to N_0$ in $\dcat{D}(A)$
s.t.\ $\th = \opn{Q}(\be_0) \circ \th_0$, where 
$\be_0 : N_0 \to N$ is the inclusion (a monomorphism in 
$\dcat{C}_{\mrm{str}}(A)$). Define 
$\bar{N} := N / N_0 = \opn{Coker}(\be_0)$,
and let $\de : N \to \bar{N}$ be the projection
(an epimorphism in $\dcat{C}_{\mrm{str}}(A)$). 
Since $\de \circ \be_0 = 0$ in $\dcat{C}_{\mrm{str}}(A)$, 
it follows that 
\[ \opn{Q}(\de) \circ \th = \opn{Q}(\de) \circ \opn{Q}(\be_0) \circ \th_0 = 
\opn{Q}(\de \circ \be_0) \circ \th_0 = 0 \]
in $\dcat{D}(A)$. 

The subset $Y \cup X_0$ of $Z_{k - 1}$ is finite. 
We define 
\begin{equation} \label{eqn:3415}
P_{k - 1}^{\natural} := F_{k - 2}(P_k)^{\natural} \oplus
\Bigl( \bigoplus\nolimits_{z \in Y \cup X_0} \, A^{\natural}\cd p_z \Bigr) 
\oplus \Bigl( \bigoplus\nolimits_{z \in Z} \, A^{\natural} \cd p_z \Bigr) . 
\end{equation}
Then $P_{k - 1}^{\natural}$ is the underlying graded module of a DG submodule 
$P_{k - 1} \sub P_k$, slightly bigger than $Q$.
(Compare this formula to (\ref{eqn:5085}) and (\ref{eqn:5086}).)
By construction the DG module $P_{k - 1}$ has properties (a)-(b) 
with index $k - 1$. Indeed, 
\[ \opn{Gr}^F_j(P_{k - 1}) \cong 
\begin{cases}
\opn{Gr}^F_j(P_{k}) & \tup{if} \ j \geq k \ \tup{or} \ j \leq k - 2 
\\[0.4em]
\bigoplus\nolimits_{z \in Y \cup X_0} \, A \cd \bar{p}_z 
& \tup{if} \ j = k - 1 
\end{cases} \]

Let us denote by $\ep : P_{k - 1} \to P_{k}$ the inclusion. 
There is a short exact sequence
$0 \to P_{k - 1} \xar{\ep} P_k \xar{\de \circ \al} \bar{N} \to  0$
in $\dcat{C}_{\mrm{str}}(A)$, and this becomes a distinguished triangle 
$P_{k - 1} \xar{\ep} P_k \xar{\de \circ \al} \bar{N} \xar{\triangle}$
in $\dcat{D}(A)$. Applying the cohomological functor 
$\opn{Hom}_{\dcat{D}(A)}(L, -)$ to this distinguished triangle gives rise to a 
long exact sequence 
\[ \cdots \to \opn{Hom}_{\dcat{D}(A)}(L, P_{k - 1}) \to 
\opn{Hom}_{\dcat{D}(A)}(L, P_{k}) \to 
\opn{Hom}_{\dcat{D}(A)}(L, \bar{N}) \to \cdots . \]
The morphism $\phi_k$ belongs to the middle term here. 
We know that 
\[ \opn{Q}(\de \circ \al) \circ \phi_k = 
\opn{Q}(\de) \circ \opn{Q}(\al) \circ \phi_k = 
\opn{Q}(\de) \circ \th = 0 . \]
Thus there is a morphism 
$\phi_{k - 1} : L \to P_{k - 1}$
in $\dcat{D}(A)$ s.t.\ 
$\phi_k = \opn{Q}(\ep) \circ \phi_{k - 1}$. The morphism $\phi_{k - 1}$
has property (c) with index $k - 1$. 
\end{proof}

The derived tensor-evaluation morphism $\opn{ev}^{\mrm{R, L}}_{L, M, N}$ was 
introduced in Theorem \ref{thm:4320}. 
Finite semi-free, algebraically perfect and compact DG modules were defined 
in Definitions \ref{dfn:3402}(3), \ref{dfn:3401} and 
\ref{dfn:3405}, respectively. 

\begin{thm} \label{thm:3400}
Let $A$ be a DG ring and let $L$ be a DG $A$-module. The following four 
conditions are equivalent.
\begin{itemize}
\rmitem{i} $L$ is algebraically perfect. 

\rmitem{ii} $L$ is a direct summand in $\dcat{D}(A)$ of a finite 
semi-free DG 
$A$-module. 

\rmitem{iii} For every DG ring $B$, every DG module
$M \in \dcat{D}(A \ot_{} B^{\mrm{op}})$ and every DG module
$N \in \dcat{D}(B)$, the tensor-evaluation morphism 
\[ \opn{ev}^{\mrm{R, L}}_{L, M, N} : 
\opn{RHom}_{A}(L, M) \ot^{\mrm{L}}_{B} N \to 
\opn{RHom}_{A}(L, M \ot^{\mrm{L}}_{B} N) \]
in $\dcat{D}(\K)$ is an isomorphism. 

\rmitem{iv} $L$ is a compact object of $\dcat{D}(A)$.
\end{itemize}
\end{thm}

\begin{proof} \mbox{} 

\smallskip \noindent
(ii) $\Rightarrow$ (i): We use Proposition \ref{prop:4470}. A finite free DG 
$A$-module is algebraically perfect. A finite semi-free DG module is 
obtained by finitely many cones from finite free DG modules, so it is also 
algebraically perfect. Finally, our DG module $L$ is a direct summand of a 
finite semi-free DG module, so it is algebraically perfect.

\medskip \noindent 
(i) $\Rightarrow$ (iii): Fixing $M$ and $N$, we have two triangulated functors 
$F, G : \dcat{D}(A) \to \dcat{D}(\K)$,
namely 
$F := \opn{RHom}_{A}(-, M) \ot^{\mrm{L}}_{B} N$
and
$G := \opn{RHom}_{A}(-, M \ot^{\mrm{L}}_{B} N)$.
There is a morphism of triangulated functors 
$\ep_{} := \opn{ev}^{\mrm{R, L}}_{(-), M, N} : F \to  G$,
and we want to prove that $\ep_L : F(L) \to  G(L)$ is an isomorphism.
Trivially for $L := A$ the morphism $\ep_{L}$ is an isomorphism. 
According to Proposition \ref{prop:4603} the morphism $\ep_L$ is an isomorphism 
for every algebraically perfect DG module $L$. 

\medskip \noindent 
(iii) $\Rightarrow$ (iv): Let $\{ N_z \}_{z \in Z}$ be a collection of objects 
of $\dcat{D}(A)$. Define $N := \bigoplus_{z \in Z} N_z$.
We have to prove that the canonical homomorphism of $\K$-modules 
\[ \Phi_L : \bigoplus\nolimits_{z \in Z} \opn{Hom}_{\dcat{D}(A)}(L, N_z) \to 
\opn{Hom}_{\dcat{D}(A)} (L, N) \]
is bijective. Take $B := A$ and 
$M := A \in \dcat{D}(A \ot_{} A^{\mrm{op}})$. 
Consider the commutative diagram of $\K$-modules (\ref{eqn:4800}).
\begin{figure}
\begin{equation} \label{eqn:4800} 
\UseTips  \xymatrix @C=6ex @R=6ex {
\bigoplus_{z \in Z} \opn{H}^0 \bigl( \opn{RHom}_A(L, A) \ot^{\mrm{L}}_{A} N_z
\bigr) 
\ar[r]^(0.56){\th}_(0.56){\cong}
\ar[d]_{\bigoplus \opn{H}^0(\opn{ev}^{\mrm{R, L}}_{L, M, N_z})}^{\cong}
&
\opn{H}^0 \bigl( \opn{RHom}_A(L, A) \ot^{\mrm{L}}_{A} N \bigr) 
\ar[d]^{\opn{H}^0(\opn{ev}^{\mrm{R, L}}_{L, M, N})}_{\cong}
\\
\bigoplus_{z \in Z}  \opn{H}^0 \bigl( \opn{RHom}_A(L, A \ot^{\mrm{L}}_{A} N_z)
\bigr) 
\ar[r]
\ar[d]^{\cong}_{\bigoplus \ze_{N_z}}
&
\opn{H}^0 \bigl( \opn{RHom}_A(L, A \ot^{\mrm{L}}_{A} N) \bigr) 
\ar[d]_{\cong}^{\ze_{N}}
\\
\bigoplus_{z \in Z} \opn{Hom}_{\dcat{D}(A)}(L, N_z)
\ar[r]^(0.55){\Phi_L}
&
\opn{Hom}_{\dcat{D}(A)}(L, N)
}
\end{equation}
\end{figure}
The homomorphism $\th$ is bijective because both $(- \ot^{\mrm{L}}_{A} -)$ 
and $\opn{H}^0$ respect all direct sums 
(see Theorem \ref{thm:3140} and Propositions \ref{prop:1585} and 
\ref{prop:4300}).
The homomorphisms $\opn{H}^0(\opn{ev}^{\mrm{R, L}}_{L, M, N_z})$ and 
$\opn{H}^0(\opn{ev}^{\mrm{R, L}}_{L, M, N})$ are bijective by condition (iii). 
The homomorphisms $\ze_{N_z}$ and $\ze_{N}$ are bijective because 
$A \ot^{\mrm{L}}_{A} (-) \cong \opn{Id}$ and by Corollary \ref{cor:2120}. 
We conclude that $\Phi_L$ is bijective. 

\medskip \noindent 
(iv) $\Rightarrow$ (ii): 
Choose a semi-free resolution $\rho : P \to L$ in $\dcat{C}_{\mrm{str}}(A)$.
We have an isomorphism 
$\phi := \opn{Q}(\rho)^{-1} : L \to P$
in $\dcat{D}(A)$. Let $\{ F_j(P) \}_{j \geq -1}$ be a semi-free filtration on 
$P$. According to Lemma \ref{lem:3406} there is an index $j_1$ such that, 
letting $P' := F_{j_1}(P)$, and letting $\ga' : P' \to P$ be the inclusion, 
there exists a morphism $\phi' : L \to P'$ in $\dcat{D}(A)$ satisfying 
$\phi = \opn{Q}(\ga') \circ \phi'$. 

Now the DG module $P'$ is semi-free of finite extension length. Lemma 
\ref{lem:3412} says that there is a  finite semi-free DG $A$-module $P''$, with 
a morphism $\phi'' : L \to P''$ in $\dcat{D}(A)$, and a monomorphism 
$\ga'' : P'' \to P'$ in $\dcat{C}_{\mrm{str}}(A)$, such that 
$\phi' = \opn{Q}(\ga'') \circ \phi''$. 

Consider the morphism 
$\psi'' : P'' \to L$
in $\dcat{D}(A)$ defined by 
$\psi'' := \opn{Q}(\rho) \circ \opn{Q}(\ga') \circ \opn{Q}(\ga'')$.
We have equality 
\[ \psi'' \circ \phi'' = 
\opn{Q}(\rho) \circ \opn{Q}(\ga') \circ \opn{Q}(\ga'') \circ \phi'' = 
\opn{Q}(\rho) \circ \phi =  \opn{id}_L . \]
Thus $L$ is a retract in $\dcat{D}(A)$ of the finite semi-free DG module $P''$. 
But by Theorem \ref{thm:4160}, $L$ is then a direct summand of $P''$ in
$\dcat{D}(A)$.
\end{proof}

\begin{cor} \label{cor:3480}
Let $A$ and $B$ be DG rings, let 
$F : \dcat{D}(A) \to \dcat{D}(B)$ 
be an equivalence of triangulated categories, and let 
$L \in \dcat{D}(A)$. Then $L$ is an algebraically perfect DG $A$-module if and 
only if $F(L)$ is an algebraically perfect DG $B$-module.
\end{cor}

\begin{proof}
Combine Theorem \ref{thm:3400} and Proposition \ref{prop:3545}. 
\end{proof}

\begin{rem} \label{rem:3425}
Assume $A$ is a {\em commutative DG ring} (i.e.\ nonpositive and strongly 
commutative, see Definitions \ref{dfn:3091} and \ref{dfn:3090})  .
Then $\bar{A} := \opn{H}^0(A)$ is a commutative ring, and there is a DG ring 
homomorphism $A \to \bar{A}$. 
As explained in \cite{Ye10}, 
given a multiplicatively closed set $\bar{S} \sub \bar{A}$, letting 
$S \sub A^0$ be the preimage of $\bar{S}$, the localized commutative DG ring 
$A_S := A \ot_{A^0} A^0_S$ satisfies 
$\opn{H}^0(A_S) = \bar{A}_{\bar{S}}$. 
We are going to use the convenient notation $A_{\bar{S}} := A_S$.
In case $\bar{S} = \{ \bar{s}^k \}_{k \in \N}$
for  some element $\bar{s} \in \bar{A}$, we shall write 
$A_{\bar{s}} := A_{\bar{S}}$.

A sequence of elements $(\bar{s}_1, \ldots, \bar{s}_n)$ in $\bar{A}$ is called 
a {\em covering sequence} if 
$\opn{Spec}(\bar{A}) = \bigcup_{i} \, \opn{Spec}(\bar{A}_{\bar{s}_i})$.
This just means that 
$\bar{A} = \sum_{i} \bar{A} \cd \bar{s}_i$. 
For each $i$ there is a localized DG ring $A_{\bar{s}_i}$. 

A DG module $L \in \dcat{D}(A)$ is called {\em geometrically perfect} 
\index{Differential graded module! geometrically perfect}
if for some covering sequence $(\bar{s}_1, \ldots, \bar{s}_n)$ of $\bar{A}$
there are finite semi-free DG $A_{\bar{s}_i}$-modules $P_i$, and isomorphisms 
$P_i \cong A_{\bar{s}_i} \ot_A L$ in $\dcat{D}(A_{\bar{s}_i})$. 

In case $A$ is a commutative ring, for which $A = A^0 = \bar{A}$, the 
definition becomes simpler: a complex $L \in \dcat{D}(A)$ is geometrically 
perfect if there is a covering sequence $(s_1, \ldots, s_n)$ of 
$A$, bounded complexes of finite free $A_{s_i}$-modules $P_i$,
and isomorphisms $P_i \cong A_{s_i} \ot_A L$ in $\dcat{D}(A_{s_i})$. 
This is the classical definition of a perfect complex of $A$-modules, see 
\cite{SGA6}. 

According to \cite[Corollary 5.21]{Ye10}, a DG $A$-module $L$ is geometrically 
perfect iff it is algebraically perfect. And according to 
\cite[Theorem 5.11]{Ye10}, $L$ is geometrically perfect iff 
$L \in \dcat{D}^-(A)$, and its derived reduction 
$\bar{A} \ot^{\mrm{L}}_{A} L \in \dcat{D}(\bar{A})$ is geometrically perfect.
\end{rem}

In the special case of a ring we can say more. 

\begin{thm} \label{thm:3415}
Let $A$ be a ring and let $L$ be a DG $A$-module. The following two conditions 
are equivalent.
\begin{itemize}
\rmitem{i}  $L$ is an algebraically perfect DG $A$-module. 

\rmitem{ii} $L$ is isomorphic in $\dcat{D}(A)$ to a bounded complex of 
finitely generated projective $A$-modules.
\end{itemize}
\end{thm}

The proof is after this lemma. The concentration $\opn{con}(M)$ of a graded 
module was defined in Definition \ref{dfn:2125}. 

\begin{lem} \label{lem:3415}
In the situation of the theorem, if $L$ is algebraically perfect
and $\opn{H}(L) \neq 0$, then 
$\opn{con}(\opn{H}(L)) = [i_0, i_1]$ for some integers $i_0 \leq i_1$, and 
$\opn{H}^{i_1}(L)$ is a finitely presented $A$-module. 
\end{lem}

\begin{proof}
According to Theorem \ref{thm:3400}, $L$ is a direct summand in $\dcat{D}(A)$ 
of a finite semi-free DG $A$-module $P$. But $A$ is a ring, so $P$ is a bounded 
complex of (finitely generated free) $A$-modules. This implies that 
$\opn{H}(P)$ is bounded. But $\opn{H}(L)$ is a direct summand, in 
$\dcat{G}_{\mrm{str}}(A)$, of $\opn{H}(P)$; so it is also bounded. 
Because $\opn{H}(L)$ is nonzero, its concentration is a bounded nonempty integer
interval $[i_0, i_1]$. Note that 
$i_1 = \opn{sup}(\opn{H}(L))$. 

Corollary \ref{cor:3340}
says that there is a semi-free resolution $Q \to L$ with 
$\opn{sup}(Q) = i_1$. As in the proof of the implication 
(iv) $\Rightarrow$ (ii) in the proof of Theorem \ref{thm:3400}, 
there is a finite semi-free DG $A$-module $Q''$, with a monomorphism 
$Q'' \to Q$ in $\dcat{C}_{\mrm{str}}(A)$, 
such that $L$ is a direct summand of $Q''$ in $\dcat{D}(A)$. 
Now $Q''$ is a bounded complex of finitely generated free $A$-modules, with 
$\opn{sup}(Q'') = i_1$, and hence $\opn{H}^{i_1}(Q'')$ is a finitely 
presented $A$-module. But $\opn{H}^{i_1}(L)$ is a direct summand of 
$\opn{H}^{i_1}(Q'')$, so it too is a finitely presented $A$-module.
\end{proof}

\begin{proof}[Proof of the Theorem]
The implication (ii) $\Rightarrow$ (i) is easy; compare to the proof of the 
implication (ii) $\Rightarrow$ (i) of Theorem \ref{thm:3400}.

The implication (i) $\Rightarrow$ (ii) will be proved in two steps. 

\smallskip \noindent
Step 1. Here $L$ is an $A$-module, which is algebraically perfect as a DG 
$A$-module. According to Proposition \ref{prop:3415}, the DG module $L$ has 
finite projective dimension, say $d \in \N$. (The case $L \cong 0$ can be 
excluded.) This is also the projective dimension of the module $L$ in the 
classical sense (see Proposition \ref{prop:4290}). 
By Lemma \ref{lem:3415} the module $L$ is finitely generated. 
We will proceed by induction on $d$. 

If $d = 0$ then $P := L$ is already a finitely generated projective $A$-module. 
If $d \geq 1$, then we can find a finite free $A$-module $P^0$, with a 
surjection $\eta : P^0 \to L$ in $\dcat{M}(A)$. Let 
$L' := \opn{Ker}(\eta)$, so we have a short exact sequence 
$0 \to L' \to P^0 \xar{\eta} L \to 0$
in $\dcat{M}(A)$. 
According to Proposition \ref{prop:2165} this becomes a distinguished triangle 
\begin{equation} \label{eqn:3428}
L' \to P^0 \xar{\opn{Q}(\eta)} L \xar{\triangle}
\end{equation}
in $\dcat{D}(A)$.
Because both $L$ and $P^0$ are perfect, we see that the DG $A$-module $L'$ is 
also perfect. Thus $L'$ is a finitely generated $A$-module of finite 
projective dimension. However, by the usual syzygy argument 
(cf.\ \cite[Section 5.1.1]{Row}) -- or by examination of the long exact 
sequence gotten by applying 
$\opn{Hom}_{\dcat{D}(A)}(-, N)$, for an $A$-module $N$, 
to the distinguished triangle (\ref{eqn:3428}) -- we see that 
the projective dimension of $L'$ is $\leq d - 1$. Induction says 
that there is an isomorphism $L' \cong P'$ in $\dcat{D}(A)$ for some bounded 
complex of finitely generated projective $A$-modules $P'$.
Plugging this isomorphism into the distinguished triangle (\ref{eqn:3428})
we obtain a distinguished triangle 
\begin{equation} \label{eqn:3421}
P' \xar{\psi} P^0 \xar{\opn{Q}(\eta)} L \xar{\triangle} 
\end{equation}
in $\dcat{D}(A)$. Because $P'$ is K-projective, there is a homomorphism 
$\til{\psi} : P' \to P^0$ in $\dcat{C}_{\mrm{str}}(A)$ such that 
$\psi = \opn{Q}(\til{\psi})$. Let 
$P := \opn{Cone}(\til{\psi})$, the standard cone on $\til{\psi}$. 
Then $P$ is a bounded complex of finitely generated projective $A$-modules, and 
the distinguished triangle (\ref{eqn:3421}) gives rise to an isomorphism 
$P \cong L$ in $\dcat{D}(A)$.

\medskip \noindent
Step 2. Here $L$ is an arbitrary algebraically perfect DG $A$-module, nonzero 
in $\dcat{D}(A)$. By Lemma \ref{lem:3415} we have 
$\opn{con}(\opn{H}(L)) = [i_0, i_1]$ for integers $i_0 \leq i_1$, and 
$\opn{H}^{i_1}(L)$ is a finitely generated $A$-module. 
Let $n := \opn{amp}(\opn{H}(L)) = i_1 - i_0 \in \N$.
We proceed by induction on $n$. 

If $n = 0$ then $L$ is isomorphic to the translation of an $A$-module, and we 
are done by step 1. On the other hand, if $n \geq 1$, then choose a finite free 
$A$-module $P^{i_1}$, with a surjection 
$\bar{\eta} : P^{i_1} \to \opn{H}^{i_1}(L)$ in $\dcat{M}(A)$. 
This can be lifted to a homomorphism 
$\eta : P^{i_1}[-i_1] \to L$ in $\dcat{C}_{\mrm{str}}(A)$. 
Consider the DG $A$-module 
$L' := \opn{Cone}(\eta)$.
There is a distinguished triangle 
\begin{equation} \label{eqn:3423}
P^{i_1}[-i_1] \xar{\opn{Q}(\eta)} L \xar{} L' \xar{\triangle} 
\end{equation}
in $\dcat{D}(A)$. 
The DG module $L'$ is also algebraically perfect. There is a long exact 
cohomology sequence 
\[ \begin{aligned}
& \cdots \to 0 \to \opn{H}^{i_0}(L) \to \opn{H}^{i_0}(L') \to \cdots 
\to 0 \to \opn{H}^{i_1 - 2}(L) \to \opn{H}^{i_1 - 2}(L')
\\
& \quad 
\to 0 \to \opn{H}^{i_1 - 1}(L) \to \opn{H}^{i_1 - 1}(L') \to
P^{i_1} \xar{\bar{\eta}} \opn{H}^{i_1}(L) \to \opn{H}^{i_1}(L') \to 0 \to \cdots
\end{aligned} \]
Because $\bar{\eta}$ is surjective we know that $\opn{H}^{i_1}(L') = 0$.
Hence $\opn{con}(\opn{H}(L')) \sub [i_0, i_1 - 1]$
and $\opn{amp}(\opn{H}(L')) \leq n - 1$.

The induction assumption says that there is an isomorphism $L' \cong P'$
in $\dcat{D}(A)$ for some bounded complex of 
finitely generated projective $A$-modules $P'$.
Plugging this isomorphism into the triangle (\ref{eqn:3423}), and then turning 
it, we obtain a distinguished triangle 
\begin{equation} \label{eqn:3424} 
P'[-1] \xar{\psi} P^{i_1}[-i_1] \xar{\opn{Q}(\eta)} L \xar{\triangle} 
\end{equation}
in $\dcat{D}(A)$. Because $P'[-1]$ is K-projective, there is a homomorphism 
$\til{\psi} : P'[-1] \to P^{i_1}[-i_1]$ in $\dcat{C}_{\mrm{str}}(A)$ such that 
$\psi = \opn{Q}(\til{\psi})$. Let 
$P := \opn{Cone}(\til{\psi})$, the standard cone on $\til{\psi}$. 
Then $P$ is a bounded complex of finitely generated projective $A$-modules, and 
the distinguished triangle (\ref{eqn:3424}) gives rise to an isomorphism 
$P \cong L$ in $\dcat{D}(A)$.
\end{proof}

\begin{cor} \label{cor:4170}
Let $A$ be a ring and let $L$ be a complex of $A$-modules. The following two 
conditions are equivalent.
\begin{itemize}
\rmitem{i}  $L$ is a compact object of $\dcat{D}(A)$.

\rmitem{ii} $L$ is isomorphic in $\dcat{D}(A)$ to a bounded complex of 
finitely generated projective $A$-modules.
\end{itemize} 
\end{cor}

\begin{proof}
Combine Theorems \ref{thm:3415} and \ref{thm:3400}.
\end{proof}

Another special case is when the DG ring $A$ is nonpositive and cohomologically 
left pseudo-noetherian; see Definition \ref{dfn:3180}.

\begin{thm} \label{thm:3416}
Let $A$ be a cohomologically left pseudo-noetherian nonpositive DG ring,
and let $L$ be a DG $A$-module. The following two conditions are 
equivalent.
\begin{itemize}
\rmitem{i}  $L$ is an algebraically perfect DG $A$-module. 

\rmitem{ii} $L$ has finite projective dimension and belongs to 
$\dcat{D}^{-}_{\mrm{f}}(A)$.
\end{itemize}
\end{thm}

% \newpage
\begin{proof} \mbox{}

\nopagebreak
\smallskip \noindent
(i) $\twoto$ (ii): Since $A \in \dcat{D}^{-}_{\mrm{f}}(A)$, and 
$\dcat{D}^{-}_{\mrm{f}}(A)$ is a saturated full triangulated subcategory of 
$\dcat{D}(A)$, we see that $L \in \dcat{D}^{-}_{\mrm{f}}(A)$. By Proposition 
\ref{prop:3415}, $L$ has finite projective dimension. 

\medskip \noindent
(ii) $\twoto$ (i): We may assume that $\opn{H}(L) \neq 0$. 
Let $i_1 := \opn{sup}(\opn{H}(L)) \in \Z$, and let $n \in \N$ be the 
projective dimension of $L$.  By Theorem \ref{thm:3340}
there is a quasi-isomorphism $\rho : P \to L$ in $\dcat{C}_{\mrm{str}}(A)$ from 
a pseudo-finite semi-free DG $A$-module $P$ such that 
$\opn{sup}(P) = i_1$. Let $\{ F_j(P) \}_{j \geq -1}$ be a pseudo-finite 
semi-free filtration of $P$ as in Definition \ref{dfn:1925}(1),
with the same upper bound $i_1$.
Define $P' := F_{i_1 + n}(P)$ and $P'' := P / P'$. 
The DG $A$-module $P''$ is concentrated in the degree interval 
$[-\infty, i_1 -n - 1]$, and therefore 
\[ \opn{con} \bigl( \opn{H} ( \opn{RHom}_A(L, P'')) \bigr) \sub
[-\infty, - 1] . \]
This implies that 
\begin{equation} \label{eqn:3425}
 \opn{Hom}_{\dcat{D}(A)}(L, P'') \cong 
\opn{H}^0 \bigl( \opn{RHom}_A(L, P'') \bigr) = 0 .
\end{equation}

The DG $A$-module $P'$ is finite semi-free. 
Let $\ga' : P' \to P$ be the inclusion. 
From the distinguished triangle 
$P' \xar{\opn{Q}(\ga')} P \to P'' \xar{\triangle}$
in $\dcat{D}(A)$ we get an exact sequence of $\K$-modules
\[ \opn{Hom}_{\dcat{D}(A)}(L, P') \to 
\opn{Hom}_{\dcat{D}(A)}(L, P) \to 
\opn{Hom}_{\dcat{D}(A)}(L, P'') . \]
The isomorphism 
$\phi := \opn{Q}(\rho)^{-1} : L \to P$
belongs to the middle term above. But by (\ref{eqn:3425}) we see that $\phi$ 
comes from some morphism 
$\phi' : L \to P'$ in $\dcat{D}(A)$. 
Because $\opn{Q}(\rho \circ \ga') \circ \phi' = \opn{id}_L$, and using 
Theorem \ref{thm:4160},
we see that $L$ is a direct summand of $P'$ in $\dcat{D}(A)$. 
Theorem \ref{thm:3400} says that $L$ is algebraically perfect. 
\end{proof}

\begin{cor} \label{cor:3425}
Let $A$ be a left noetherian ring, and let $L$ be a DG $A$-module. The 
following 
two conditions are equivalent.
\begin{itemize}
\rmitem{i}  $L$ is an algebraically perfect DG $A$-module. 

\rmitem{ii} $L$ has finite projective dimension and belongs to 
$\dcat{D}^{\mrm{b}}_{\mrm{f}}(A)$.
\end{itemize}
\end{cor}

\begin{proof}
Combine Theorems \ref{thm:3415} and \ref{thm:3416}.
\end{proof}

\begin{rem} \label{rem:3426}
Here are some historical notes on perfect and compact objects.

Perfect complexes in algebraic geometry were introduced by A. Grothen\-dieck 
et.\ al.\ in \cite{SGA6}. For a scheme $X$, let $\dcat{D}_{\mrm{qc}}(X)$
be the derived category  of complexes of $\OO_X$-modules with 
quasi-coherent cohomology. The definition of a perfect complex given in 
\cite{SGA6} was this: $\LL$ is perfect if locally it is isomorphic to a bounded 
complex of finite rank free $\OO_X$-modules. This coincides with our notion of 
geometrically perfect complex, as in Remark \ref{rem:3425}, when $A$ is a 
commutative ring and $X = \opn{Spec}(A)$.

The idea of defining finiteness properties of an object $L$ in a category 
$\cat{D}$ via the commutation of the functor $\opn{Hom}_{\cat{D}}(L, -)$ with 
suitable coproducts goes back to the early days of category theory. Objects $L$ 
such that $\opn{Hom}_{\cat{D}}(L, -)$ commutes with coproducts were called 
{\em small}. Grothendieck \cite{Gro} used small objects in his construction of 
injective resolutions in what was eventually called a Grothendieck abelian 
category. P. Freyd \cite{Fre} used small objects in the proof of the 
Freyd-Mitchell Theorem. In algebraic topology, small objects were used by E. 
Brown \cite{Brown} in his celebrated Representability Theorem.

J. Rickard \cite{Ric1} proved a slightly weaker version of our Corollary 
\ref{cor:4170}. Shortly afterwards B. Keller \cite{Kel1} deduced from A. 
Neeman's work \cite{Ne0} a result that 
is essentially our Theorem \ref{thm:3400} -- he did not have condition (iii), 
but on the other hand he considered the more general derived category 
$\dcat{D}(\cat{A})$ of DG modules over a DG category $\cat{A}$, whereas we only 
consider $\dcat{D}(A)$ for a DG ring $A$, which is a single-object DG category. 

R.W. Thomason \cite{Thom} discovered that the perfect complexes in 
$\dcat{D}_{\mrm{qc}}(X)$, when $X$ is a quasi-projective scheme over a 
ring $\K$, are precisely the compact objects in $\dcat{D}_{\mrm{qc}}(X)$. 
Neeman \cite{Ne0} realized the connection between the work of Thomason and 
that of the topologists  A.K. Bousfield and D. Ravenel.
The terminology switch from ``small object'' to ``compact object'' seems to 
have taken place in \cite{Ne0}. 
Neeman's book \cite{Ne1} is  essentially devoted 
to the study of $\al$-compactly generated triangulated categories, for a 
regular cardinal number $\al$, and to generalizations of the Brown 
Representability Theorem. Note that what we called ``compact object'' in 
Definition \ref{dfn:3405} is, in the framework of \cite{Ne0}, an 
``$\aleph_0$-compact object''. 

The paper \cite{BonVdB} by A. Bondal and M. Van den Bergh was very influential 
in promoting the role of compact objects in algebraic geometry. 
In the last 20 years there has been a proliferation in the presence of compact 
(or perfect) objects in research in the areas of derived algebraic geometry,  
noncommutative algebraic geometry and mathematical physics -- much of this due 
to the influence of M. Kontsevich. For more details we recommend looking in the 
online reference \cite{nLab}. 
\end{rem}



\mysubsection{Derived Morita Theory} \label{subsec:der-morita}

Recall that Convention \ref{conv:3485} is in effect.
In particular all DG rings are $\K$-central, and 
$\ot = \ot_{\K}$.

\begin{dfn} \label{dfn:3525}
Let $A$ and $B$ be DG rings. The objects of the category 
$\dcat{C}(A \ot B^{\mrm{op}})$
are called 
\index{Differential graded bimodule}
{\em DG $A$-$B$-bimodules}. 
\end{dfn}

There is a commutative diagram (\ref{eqn:3535}) 
in $\catt{DGRng} \centover \K$.

\begin{equation} \label{eqn:3535}
\UseTips  \xymatrix @C=3ex @R=3ex {
&
\ A \ot B^{\mrm{op}}
\\
A
\ar[ur]
&
&
B^{\mrm{op}}
\ar[ul]
\\
&
\K 
\ar[ul]
\ar[ur]
\ar[uu]
}
\end{equation}

By restriction there is a commutative diagram of DG 
functors, and a commutative diagram of triangulated functors, see 
(\ref{eqn:3536}).

\begin{equation} \label{eqn:3536}
\UseTips  \xymatrix @C=3ex @R=3ex {
&
\ \dcat{C}(A \ot B^{\mrm{op}})
\ar[dl]_{\opn{Rest}_{A}}
\ar[dr]^{\opn{Rest}_{B^{\mrm{op}}}}
\ar[dd]^{\opn{Rest}_{\K}}
\\
\dcat{C}(A)
\ar[dr]_{\opn{Rest}_{\K}}
&
&
\dcat{C}(B^{\mrm{op}})
\ar[dl]^{\opn{Rest}_{\K}}
\\
&
\dcat{C}(\K) 
}
\quad 
\UseTips  \xymatrix @C=3ex @R=3ex {
&
\ \dcat{D}(A \ot B^{\mrm{op}})
\ar[dl]_{\opn{Rest}_{A}}
\ar[dr]^{\opn{Rest}_{B^{\mrm{op}}}}
\ar[dd]^{\opn{Rest}_{\K}}
\\
\dcat{D}(A)
\ar[dr]_{\opn{Rest}_{\K}}
&
&
\dcat{D}(B^{\mrm{op}})
\ar[dl]^{\opn{Rest}_{\K}}
\\
&
\dcat{D}(\K) 
}
\end{equation}

There are several manipulations of DG bimodules that we are going to use. 
First we note that $(A^{\mrm{op}})^{\mrm{op}} = A$.
Next, given DG $\K$-modules $M$ and $N$, let 
\begin{equation} \label{eqn:3550} 
\opn{br}_{M, N} : M \ot N \iso N \ot M
\end{equation}
be the isomorphism in $\dcat{C}(\K)$ defined by 
\begin{equation} \label{eqn:3553} 
\opn{br}_{M, N}(m \ot n) := (-1)^{i \cd j} \cd n \ot m
\end{equation}
for homogeneous elements $m \in M^i$ and $n \in N^j$. 
In Subsection \ref{subsec:gr-alg} the isomorphism $\opn{br}_{M, N}$ was called 
the braiding of the symmetric monoidal category $\dcat{C}(\K)$.
If $M \in \dcat{C}(A^{\mrm{op}})$, then we can view $M$ either as a 
left DG $A^{\mrm{op}}$-module, or as a right DG $A$-module.
The opposite holds for $N \in \dcat{C}(A)$. It is not hard to see that there 
is an isomorphism 
\begin{equation} \label{eqn:3551} 
M \ot_A N \iso N \ot_{A^{\mrm{op}}} M
\end{equation}
in $\dcat{C}(\K)$ with the same formula (\ref{eqn:3553}).
This works also for DG rings: there are isomorphisms 
\begin{equation} \label{eqn:3552} 
A \ot B^{\mrm{op}} \iso B^{\mrm{op}} \ot A
\end{equation}
and
\begin{equation} \label{eqn:3554} 
A \ot B^{\mrm{op}} \iso (B \ot A^{\mrm{op}})^{\mrm{op}}
\end{equation}
in $\catt{DGRng} \centover \K$, with a formulas like (\ref{eqn:3553}). 

There is a DG bifunctor
\begin{equation} \label{eqn:3537}
(- \ot^{}_{B} -) :  \dcat{C}(A \ot B^{\mrm{op}}) \times \dcat{C}(B)
\to \dcat{C}(A)  . 
\end{equation}
Here is a variation of Theorem \ref{thm:2107}.

\begin{prop} \label{prop:3535}
The bifunctor \tup{(\ref{eqn:3537})} has a triangulated left derived bifunctor
\[ (- \ot^{\mrm{L}}_{B} -) :
\dcat{D}(A \ot B^{\mrm{op}}) \times \dcat{D}(B)
\to \dcat{D}(A) . \]
If $P \in \dcat{D}(B)$ is a K-flat DG module, then for every 
$M \in \dcat{D}(A \ot B^{\mrm{op}})$
the morphism 
\[ \eta^{\mrm{L}}_{M, P} : M \ot^{\mrm{L}}_{B} P \to M \ot^{}_{B} P \]
in $\dcat{D}(A)$ is an isomorphism. 
\end{prop}

\begin{proof}
This is because $\dcat{C}(B)$ has enough K-flat objects. See
Theorem \ref{thm:2106} and Lemma \ref{lem:2109}.
\end{proof}

\begin{rem} \label{rem:3431}
There is a delicate issue here. Even though the category 
$\dcat{C}(A \ot B^{\mrm{op}})$ has enough K-flat objects, 
they can not be used to calculate $(- \ot^{\mrm{L}}_{B} -)$. 
The reason is this: given $N \in \dcat{C}(B)$,
the DG $(A \ot B^{\mrm{op}})$-modules that are acyclic 
for the DG functor $(-) \ot_B N$ are those that are 
{\em K-flat over $B^{\mrm{op}}$}. In general, a K-flat DG 
$(A \ot B^{\mrm{op}})$-module $M$ is not K-flat as a $B^{\mrm{op}}$-module. 

A sufficient condition for a K-flat DG $(A \ot B^{\mrm{op}})$-module 
$M$ to be K-flat over $B^{\mrm{op}}$ is that $A$ is K-flat as a DG 
$\K$-module. In Subsection \ref{subsec:flat-dg-rng} we will make heavy use of 
this fact. 
\end{rem}

There is another DG bifunctor we shall want to use:
\begin{equation} \label{eqn:3538}
\opn{Hom}_A(-, -) :  \dcat{C}(A \ot B^{\mrm{op}})^{\mrm{op}} \times 
\dcat{C}(A) \to \dcat{C}(B)  . 
\end{equation}
Here is a variation of Theorem \ref{thm:2005}. 

\begin{prop} \label{prop:3536}
The bifunctor \tup{(\ref{eqn:3538})} has a triangulated right derived 
bifunctor
\[ \opn{RHom}_A(-, -) :
\dcat{D}(A \ot B^{\mrm{op}})^{\mrm{op}} \times \dcat{D}(A)
\to \dcat{D}(B) . \]
If $I \in \dcat{D}(A)$ is a K-injective DG module, then for every 
$M \in \dcat{D}(A \ot B^{\mrm{op}})$
the morphism 
\[ \eta^{\mrm{R}}_{M, I} : \opn{Hom}_A(M, I) \to \opn{RHom}_A(M, I) \]
in $\dcat{D}(B)$ is an isomorphism. 
\end{prop}

\begin{proof}
This is because $\dcat{C}(A)$ has enough K-injective objects. See
Theorem \ref{thm:4225}, Theorem \ref{thm:2000} and 
Lemma \ref{lem:2000}.
\end{proof}

\begin{rem} \label{rem:3432}
Like in Remark \ref{rem:3431}, even though the category 
$\dcat{C}(A \ot  B^{\mrm{op}})$ has enough K-projective 
DG modules, they can not be used to calculate \lb $\opn{RHom}_A(-, N)$. 
This is because a K-projective DG $(A \ot  B^{\mrm{op}})$-module $P$ is 
not, in general, K-projective over $A$. 

A sufficient condition for a K-projective DG $(A \ot B^{\mrm{op}})$-module 
$P$ to be K-projective over $A$ is that $B$ is K-projective over $\K$. 
This is a bit stronger than the condition that $B$ is K-flat over $\K$.
\end{rem}

\begin{prop} \label{prop:3537}
Let $L \in \dcat{D}(A \ot B^{\mrm{op}})$, and consider the $\K$-linear 
triangulated functors
\[ F := \opn{RHom}_A(L, -) : \dcat{D}(A) \to \dcat{D}(B) \]
and 
\[ G := L \ot^{\mrm{L}}_{B} (-) : \dcat{D}(B) \to \dcat{D}(A) . \]
Then\tup{:}
\begin{enumerate}
\item The functor $F$ is a right adjoint of $G$.

\item The functor $F$ is an equivalence if and only if the functor $G$ is an 
equivalence. 
\end{enumerate}
\end{prop}

\begin{proof} \mbox{}

\smallskip \noindent
(1) Given $M \in \dcat{D}(B)$ and $N \in \dcat{D}(A)$ we choose 
a K-projective resolution $P \to M$ in $\dcat{C}_{\mrm{str}}(B)$ and a 
K-injective resolution $N \to I$ in $\dcat{C}_{\mrm{str}}(A)$.
We get isomorphisms 
$F(N) \cong F(I) \cong \opn{Hom}_A(L, I)$ 
and
$G(M) \cong G(P) \cong L \ot_B P$ 
in $\dcat{D}(B)$ and $\dcat{D}(A)$ respectively.
These give rise to isomorphisms of $\K$-modules
\[ \begin{aligned}
& \opn{Hom}_{\dcat{D}(B)} \bigl( M, F(N) \bigr) \cong^{\diamond}
\opn{H}^0 \bigl( \opn{RHom}_B \bigl( M, F(N) \bigr) \bigr)
\\
& \qquad \cong
\opn{H}^0 \bigl( \opn{Hom}_B \bigl( P, \opn{Hom}_A(L, I) \bigr) \bigr)
\cong^{\dag}
\opn{H}^0 \bigl( \opn{Hom}_A(L \ot_B P, I) \bigr)
\\
& \qquad \cong
\opn{H}^0 \bigl( \opn{RHom}_A \bigl( G(M), N \bigr) \bigr)
\cong^{\diamond} \opn{Hom}_{\dcat{D}(A)} \bigl( G(M), N \bigr) 
\end{aligned} . \]
The isomorphism $\cong^{\dag}$ is due to the noncommutative Hom-tensor 
adjunction (see \cite[Theorem 2.11]{Rot}), and the 
isomorphisms $\cong^{\diamond}$ are by Corollary \ref{cor:2120}.
The composed isomorphism 
\[ \opn{Hom}_{\dcat{D}(B)} \bigl( M, F(N) \bigr) \cong
\opn{Hom}_{\dcat{D}(A)} \bigl( G(M), N \bigr) \]
is functorial in $M$ and $N$. 

\medskip \noindent 
(2) Clear from (1). 
\end{proof}

\begin{dfn} \label{dfn:3430}
Let $A$ and $B$ be DG rings. A DG module
$L \in \dcat{D}(A \ot B^{\mrm{op}})$
is called a {\em pretilting DG $A$-$B$-bimodule}% 
\index{Tilting! pre{\indash} DG bimodule}
if it satisfies the equivalent 
conditions in Proposition \ref{prop:3537}(2). 
\end{dfn}

\begin{rem} \label{rem:3490}
The literature has several different meanings for the word ``tilting''. 
See Remark \ref{rem:3436} for a historical survey. 
The commonly perceived meaning of a
tilting object, say as in the book \cite{AnHaKr}, is very close to what we call 
a pretilting DG bimodule (when $A$ and $B$ are rings). 

Notice that there is a lack of symmetry in the 
definition of a pretilting DG $A$-$B$-bimodule $L$.
This lack of symmetry will disappear when we talk about {\em tilting DG 
bimodules} in Subsection \ref{subsec:flat-dg-rng}. 
\end{rem}

\begin{dfn} \label{dfn:3446}
Let $A$ and $B$ be DG rings, and let 
$M \in \dcat{C}(A \ot B^{\mrm{op}})$. We say that $M$ is {\em K-flat} 
(resp.\ {\em K-injective}, resp.\ {\em K-projective}, resp.\ {\em semi-free},
resp.\ {\em algebraically perfect}) {\em over $A$}, or {\em on the $A$ side},
if $\opn{Rest}_A(M) \in \dcat{C}(A)$ is K-flat (resp.\  K-injective, resp.\ 
K-projective, resp.\ semi-free, resp.\ algebraically perfect).
Likewise we define the properties {\em on the $B^{\mrm{op}}$ side}. 
\end{dfn}

\begin{prop} \label{prop:3431}
Suppose $L \in \dcat{D}(A \ot B^{\mrm{op}})$ is a pretilting DG 
$A$-$B$-bimodule. Then $L$ is an algebraically perfect DG $A$-module.
\end{prop}

\begin{proof}
Under the equivalence $G : \dcat{D}(B) \to \dcat{D}(A)$ from Proposition 
\ref{prop:3537} we have $G(B) \cong L$. Of course $B \in \dcat{D}(B)$ is an 
algebraically perfect DG $B$-module. Now use Corollary \ref{cor:3480}. 
\end{proof}

Observe that an object $N \in \dcat{D}(A)$ is {\em nonzero} if it is 
not a zero object of the category $\dcat{D}(A)$, i.e.\ if $N \not\cong 0$ in
$\dcat{D}(A)$. This is equivalent to the condition that 
$\opn{H}(N) \neq 0$. The clarification is important for the next definition. 

\begin{dfn} \label{dfn:3431}
Let $\cat{E} \sub \dcat{D}(A)$ be a full triangulated subcategory that's closed 
under infinite direct sums. 
\begin{enumerate}
\item An object $L \in \dcat{D}(A)$ is called a 
{\em compact object relative to $\cat{E}$}%
\index{Differential graded module! compact}
if the functor 
\[ \opn{Hom}_{\dcat{D}(A)}(L, -) : \cat{E} \to \dcat{M}(\K) \]
commutes with infinite direct sums. 

\item An object $L \in \dcat{D}(A)$ is called a 
{\em generator relative to $\cat{E}$}%
\index{Differential graded module! generator}
if for every nonzero object $N \in \cat{E}$ there exists some 
$i \in \Z$ such that 
$\opn{Hom}_{\dcat{D}(A)}(L, N[i]) \neq 0$. 

\item An object $L \in \cat{E}$ is called a 
{\em compact generator of $\cat{E}$} 
if it is both a compact object relative to 
$\cat{E}$ and a generator 
relative to $\cat{E}$.
\end{enumerate}
\end{dfn}

\begin{exa} \label{exa:3430}
For every positive integer $r$ the free DG module 
$P := A^r$ is a compact generator of $\cat{E} := \dcat{D}(A)$.
\end{exa}

\begin{exa} \label{exa:3431}
Assume $A$ is a commutative noetherian ring and 
$\bsym{a} = \lb (a_1, \ldots, a_n)$ is a finite sequence of elements in it. 
(Or, more generally, $A$ is commutative and $\bsym{a}$ is a {\em weakly 
proregular sequence}, as in \cite{PSY}.)
Let $\a \sub A$ be the ideal generated by $\bsym{a}$. 
A complex $M \in \dcat{D}(A)$ is called {\em cohomologically $\a$-torsion}
if all its cohomology modules $\opn{H}^i(M)$ are $\a$-torsion. 
Consider the full subcategory 
$\cat{E} := \dcat{D}_{\a\tup{-tor}}(A)$ of $\dcat{D}(A)$ on the 
cohomologically $\a$-torsion complexes. 
This is full triangulated subcategory that's closed 
under infinite direct sums. 
The Koszul complex $\opn{K}(A; \bsym{a})$ is a compact generator of 
$\dcat{D}_{\a\tup{-tor}}(A)$. 
See \cite[Proposition 5.1]{PSY2}.
\end{exa}

\begin{prop} \label{prop:3547}
Let $A$ and $B$ be DG rings, $F : \dcat{D}(A) \to \dcat{D}(B)$ an 
equivalence of triangulated categories, 
$\cat{E} \sub \dcat{D}(A)$ a full triangulated subcategory that's closed 
under infinite direct sums, and $L \in \dcat{D}(A)$. 
\begin{enumerate}
\item $L$ is a compact object relative to $\cat{E}$ if and only if 
$F(L)$ is a compact object relative to $F(\cat{E})$. 

\item $L$ is a generator relative to $\cat{E}$ if and only if 
$F(L)$ is a generator relative to $F(\cat{E})$. 
\end{enumerate}
\end{prop}

\begin{proof}
For item (1) use the proof of Proposition \ref{prop:3545}. Item (2) is trivial. 
\end{proof}

\begin{dfn} \label{dfn:3465}
Let $A$ and $B$ be DG rings, let $\cat{E} \sub \dcat{D}(A)$ be a 
full triangulated subcategory that is closed under infinite direct sums, and 
let $L \in \dcat{D}(A \ot B^{\mrm{op}})$. 
\begin{enumerate}
\item We say that $L$ is a {\em compact object relative to $\cat{E}$ on the 
$A$ side} if $\opn{Rest}_A(L) \in \dcat{D}(A)$ is a compact object relative to 
$\cat{E}$, in the sense of Definition \lb \ref{dfn:3431}(1).

\item We say that $L$ is a {\em generator relative to $\cat{E}$ on the 
$A$ side} if $\opn{Rest}_A(L) \in \dcat{D}(A)$ is a generator relative to 
$\cat{E}$, in the sense of Definition \ref{dfn:3431}(2).

\item We say that $L$ is a {\em compact generator of $\cat{E}$ on the $A$ side} 
if $\opn{Rest}_A(L) \in \dcat{D}(A)$ is 
a compact generator of $\cat{E}$, in the sense of Definition \ref{dfn:3431}(3).
\end{enumerate}
\end{dfn}

\begin{lem} \label{lem:3436}
Let $\cat{E} \sub \dcat{D}(A)$ be a full triangulated subcategory that is 
closed under infinite direct sums, and let 
$L \in \dcat{D}(A \ot B^{\mrm{op}})$. The following 
two conditions are equivalent.
\begin{itemize}
\rmitem{i} The functor
$\opn{RHom}_A(L, -)|_{\cat{E}} : \cat{E} \to \dcat{D}(B)$
commutes with infinite direct sums.

\rmitem{ii} $L$ is a compact object relative to $\cat{E}$ on the $A$ side. 
\end{itemize}
\end{lem}

\begin{proof}
Let's write 
$F := \opn{RHom}_A(L, -)$. We know that for every $M \in \dcat{D}(A)$ and 
$j \in \Z$ there are isomorphisms
\[ \opn{Hom}_{\dcat{D}(A)}(L, M[j]) \cong 
\mrm{H}^0(F(M[j])) \cong \mrm{H}^j (F(M))  \]
in $\dcat{M}(\K)$, and they are functorial in $M$.
We see that $L$ is a compact object relative to $\cat{E}$ on the $A$ side iff
the functor
$\mrm{H} \circ F|_{\cat{E}} : \cat{E} \to \dcat{G}_{\mrm{str}}(\K)$
commutes with infinite direct sums. But the functor 
$ \mrm{H} : \dcat{D}(B) \to \dcat{G}_{\mrm{str}}(\K)$ 
commutes with infinite direct sums and is conservative (see Corollary 
\ref{cor:2145}). 
So $\mrm{H} \circ F|_{\cat{E}}$ commutes with infinite direct 
sums iff $F|_{\cat{E}}$ does. 
\end{proof}

\begin{lem} \label{lem:3435}
Let $A$ and $B$ be DG rings, let 
$F, G : \dcat{D}(A) \to \dcat{D}(B)$ be
triangulated functors that commute with infinite direct sums,
and let $\eta : F \to G$ be a morphism of triangulated functors. Assume that
$\eta_A: F(A) \to G(A)$ is an isomorphism. Then $\eta$ is an isomorphism. 
\end{lem}

\begin{proof}
Let us denote by $\cat{E}$ the full subcategory of $\dcat{D}(A)$ on the objects 
$M$ such that $\eta_M: F(M) \to G(M)$ is an isomorphism. We have to prove that 
$\cat{E} = \dcat{D}(A)$. 

By Proposition \ref{prop:4603}, 
the category $\cat{E}$ is a full triangulated subcategory of $\dcat{D}(A)$.
Since both functors $F, G$ commute with infinite direct sums,
we know that $\cat{E}$ is closed under infinite direct sums.

We are given that the free DG module $A$ belongs to $\cat{E}$.
Hence every free DG $A$-module $P$ is in $\cat{E}$.
Because $\cat{E}$ is triangulated, it follows that every semi-free DG 
$A$-module $P$ of finite extension length (Definition \ref{dfn:3411}) belongs 
to 
it. 

Next consider an arbitrary semi-free DG $A$-module $P$.
Let $\{ F_j(P) \}_{j \geq -1}$ be a semi-free filtration of $P$
(see Definition \ref{dfn:1575}). 
Then $P$ is a homotopy colimit of the direct system 
$\{ F_j(P) \}_{j \geq -1}$, and we have a distinguished triangle 
\begin{equation} \label{eqn:3440}
\bigoplus\nolimits_{j \in \N} F_j(P) \xar{\ \phi \ } 
\bigoplus\nolimits_{j \in \N} F_j(P) 
\xar{ \ \ } P \xar{\triangle}
\end{equation}
in $\dcat{D}(A)$; see Definition \ref{dfn:3410}. 
Each $F_j(P)$ is a semi-free DG $A$-module of finite extension length. It 
follows that $P \in \cat{E}$. 

Finally, every DG $A$-module $M$ admits a quasi-isomorphism 
$P \to M$ with $P$ semi-free. Therefore $M \in \cat{E}$.
\end{proof}

\begin{lem} \label{lem:3470}
Let $A$ and $B$ be DG rings, let $\cat{E}$ be a full triangulated 
subcategory of $\dcat{D}(A)$ which is closed under infinite direct sums
and isomorphisms, and let $G : \dcat{D}(B) \to \dcat{D}(A)$ be a
triangulated functor that commutes with infinite direct sums.
Assume that $G(B) \in \cat{E}$. Then the essential image of $G$ is contained 
in $\cat{E}$. 
\end{lem}

\begin{proof}
This is like the proof of Lemma \ref{lem:3435}.
\end{proof}

Let $A$ be a DG ring. Given a DG $A$-module $M$, the DG ring 
$\opn{End}_A(M)$ acts on $M$ from the left; and thus 
$B := \opn{End}_A(M)^{\mrm{op}}$ 
acts on $M$ from the right. Because the right action of $B$ on $M$ commutes 
with the left action of $A$ on it, we see that 
$M \in \dcat{C}(A \ot B^{\mrm{op}})$.

Recall that the category $\dcat{C}(A)$ has enough K-projective and K-injective 
objects. 

\begin{thm}[Derived Morita] \label{thm:3430}
\index{Derived Morita theory}
Let $A$ be a DG ring, let 
$\cat{E} \sub \dcat{D}(A)$ be a full triangulated subcategory which is closed 
under infinite direct sums, and let $L$ be a compact 
generator of $\cat{E}$. Choose an isomorphism 
$L \cong P$ in $\dcat{D}(A)$, where the DG $A$-module $P$ is either 
K-projective or K-injective, and define the DG ring
$B := \opn{End}_A(P)^{\mrm{op}}$.
Consider the triangulated functors 
\[ F := \opn{RHom}_A(P, -) : \dcat{D}(A) \to \dcat{D}(B) \] 
and 
\[ G := P \ot_B^{\mrm{L}} (-) : \dcat{D}(B) \to \dcat{D}(A) . \]
Then the functor 
\[ F|_{\cat{E}} : \cat{E} \to \dcat{D}(B) \]
is an equivalence, with quasi-inverse $G$. 
\end{thm}

This theorem is a variant of \cite[Theorem 8.2]{Kel1} by B. Keller.
In fact, Keller worked in the more general setting of a DG category $\cat{A}$, 
whereas we only treat a DG ring $A$. For the history of these ideas see Remark 
\ref{rem:3436}.
A comparison to classical Morita theory can be found in Example \ref{exa:3435}.

Note that if we choose $P$ to be a K-projective DG $A$-module, then 
$F \cong \lb \opn{Hom}_A(P, -)$. 
However, even then, $P$ will rarely be K-flat as a DG $B^{\mrm{op}}$-module, so 
$G$ is only the left derived functor of $P \ot_B (-)$. 

\begin{proof} 
The proof is in four steps. 

\smallskip \noindent 
Step 1. We can assume that the category $\cat{E}$ is closed under isomorphisms 
in $\dcat{D}(A)$. The functor $G$ commutes with infinite direct sums (see 
Proposition \ref{prop:4300}), 
and $G(B) \cong L \in \cat{E}$. According to Lemma \ref{lem:3470} the image of 
$G$ is contained in $\cat{E}$. 

\medskip \noindent 
Step 2. We already know that the functor $F$ is right adjoint to $G$, by 
Proposition \ref{prop:3537}. 
The corresponding morphisms of triangulated functors are denoted by
\begin{equation} \label{eqn:3471}
\th : \opn{Id}_{\dcat{D}(B)} \to F \circ G \quad \tup{and} \quad
\ze : G \circ F \to \opn{Id}_{\dcat{D}(A)} . 
\end{equation}
By step 1 we see that  
$F|_{\cat{E}} : \cat{E} \to \dcat{D}(B)$ 
is the right adjoint of 
$G : \dcat{D}(B) \to \cat{E}$.
Hence (\ref{eqn:3471}) restricts to morphisms of triangulated functors
\begin{equation} \label{eqn:3472}
\th : \opn{Id}_{\dcat{D}(B)} \to F|_{\cat{E}} \circ G \quad \tup{and} \quad
\ze : G \circ F|_{\cat{E}} \to \opn{Id}_{\cat{E}}  . 
\end{equation}
We will prove that $\th$ and $\ze$ in (\ref{eqn:3472}) are isomorphisms. 
By Lemma \ref{lem:3436} the functor $F|_{\cat{E}}$ commutes with infinite 
direct sums. Therefore both $F|_{\cat{E}} \circ G$ and 
$G \circ F|_{\cat{E}}$ commute with infinite direct sums.

\medskip \noindent 
Step 3. Now we will prove that $\th$ is an isomorphism of functors; i.e.\
for every $N \in \dcat{D}(B)$ the morphism
$\th_N :  N \to (F|_{\cat{E}} \circ G) (N)$ 
is an isomorphism. The functors $\opn{Id}_{\dcat{D}(B)}$ and 
$F|_{\cat{E}} \circ G$ both commute with infinite direct sums.
Therefore, by Lemma \ref{lem:3435}, 
it suffices to check that $\th_B$ is an isomorphism in $\dcat{D}(B)$. 
But $\th_B$ is represented (both when $P$ is K-projective and when it is  
K-injective) by the canonical homomorphism 
$\til{\th}_B : B \to \opn{Hom}_{A}(P, P \otimes_B B)$,
in $\dcat{C}_{\mrm{str}}(B)$, 
which is clearly bijective. 

\medskip \noindent 
Step 4. Finally we will prove that $\ze$ is an isomorphism of functors. 
Take any $M \in \cat{E}$, and 
consider the distinguished triangle
\begin{equation} \label{eqn:3473}
(G \circ F|_{\cat{E}})(M) \xar{\zeta_{M}} M \to M' \xar{\triangle}
\end{equation}
in $\cat{E}$, in which $M' \in \cat{E}$ is the cone of $\zeta_{M}$
(see Definition \ref{dfn:3600}).

Applying $F$ and using the functorial isomorphism $\th$ we get a 
distinguished triangle
\begin{equation} \label{eqn:3474}
F|_{\cat{E}}(M) \xar{\opn{id}_{F|_{\cat{E}}(M)}} F|_{\cat{E}}(M)
\to F|_{\cat{E}}(M') \xar{\triangle}
\end{equation}
in $\dcat{D}(B)$. This implies that 
$F|_{\cat{E}}(M') \cong 0$ in $\dcat{D}(B)$, and thus, after applying the 
forgetful functor, we get $F|_{\cat{E}}(M') \cong 0$ in $\dcat{D}(\K)$. But 
$F|_{\cat{E}}(M') \cong \opn{RHom}_{A}(L, M')$,
so 
$\opn{Hom}_{\dcat{D}(A)}(L, M'[i]) = 0$ 
for all $i$. Because $L$ is generator of $\cat{E}$, it follows that 
$M' \cong 0$ in $\cat{E}$. Going back to the distinguished 
triangle (\ref{eqn:3473}) we conclude that $\ze_M$ is an isomorphism.
\end{proof}

\begin{rem} \label{rem:3435}
Suppose that in the theorem we were to choose some other isomorphism 
$L \cong P'$ in $\dcat{D}(A)$ to a K-projective or K-injective DG $A$-module 
$P$. Then the DG ring $B' := \opn{End}_A(P')^{\mrm{op}}$
would be related to $B$ as follows. Without loss of generality we can assume 
that either $P$ is K-projective or $P'$ is K-injective. Then there is a 
quasi-isomorphism $\phi : P \to P'$ in $\dcat{C}_{\mrm{str}}(A)$, 
unique up to homotopy, that respects the given isomorphisms 
$L \cong P$ and $L \cong P'$ in $\dcat{D}(A)$. 
Define the DG bimodule 
$N := \opn{Hom}_A(P, P') \in \dcat{C}(B \ot B'^{\, \mrm{op}})$, 
and the matrix DG ring 
$C := \sbmat{B & \, N[-1] \\[0.2em] 0 & B'}$,
which has a differential depending on $\phi$. 
Then the obvious DG ring homomorphisms $C \to B$ and $C \to B'$ are 
quasi-isomorphisms. See \cite[Proposition 3.3]{PSY2}. 
\end{rem}

\begin{exa}  \label{exa:3432}
In case $\cat{E} = \dcat{D}(A)$, the theorem shows that the DG bimodule 
$P \in \dcat{D}(A \ot B^{\mrm{op}})$ 
is a pretilting DG $A$-$B$-bimodule. 
\end{exa}

\begin{exa}  \label{exa:3435}
Assume $A$ is a ring, and $P \in \dcat{M}(A) = \cat{Mod} A$ is a 
{\em progenerator}, i.e.\ $P$ is a finitely generated projective $A$-module, 
such that every nonzero $A$-module $N$ admits a nonzero homomorphism 
$P \to N$. Let $B := \opn{End}_A(P)^{\mrm{op}}$. 
Because $P$ is a compact generator of $\dcat{D}(A)$, the theorem 
says that 
\[ F := \opn{Hom}_A(P, -) : \dcat{D}(A) \to \dcat{D}(B) \]
is an equivalence of triangulated categories. 

Classical Morita Theory 
(see e.g.\ \cite[Section 4.1]{Row}) says that $F$ restricts to an equivalence 
of abelian categories 
\[ F = \opn{Hom}_A(P, -) : \dcat{M}(A) \to \dcat{M}(B) . \]
Furthermore, classical Morita theory tells us that there is a bijection between 
the set of isomorphism classes of such $A$-$B$-bimodules $P$, and the set of 
isomorphism classes of linear equivalences 
$F : \dcat{M}(A) \to \dcat{M}(B)$. 
This last assertion is an open problem in derived Morita theory; see Remark 
\ref{rem:3565}. 
\end{exa}

\begin{rem} \label{rem:3470}
Here is a geometric variant of Theorem \ref{thm:3430}. Let $(X, \OO_X)$ be a 
scheme. We denote by 
$\dcat{D}_{\mrm{qc}}(X) = \dcat{D}_{\mrm{qc}}(\cat{Mod} \OO_X)$ 
the derived category of 
(unbounded) complexes of $\OO_X$-modules with quasi-coherent cohomology 
sheaves. This is a triangulated category with infinite direct sums, and with 
enough K-injective resolutions (by \cite[Theorem 4.5]{Spa}). 

Suppose $\LL \in \dcat{D}_{\mrm{qc}}(X)$ is a compact generator. Choose a 
K-injective resolution $\LL \to \II$ in 
$\dcat{C}_{\mrm{str}}(X)$, and define the DG ring 
$B := \opn{End}_X(\II)^{\mrm{op}}$.
Then the functor 
\begin{equation} \label{eqn:5088}
\opn{RHom}_X(\II, -) : \dcat{D}_{\mrm{qc}}(X) \to \dcat{D}(B)
\end{equation}
is an equivalence, with quasi-inverse 
\[ \II \ot^{\mrm{L}}_{B} (-) : \dcat{D}(B) \to \dcat{D}_{\mrm{qc}}(X) . \]
This result is stated as \cite[Corollary 3.1.8]{BonVdB}, without such a precise 
formulation, and the proof is attributed to B. Keller. Presumably the proof of 
Theorem \ref{thm:3430} above, with slight changes, would work also in this 
geometric context. 

Having a compact generator of $\dcat{D}_{\mrm{qc}}(X)$ is quite a general 
feature -- by \cite[Theorem 3.1.1]{BonVdB} this is true if $X$ is quasi-compact 
and quasi-separated. 

The first such result is perhaps due to A. Beilinson \cite{Bei}, who showed 
that for $X = \mbf{P}^n_{\K}$, the $n$-dimensional projective space over a 
field $\K$, the sheaf 
\[ \LL := \bigoplus\nolimits_{i = 0}^n \, \OO_X(i) \]
is a compact generator of $\dcat{D}_{\mrm{qc}}(X)$. The DG ring $B$ appearing 
in the equivalence (\ref{eqn:5088}) is actually a finite $\K$-ring (the path 
ring of a Kronecker quiver modulo commutation relations).  
\end{rem}

\begin{rem} \label{rem:3436}
Tilting theory and derived Morita theory have their origins in the 
{\em representation theory of finite dimensional algebras} (in our terminology 
these are ``finite central $\K$-rings'', where $\K$ is a base field). Among the 
first examples of tilting are the {\em reflection functors} in the paper
\cite{BeGePo} by I.N. Bernstein, I.M. Gelfand and V.A. Ponomarev. Later these 
functors were understood to be of the form $\opn{Hom}_{A}(T, -)$ for a {\em 
tilting $A$-module} $T$. 

These ideas were generalized by M. Auslander, M. Platzeck and I. Reiten 
\cite{AuPlRe}, and later by S. Brenner and M.C.R. Butler \cite{BrnBu}, who 
coined the term {\em tilting functor}. 
These concepts were further clarified by  D. Happel, C.M. Ringel and 
K. Bongartz (see \cite{Hap0}, \cite{Hap} and \cite{HapRi}). In \cite{Hap0} 
Happel showed that for a tilting $A$-module $T$, 
with opposite endomorphism ring $B := \opn{End}_{A}(T)^{\mrm{op}}$, 
the functor $\opn{RHom}_{A}(T,-)$ is an equivalence of triangulated categories 
$\dcat{D}_{\mrm{f}}^{\mrm{b}}(A) \to \dcat{D}_{\mrm{f}}^{\mrm{b}}(B)$.  
This result was slightly generalized by E. Cline, B. Parshall and L. Scott
\cite{ClPaSc}. 

J. Rickard \cite{Ric1}, \cite{Ric2} was the first to talk about {\em tilting 
complexes} (as opposed to tilting modules). He introduced two-sided tilting 
complexes (that we will discuss in Subsection \ref{subsec:tilting}), and proved 
the celebrated Theorem \ref{thm:3565} (under some boundedness conditions). The 
generalization of derived Morita theory to unbounded derived categories, 
and from rings to DG categories, was done by B. Keller \cite{Kel1}.
In this seminal paper Keller also gave new construction of two-sided 
tilting complexes,  and characterized algebraic triangulated categories with a 
compact generator as those that are equivalent to derived categories of DG 
rings. 

For a thorough survey of tilting theory see the book \cite{AnHaKr}.
\end{rem}





\mysubsection{DG Bimodules over K-Flat DG Rings} \label{subsec:flat-dg-rng}

Recall that Convention \ref{conv:3485} is in effect.
In particular all DG rings are $\K$-central, and 
$\ot = \ot_{\K}$.

\begin{dfn} \label{dfn:3560}
Let $A$ be a DG $\K$-ring. We call $A$ a {\em K-flat DG $\K$-ring} 
\index{Differential graded ring! K-flat}
if $A$ is K-flat as a DG $\K$-module (see Definition \ref{dfn:1525}). 
The category of K-flat central DG $\K$-rings is denoted by 
$\catt{DGRng} \fcentover \K$. 
\end{dfn}

\begin{exa} \label{exa:3445}
If $\K$ is a {\em field}, then every DG  $\K$-ring $A$ is K-flat over 
$\K$.
\end{exa}

\begin{exa} \label{exa:3501}
If $A$ is a {\em semi-free} DG $\K$-ring (either commutative or 
noncommutative), then $A$ is a K-flat DG $\K$-ring. 
See Definition \ref{dfn:4309}, Proposition \ref{prop:4305} and Remark 
\ref{rem:4945}. 
\end{exa}

\begin{exa} \label{exa:3500}
If $A$ is a {\em nonpositive} DG $\K$-ring, and each 
$A^i$ is a {\em flat} $\K$-module, then $A$ is a K-flat DG $\K$-ring.
This includes the case of a flat $\K$-ring $A$. 
See Proposition \ref{prop:4535}. 
\end{exa}

\begin{exa} \label{exa:3525}
A very special case of Examples \ref{exa:3445} and \ref{exa:3500} is when 
$\K$ is a field and $A$ is a $\K$-ring (better known as a unital 
associative $\K$-algebra). 
This setting is the one commonly used in ring theory, see 
\cite{ATV}, \cite{ArZh}, \cite{Ye1}, \cite{VdB}, \cite{StaVdB}, \cite{YeZh1}. 
\end{exa}

\begin{dfn} \label{dfn:3595}
Let $A$ and $B$ be K-flat DG $\K$-rings. The 
{\em derived category of DG $B$-$A$-bimodules} 
\index{Derived category! of DG bimodules}
is the $\K$-linear triangulated category 
$\dcat{D}(B \ot A^{\mrm{op}})$. 
\end{dfn}

See Remark \ref{rem:1420} regarding the derived category of DG 
$B$-$A$-bimodules in the nonflat case.

\begin{dfn} \label{dfn:3450}
For a K-flat DG $\K$-ring $A$ we write 
$A^{\mrm{en}} := A \ot_{} A^{\mrm{op}}$, 
and call it the {\em enveloping DG ring} 
\index{Differential graded ring! enveloping}
of $A$ (relative to $\K$).
\end{dfn}

A explained in formula (\ref{eqn:3554}), the enveloping DG ring has a canonical 
isomorphism 
$A^{\mrm{en}} \iso (A^{\mrm{en}})^{\mrm{op}}$.

From here to the end of this section we assume the next convention, that 
strengthens Convention \ref{conv:3485}.  

\begin{conv} \label{conv:3445} 
In addition to the stipulations of Convention \ref{conv:3485}, we also 
assume by default that all DG rings are {\em K-flat central DG 
$\K$-rings}.
\end{conv}

Recall that a homomorphism of DG rings $f : A \to B$ induces a forgetful 
functor
$\opn{Rest}_f : \dcat{C}(B) \to \dcat{C}(A)$ 
that we call {\em restriction}. This is an exact DG functor, and thus 
it induces a triangulated functor  
\begin{equation} \label{eqn:3446}
\opn{Rest}_f : \dcat{D}(B) \to \dcat{D}(A) .
\end{equation}

\begin{prop} \label{prop:3445}
Let $f : A \to B$ be a DG ring homomorphism.
The triangulated functor $\opn{Rest}_f$ is conservative. Namely a morphism 
$\phi : M \to N$ in $\dcat{D}(B)$ is an isomorphism if and only if the morphism 
\[ \opn{Rest}_f(\phi) : \opn{Rest}_f(M) \to \opn{Rest}_f(N) \]
in  $\dcat{D}(A)$ is an isomorphism.
\end{prop}

\begin{proof}
We know by Corollary \ref{cor:2145} that the functor 
$\opn{H} : \dcat{D}(B) \to \dcat{G}(\K)$
is conservative. But 
$\opn{H} = \opn{H} \circ \opn{Rest}_f$.
\end{proof}

\begin{lem} \label{lem:3445}
Let $A$ and $B$ be DG rings.
\begin{enumerate}
\item If $P \in \dcat{C}(A \ot B^{\mrm{op}})$ is K-flat, then $P$ is K-flat 
over 
$A$.

\item  If $I \in \dcat{C}(A \ot B^{\mrm{op}})$ is K-injective, then $I$ is 
K-injective over $A$.

\item If $B$ is K-projective as a DG $\K$-module, and if 
$P \in \dcat{C}(A \ot_{\K} B^{\mrm{op}})$ is K-projective, 
then $P$ is K-projective over $A$.

\item If $B$ is a semi-free DG $\K$-ring, and if 
$P \in \dcat{C}(A \ot_{\K} B^{\mrm{op}})$ is semi-free, 
then $P$ is semi-free over $A$.
\end{enumerate}
\end{lem}

\begin{proof}
(1) and (2) are direct consequences of the canonical isomorphisms  
\[ M \ot_{A} P \cong (M \ot_{} B) \ot_{A \ot_{} B} P \]
and
\[ \opn{Hom}_A(N, I) \cong  \opn{Hom}_{A \ot_{} B}(N \ot_{} B, I)  \]
in $\dcat{C}_{\mrm{str}}(\K)$, 
for $M \in \dcat{C}(A^{\mrm{op}})$ and $N \in \dcat{C}(A)$,
together with the fact that $B$ is K-flat over $\K$. 
Items (3) and (4) are left as an exercise. 
\end{proof}

\begin{exer} \label{exer:3445}
Prove items (3) and (4) in this lemma. 
\end{exer}

\begin{prop} \label{prop:3450}
Let $A$, $B$ and $C$ be DG rings.
\begin{enumerate}
\item The DG bifunctor 
\[ (- \ot_{B} -) :  \dcat{C}(A \ot B^{\mrm{op}}) \times 
\dcat{C}(B \ot C^{\mrm{op}}) \to \dcat{C}(A \ot C^{\mrm{op}})  \]
has a triangulated left derived bifunctor 
\[ \bigl(  (- \ot^{\mrm{L}}_{B} -), \eta^{\mrm{L}} \bigr) :  
\dcat{D}(A \ot_{} B^{\mrm{op}}) \times 
\dcat{D}(B \ot_{} C^{\mrm{op}}) \to \dcat{D}(A \ot_{} C^{\mrm{op}}) . \]

\item Given $M \in \dcat{C}(A \ot_{} B^{\mrm{op}})$
and 
$N \in \dcat{C}(B \ot_{} C^{\mrm{op}})$,
such that $M$ is K-flat over $B^{\mrm{op}}$ or $N$ is K-flat over $B$, the 
morphism 
\[ \eta^{\mrm{L}}_{M, N} : M \ot_B^{\mrm{L}} N \to M \ot_{B} N \]
in $\dcat{D}(A \ot_{} C^{\mrm{op}})$ is an isomorphism. 

\item Suppose we are given DG ring homomorphisms 
$A' \to A$, $f : B' \to B$ and $C' \to C$, such that $f$ is a 
quasi-isomorphism. 
Then the diagram 
\[ \UseTips \xymatrix @C=10ex @R=6ex {
\dcat{D}(A \ot_{} B^{\mrm{op}}) \times \dcat{D}(B \ot_{} C^{\mrm{op}})
\ar[r]^(0.6){ (- \ot^{\mrm{L}}_{B} -) }
\ar[d]_{\opn{Rest} \times \opn{Rest}}
&
\dcat{D}(A \ot_{} C^{\mrm{op}})
\ar[d]^{\opn{Rest}}
\\
\dcat{D}(A' \ot_{} B'^{\, \mrm{op}}) \times \dcat{D}(B' \ot C'^{\, \mrm{op}})
\ar[r]^(0.6){ (- \ot^{\mrm{L}}_{B'} -) }
&
\dcat{D}(A' \ot_{} C'^{\, \mrm{op}})
} \]
is commutative up to an isomorphism of triangulated bifunctors. 

\item Suppose $D$ is another DG ring. Then there is an isomorphism  
\[ ((- \ot^{\mrm{L}}_{B} -) \ot^{\mrm{L}}_{C} -) \cong 
(- \ot^{\mrm{L}}_{B} (- \ot^{\mrm{L}}_{C} -)) \]
of triangulated trifunctors 
\[ \dcat{D}(A \ot_{} B^{\mrm{op}}) \times \dcat{D}(B \ot_{} C^{\mrm{op}}) 
\times 
\dcat{D}(C \ot_{} D^{\mrm{op}}) \to \dcat{D}(A \ot_{} D^{\mrm{op}}) . \]
\end{enumerate}
\end{prop}

\begin{proof} \mbox{}

\smallskip \noindent
(1) The DG bifunctor $(- \ot_{B} -)$ induces a triangulated bifunctor 
\[ (- \ot_{B} -) :  \dcat{K}(A \ot_{} B^{\mrm{op}}) \times 
\dcat{K}(B \ot_{} C^{\mrm{op}}) \to \dcat{K}(A \ot_{} C^{\mrm{op}})  \]
on the homotopy categories, in the obvious way. 
By Corollary \ref{cor:1580}, Proposition \ref{prop:1525} and Lemma 
\ref{lem:3445}(1), every DG bimodule $N \in \dcat{C}(B \ot_{} C^{\mrm{op}})$
admits a quasi-isomorphism 
$\rho_N : P \to N$, where 
$P \in \dcat{C}(B \ot_{} C^{\mrm{op}})$
is K-flat over $B$.
Given $M \in \dcat{D}(A \ot_{} B^{\mrm{op}})$, let us define the object
\[ M \ot^{\mrm{L}}_{B} N := M \ot_{B} P \in
\dcat{D}(A \ot_{} C^{\mrm{op}}) , \]
with the morphism 
\[ \eta^{\mrm{L}}_{M, N} := \opn{id}_M \ot_B \, \rho_N : 
M \ot_B^{\mrm{L}} N \to M \ot_{B} N . \]
The pair 
$\bigl( (- \ot^{\mrm{L}}_{B} -), \eta^{\mrm{L}} \bigr)$
is a left derived bifunctor of $(- \ot_{B} -)$. 

\medskip \noindent 
(2) Under either assumption the homomorphism 
$\opn{id}_M \ot_B \, \rho_N$
is a quasi-iso\-morphism; cf.\ Lemma \ref{lem:2109}.

\medskip \noindent 
(3) This is similar to the proof of item (2) of Theorem \ref{thm:2363}. 
Take $M \in \dcat{C}(A \ot B^{\mrm{op}})$ and 
$N \in \dcat{C}(B \ot C^{\mrm{op}})$. 
Let $\rho_N : P \to N$ be the resolution from item (1).
Choose a resolution 
$\rho'_P : P' \to P$ in $\dcat{C}_{\mrm{str}}(B' \ot C'^{\, \mrm{op}})$
where $P'$ is K-flat over $B'$. Then we have a canonical isomorphism 
\[ \opn{Rest}(M) \ot^{\mrm{L}}_{B'} \opn{Rest}(N) \cong M \ot_{B'} P' \]
in $\dcat{D}(A' \ot_{} C'^{\, \mrm{op}})$.
Now in the commutative diagram 
\[ \UseTips \xymatrix @C=10ex @R=6ex {
B' \ot_{B'} P'
\ar[r]^(0.6){\cong}
\ar[d]_{f \ot \, \opn{id}_{P'}}
&
P'
\ar[d]^{\rho'_P}
\\
B \ot_{B'} P'
\ar[r]^(0.6){\opn{id}_B \ot \, \rho'_P}
&
P
} \]
in $\dcat{C}_{\mrm{str}}(B')$ the vertical arrows are quasi-isomorphisms, 
because $P'$ is K-flat over $B'$. Therefore the bottom arrow 
$\opn{id}_B \ot \, \rho'_P$ is a quasi-isomorphism. 

We now look at this commutative diagram 
\[ \UseTips \xymatrix @C=6ex @R=6ex {
M \ot_{B'} P'
\ar@(ur,ul)[rr]^{\opn{id}_{M} \ot \, \rho'_P}
\ar[r]^(0.4){\cong}
&
M \ot_B B \ot_{B'} P'
\ar[r]^(0.6){\phi}
&
M \ot_B P
} \]
in $\dcat{C}_{\mrm{str}}(A' \ot_{} C'^{\, \mrm{op}})$, where 
$\phi := \opn{id}_M \ot \opn{id}_B \ot \, \rho'_P$. 
Because both $P$ and $B \ot_{B'} P'$ are K-flat over $B$, and 
$\opn{id}_B \ot \, \rho'_P$ is a quasi-isomorphism, 
Lemma \ref{lem:2109} tells us that $\phi$ is a quasi-isomorphism. 
Hence $\opn{id}_{M} \ot \, \rho'_P$ is a quasi-isomorphism. 
We get the desired canonical isomorphism 
\[ \opn{Q}(\opn{id}_{M} \ot \, \rho'_P) : 
\opn{Rest}(M) \ot^{\mrm{L}}_{B'} \opn{Rest}(N) \iso 
\opn{Rest}(M \ot^{\mrm{L}}_{B} N) \]
in $\dcat{D}(A' \ot_{} C'^{\, \mrm{op}})$.
 
\medskip \noindent 
(4) Given $M, N, P$ as above and 
$L \in \dcat{C}(C \ot_{} D^{\mrm{op}})$, 
we choose a quasi-iso\-morphism 
$\rho_L : Q \to L$, where 
$Q \in \dcat{C}(C \ot_{} D^{\mrm{op}})$
is K-flat over $C$. A small calculation shows that 
$P \ot_C Q$ is K-flat over $B$. The desired isomorphism 
\[ (M \ot^{\mrm{L}}_{B} N) \ot^{\mrm{L}}_{C} L \cong  
M \ot^{\mrm{L}}_{B} (N \ot^{\mrm{L}}_{C} L) \]
in $\dcat{D}(A \ot_{} D^{\mrm{op}})$
comes from the obvious isomorphism 
\[ (M \ot_{B} P) \ot_{C} Q \cong  
M \ot_{B} (P \ot_{C} Q) \]
in $\dcat{C}_{\mrm{str}}(A \ot_{} D^{\mrm{op}})$.
\end{proof}

Here is some notation: suppose we are given morphisms 
$\phi : M' \to M$ in 
$\dcat{D}(A \ot B^{\mrm{op}})$
and $\psi : N' \to N$ in $\dcat{D}(B \ot C^{\mrm{op}})$.
The result of applying the bifunctor $(- \ot_B^{\mrm{L}} -)$ is the morphism
\begin{equation} \label{eqn:3475}  
 \phi \ot^{\mrm{L}}_B \psi : M' \ot_B^{\mrm{L}} N' \to M \ot_B^{\mrm{L}} N
\end{equation}
in $\dcat{D}(A \ot C^{\mrm{op}})$. 

\begin{prop} \label{prop:1031}
Let $A$, $B$ and $C$ be DG rings. 
\begin{enumerate}
\item The DG bifunctor 
\[ \opn{Hom}_B(-, -) :  \dcat{C}(B \ot A^{\mrm{op}})^{\mrm{op}} \times 
\dcat{C}(B \ot C^{\mrm{op}}) \to \dcat{C}(A \ot C^{\mrm{op}})  \]
has a triangulated right derived bifunctor 
\[ \bigl( \opn{RHom}_B(-, -), \eta^{\mrm{R}} \bigr) :  
\dcat{D}(B \ot A^{\mrm{op}})^{\mrm{op}} \times 
\dcat{D}(B \ot C^{\mrm{op}}) \to \dcat{D}(A \ot C^{\mrm{op}}) . \]

\item Given $M \in \dcat{C}(B \ot A^{\mrm{op}})$
and $N \in \dcat{C}(B \ot C^{\mrm{op}})$,
such that either $M$ is K-projective over $B$ or $N$ is K-injective over $B$, 
the morphism 
\[ \eta^{\mrm{R}}_{M, N} : \opn{Hom}_B(M, N) \to \opn{RHom}_B(M, N) \]
in $\dcat{D}(A \ot C^{\mrm{op}})$ is an isomorphism. 

\item Suppose we are given DG ring homomorphisms 
$A' \to A$, $f : B' \to B$ and $C' \to C$, such that $f$ is a 
quasi-isomorphism. 
Then the diagram 
\[ \UseTips \xymatrix @C=14ex @R=6ex {
\dcat{D}(B \ot A^{\mrm{op}})^{\mrm{op}} \times \dcat{D}(B \ot C^{\mrm{op}})
\ar[r]^(0.6){ \opn{RHom}_B(-, -) }
\ar[d]_{\opn{Rest} \times \opn{Rest}}
&
\dcat{D}(A \ot C^{\mrm{op}})
\ar[d]^{\opn{Rest}}
\\
\dcat{D}(B' \ot A'^{\, \mrm{op}})^{\mrm{op}} \times 
\dcat{D}(B' \ot C'^{\, \mrm{op}})
\ar[r]^(0.6){ \opn{RHom}_{B'}(-, -) }
&
\dcat{D}(A' \ot C'^{\, \mrm{op}})
} \]
is commutative up to an isomorphism of triangulated bifunctors. 
\end{enumerate}
\end{prop}

\begin{exer} \label{exer:3560}
Prove Proposition \ref{prop:1031}. 
Hints: this is similar to the proof of items (1)-(3) of Proposition 
\ref{prop:3450}, but now we rely on Corollary \ref{cor:1665} and Lemma 
\ref{lem:3445}(2) for the existence of resolutions $\rho_N : N \to I$ 
in $\dcat{C}_{\mrm{str}}(B \ot C^{\mrm{op}})$ by DG bimodules $I$
that are K-injective over $B$. For item (3) use Lemma \ref{lem:2000}.
\end{exer}

Again there is related notation: suppose we are given morphisms
$\phi : M \to M'$ in $\dcat{D}(B \ot A^{\mrm{op}})$
and $\psi : N' \to N$ in $\dcat{D}(B \ot C^{\mrm{op}})$.
The result of applying the bifunctor $\opn{RHom}_B(-, -)$
is the morphism
\begin{equation} \label{eqn:3476}  
\opn{RHom}_B(\phi, \psi) : \opn{RHom}_B(M', N') \to \opn{RHom}_B(M, N) 
\end{equation}
in $\dcat{D}(A \ot C^{\mrm{op}})$.

\begin{prop} \label{prop:3550}
Let $A$, $B$, $C$ and $D$ be DG rings. For 
$M \in \dcat{D}(B \ot A^{\mrm{op}})$,
$N \in \dcat{D}(C \ot B^{\mrm{op}})$ and 
$L \in \dcat{D}(C \ot D^{\mrm{op}})$
there is an isomorphism 
\[ \opn{RHom}_B \bigl( M, \opn{RHom}_C(N, L) \bigr) \cong 
\opn{RHom}_C(N \ot_B^{\mrm{L}} M, L) \]
in $\dcat{D}(A \ot D^{\mrm{op}})$ called {\em derived Hom-tensor adjunction}. 
This isomorphism is functorial in the objects $M, N, L$.
\end{prop}

\begin{proof}
We choose a quasi-isomorphism $L \to J$ in 
$\dcat{C}_{\mrm{str}}(C \ot D^{\mrm{op}})$ into a DG module $J$ that is 
K-injective over $C$, and a quasi-isomorphism $Q \to N$ in 
$\dcat{C}_{\mrm{str}}(C \ot B^{\mrm{op}})$ from a DG module $Q$ that is K-flat 
over $B^{\mrm{op}}$. A calculation shows that the DG bimodule 
$\opn{Hom}_C(Q, J)$ is K-injective over $B$. 
The desired isomorphism comes from the obvious adjunction isomorphism 
\[ \opn{Hom}_B(M, \opn{Hom}_C(Q, J)) \cong 
\opn{Hom}_C(Q \ot_B M, J) \]
in $\dcat{C}_{\mrm{str}}(A \ot D^{\mrm{op}})$. 
\end{proof}

\begin{prop} \label{prop:3565}
Let $A$, $B$, $C$ and $D$ be DG rings. For 
$L \in \dcat{D}(A \ot C^{\mrm{op}})$,
$M \in \dcat{D}(A \ot B^{\mrm{op}})$ and 
$N \in \dcat{D}(B \ot D^{\mrm{op}})$
there is a morphism 
\[ \opn{ev}^{\mrm{R, L}}_{L, M, N} :
\opn{RHom}_A(L, M) \ot_B^{\mrm{L}} N \to 
\opn{RHom}_A(L, M \ot_B^{\mrm{L}} N) \]
in $\dcat{D}(C \ot D^{\mrm{op}})$, called {\em derived tensor-evaluation}. 
This morphism is functorial in the objects $L, M, N$.
Moreover, after applying the restriction functor  
$\dcat{D}(C \ot D^{\mrm{op}}) \to \dcat{D}(\K)$, 
the morphism $\opn{ev}^{\mrm{R, L}}_{L, M, N}$ coincides with the morphism from 
Theorem \tup{\ref{thm:4320}}. 
\end{prop}

The reason we can give stronger statement here, as compared to 
Theorem \ref{thm:4320}, is because here our DG rings are K-flat 
over $\K$. 

\begin{proof}
Choose a quasi-isomorphism $Q \to N$ in 
$\dcat{C}_{\mrm{str}}(B \ot D^{\mrm{op}})$ from a DG module $Q$ that is 
K-flat over $B$, a quasi-isomorphism $M \to I$ in 
$\dcat{C}_{\mrm{str}}(A \ot B^{\mrm{op}})$ into a DG module $I$ that is 
K-injective over $A$, and a quasi-isomorphism 
$\rho : I \ot_B Q \to J$ in $\dcat{C}_{\mrm{str}}(A \ot D^{\mrm{op}})$ into a 
DG module $J$ that is K-injective over $A$. Then 
$\opn{ev}^{\mrm{R, L}}_{L, M, N}$
is represented by the composed homomorphism 
\[  \opn{Hom}_A(L, I) \ot_B Q \xar{ \opn{ev}_{L, I, Q} }  
\opn{Hom}_A(L, I \ot_B Q) \xar{\til{\rho}} \opn{Hom}_A(L, J)  \]
in $\dcat{C}_{\mrm{str}}(C \ot D^{\mrm{op}})$,
where $\til{\rho} := \opn{Hom}(\opn{id}_L, \rho)$. 

When we forget $C$ and $D$, then we can choose a K-projective resolution 
$P \to L$ in $\dcat{C}_{\mrm{str}}(A)$. Then we have a commutative diagram 
\[ \UseTips \xymatrix @C=8ex @R=6ex {
\opn{Hom}_A(P, I) \ot_B Q
\ar[r]^{ \opn{ev}_{P, I, Q} }
&
\opn{Hom}_A(P, I \ot_B Q) 
\ar[r]^(0.55){\mrm{q.i.}}
&
\opn{Hom}_A(P, J)
\\
\opn{Hom}_A(L, I) \ot_B Q
\ar[r]^{ \opn{ev}_{L, I, Q} }
\ar[u]^{\mrm{q.i.}}
&
\opn{Hom}_A(L, I \ot_B Q) 
\ar[r]^(0.55){\til{\rho}}
\ar[u]
&
\opn{Hom}_A(L, J)
\ar[u]^{\mrm{q.i.}}
} \]
in $\dcat{C}_{\mrm{str}}(\K)$, in which the arrows marked ``q.i.'' are 
quasi-isomorphisms. By comparing to the proof of Theorem \ref{thm:4320}
this proves that last assertion.
\end{proof}

According to Propositions \ref{prop:3450}
and \ref{prop:1031} there are triangulated bifunctors 
\begin{equation} \label{eqn:3477}
(- \ot^{\mrm{L}}_{A} -) : \dcat{D}(A^{\mrm{en}}) \times 
\dcat{D}(A^{\mrm{en}}) \to \dcat{D}(A^{\mrm{en}}) ,
\end{equation}
\begin{equation} \label{eqn:3478}
(- \ot^{\mrm{L}}_{A} -) : \dcat{D}(A^{\mrm{en}}) \times 
\dcat{D}(A \ot B^{\mrm{op}}) \to \dcat{D}(A \ot B^{\mrm{op}}) 
\end{equation}
and
\begin{equation} \label{eqn:1479}
\opn{RHom}_A(-, -) : \dcat{D}(A^{\mrm{en}})^{\mrm{op}} \times 
\dcat{D}(A \ot B^{\mrm{op}}) \to \dcat{D}(A \ot B^{\mrm{op}}) .
\end{equation}

\begin{rem} \label{rem:3450}
The category $\dcat{D}(A^{\mrm{en}})$, with the tensor operation 
$(- \ot^{\mrm{L}}_{A} -)$, is a {\em monoidal category}.
Moreover, $\dcat{D}(A^{\mrm{en}})$ is a {\em biclosed monoidal category}: the 
two internal Hom operations are 
$\opn{RHom}_A(-, -)$ and $\opn{RHom}_{A^{\mrm{op}}}(-, -)$.
If $A$ is weakly commutative, then $\dcat{D}(A^{\mrm{en}})$ is a symmetric 
monoidal category. See \cite[Section 5]{VyYe}, \cite{Mac2} and \cite{nLab}. 
\end{rem}

\begin{rem} \label{rem:1420}
Here is an outline of our way to handle derived categories of DG bimodules in 
the absence of flatness. The idea is to choose K-flat resolutions 
$\til{A} \to A$ and $\til{B} \to B$ in the category 
$\catt{DGRng} \centover \K$
of DG central $\K$-rings. This can be done by 
Theorem \ref{thm:4310} and Proposition \ref{prop:4305}. 
Then the {\em derived category of DG $A$-$B$-bimodules} 
\index{Derived category! of DG bimodules}
is the triangulated 
category $\dcat{D}(\til{A} \ot_{} \til{B}^{\mrm{op}})$.
Note that the restriction functors 
$\dcat{D}(\til{A}) \to \dcat{D}(A)$ and 
$\dcat{D}(\til{B}^{\mrm{op}}) \to \dcat{D}(B^{\mrm{op}})$ 
are equivalences. This says that we have triangulated bifunctors 
\[ (- \ot^{\mrm{L}}_{\til{B}} -) : 
\dcat{D}(\til{A} \ot_{} \til{B}^{\mrm{op}}) \times 
\dcat{D}(B) \to \dcat{D}(A) \]
and
\[ \opn{RHom}_{\til{A}}(-, -) : 
\dcat{D}(\til{A} \ot_{} \til{B}^{\mrm{op}})^{\mrm{op}} \times 
\dcat{D}(A) \to \dcat{D}(B) . \]
We also have the ``tilde'' versions of the functors 
(\ref {eqn:3477}), (\ref{eqn:3478}) and (\ref{eqn:1479}). 
These are collectively called the {\em package of standard derived functors}. 

The triangulated category 
$\dcat{D}(\til{A} \ot_{} \til{B}^{\mrm{op}})$
is independent of the resolutions, up to a canonical equivalence.
The argument is this. 
Suppose we are given K-flat resolutions 
$\til{A}_i \to A$ and $\til{B}_i \to B$ in the category 
$\catt{DGRng} \centover \K$, indexed by $i \in I$ for some set $I$. 
For each $i, j \in I$, the 
the DG module $A \ot^{\mrm{L}}_{\K} B$
has an outside action by the DG ring 
$\til{A}_j \ot_{} \til{B}_j^{\mrm{op}}$
and an inside action the DG ring 
$\til{A}_i \ot_{} \til{B}_i^{\mrm{op}}$. 
(See Subsection \ref{subsec:RNCDC-uniq} regarding outside and inside actions.)
These give rise to a DG bimodule
\[ T_{i, j} := A \ot^{\mrm{L}}_{\K} B \in 
\dcat{D} \bigl( (\til{A}_j \ot_{} \til{B}_j^{\mrm{op}}) \ot_{}
(\til{A}_i \ot_{} \til{B}_i^{\mrm{op}})^{\mrm{op}} \bigr) . \]
It turns out that $T_{i, j}$ is a {\em tilting DG bimodule}, as defined in the 
next subsection; and hence the triangulated functor
\begin{equation} \label{eqn:5090}
G_{i, j} := T_{i, j} \ot^{\mrm{L}}_{\til{A}_i \ot \til{B}_i^{\mrm{op}}} (-) :
\dcat{D}(\til{A}_i \ot \til{B}_i^{\mrm{op}}) \to
\dcat{D}(\til{A}_j \ot \til{B}_j^{\mrm{op}}) 
\end{equation}
is an equivalence. For a third index $k \in I$ there is a canonical 
isomorphism
\[ T_{j, k} \ot^{\mrm{L}}_{\til{A}_j \ot_{} \til{B}_j^{\mrm{op}}} 
T_{i, j} \cong T_{i, k} 
\quad \tup{in} \quad 
\dcat{D} \bigl( (\til{A}_k \ot_{} \til{B}_k^{\mrm{op}}) \ot_{}
(\til{A}_i \ot_{} \til{B}_i^{\mrm{op}})^{\mrm{op}} \bigr) . \]
Thus there are isomorphisms of triangulated functors 
\begin{equation} \label{eqn:3501}
G_{j, k} \circ G_{i, j} \iso G_{i, k} ,
\end{equation}
and they satisfy the pentagon axiom. 
The equivalences $G_{i, j}$ respect the packages of standard derived 
functors.

We see that there is a $\K$-linear triangulated category that we will 
symbolically denote by 
$\dcat{D}(A \ot^{\mrm{L}}_{\K} B^{\mrm{op}})$,
and it is canonically equivalent to all the triangulated categories 
$\dcat{D}(\til{A}_i \ot_{} \til{B}_i^{\mrm{op}})$.
For more details see the lecture notes \cite{Ye16} 
or the paper \cite{Ye14}. 
\end{rem}

\begin{rem} \label{rem:1421}
Actually, we can say more. 
Recall the $2$-category $\catt{TrCat} / \K$ of $\K$-linear triangulated 
categories that was briefly mentioned at the end of Subsection 
\ref{subsec:2-cat-not}. 
There is a pseudofunctor 
\begin{equation} \label{eqn:3502}
\opn{DerCat} : \catt{DGRng} \centover \K \to \catt{TrCat} / \K
\end{equation}
that sends a DG ring $A$ to the triangulated category
$\opn{DerCat}(A) := \dcat{D}(A)$,
and a homomorphism $f : A \to B$ in $\catt{DGRng} \centover \K$
is sent to the triangulated functor 
\[ \opn{DerCat}(f) := \opn{LInd}_f = B \ot^{\mrm{L}}_{A} (-)
: \dcat{D}(A) \to \dcat{D}(B) . \]
Note that the category $\catt{DGRng} \centover \K$ is not linear (the morphisms 
sets have no linear structure), whereas the 
$2$-category $\catt{TrCat} / \K$ is linear (in the sense that the sets of 
$2$-morphisms are $\K$-modules). 

Define $\dcat{D}(\catt{DGRng} \centover \K)$
to be the categorical localization of the category  
$\catt{DGRng} \centover \K$ with respect to the quasi-isomorphisms in it. Then 
the pseudofunctor $\opn{DerCat}$ from (\ref{eqn:3502}) localizes, in the 
categorical sense, to a pseudofunctor
\[ \opn{DerCat} : \dcat{D}(\catt{DGRng} \centover \K) \to \catt{TrCat} / \K . \]

The (nonlinear) bifunctor 
\[ (- \ot_{\K} -) : 
(\catt{DGRng} \centover \K) \times (\catt{DGRng} \centover \K) \to
\catt{DGRng} \centover \K \]
has a left derived bifunctor 
\[ (- \ot^{\mrm{L}}_{\K} -) : 
\dcat{D}(\catt{DGRng} \centover \K) \times 
\dcat{D}(\catt{DGRng} \centover \K) \to 
\dcat{D}(\catt{DGRng} \centover \K) . \]
Given a pair of DG rings $A$ and $B$, we thus have DG ring 
$A \ot^{\mrm{L}}_{\K} B^{\mrm{op}}$, 
well-defined up to a unique isomorphism in 
$\dcat{D}(\catt{DGRng} \centover \K)$. 
The derived category of DG $A$-$B$-bimodules described at the end of Remark 
\ref{rem:1420} is canonically equivalent to 
$\opn{DerCat}(A \ot^{\mrm{L}}_{\K} B^{\mrm{op}})$
in $\catt{TrCat} / \K$. 
\end{rem}




\mysubsection{Tilting DG Bimodules} \label{subsec:tilting}

We continue with Convention \ref{conv:3445}. In particular all DG rings are 
K-flat central over the base ring $\K$, $\ot = \ot_{\K}$, 
and $A^{\mrm{en}} = A \ot A^{\mrm{op}}$ for a DG ring $A$. 

\begin{dfn} \label{dfn:3455}
Let $A$ and $B$ be DG rings.
An object $T \in \dcat{D}(B \ot_{} A^{\mrm{op}})$
is called a {\em tilting DG $B$-$A$-bimodule} 
\index{Tilting! DG bimodule}
if there exists some object $S \in \dcat{D}(A \ot_{} B^{\mrm{op}})$,
and isomorphisms 
$S \ot^{\mrm{L}}_B T \cong A$ in $\dcat{D}(A^{\mrm{en}})$
and
$T \ot^{\mrm{L}}_A S \cong B$ in $\dcat{D}(B^{\mrm{en}})$.

When $B = A$, so that $B \ot_{} A^{\mrm{op}} = A^{\mrm{en}}$,
we use the term {\em tilting DG $A$-bimodule} as an abbreviation for the term
tilting DG $A$-$A$-bimodule. 
\end{dfn}

It is clear from the symmetry of the definition that the object 
$S \in \dcat{D}(A \ot_{} B^{\mrm{op}})$ is a 
tilting DG $A$-$B$-bimodule. 

\begin{lem} \label{lem:3455}
Let $T \in \dcat{D}(B \ot A^{\mrm{op}})$ and 
$S, S' \in \dcat{D}(A \ot B^{\mrm{op}})$
satisfy 
$S \ot^{\mrm{L}}_B T \cong A$ in $\dcat{D}(A^{\mrm{en}})$ 
and
$T \ot^{\mrm{L}}_A S' \cong B$ in $\dcat{D}(B^{\mrm{en}})$.
Then 
$S' \cong S$ in $\dcat{D}(A \ot_{} B^{\mrm{op}})$.
Therefore $T$ is a tilting DG $B$-$A$-bimodule, and $S$ is a tilting 
DG $A$-$B$-bimodule. 
\end{lem}

\begin{proof}
Using the associativity of derived tensor products from Proposition \lb
\ref{prop:3450}(4)
we have isomorphisms 
\[ S \cong S \ot^{\mrm{L}}_B B \cong
S \ot^{\mrm{L}}_B (T \ot^{\mrm{L}}_A S') \cong 
(S \ot^{\mrm{L}}_B T) \ot^{\mrm{L}}_A S' \cong 
A \ot^{\mrm{L}}_A S' \cong S' \]
in $\dcat{D}(A \ot_{} B^{\mrm{op}})$. 
\end{proof}

The lemma implies that the tilting DG $A$-$B$-bimodule $S$ in Definition 
\ref{dfn:3455} is unique up to isomorphism.

\begin{dfn} \label{3460}
The tilting DG $A$-$B$-bimodule $S$ in Definition \ref{dfn:3455} is called 
the {\em quasi-inverse of $T$}. 
\index{Differential graded bimodule! quasi-inverse of tilting}
\end{dfn}

Here are some general properties of tilting DG bimodules. 

\begin{prop} \label{prop:3490}
Let $A, B, C$ be DG rings, let $T$ be a tilting DG $B$-$A$-bimodule, and let 
$S$ be a tilting DG $C$-$B$-bimodule. 
Then $S \ot^{\mrm{L}}_{B} T$ is a tilting $C$-$A$-bimodule. 
\end{prop}

\begin{proof}
Let $T^{\vee}$ be the quasi-inverse of $T$, and let 
$S^{\vee}$ be the quasi-inverse of $S$. Then
\[ \begin{aligned}
& (T^{\vee} \ot^{\mrm{L}}_{B} S^{\vee}) \ot^{\mrm{L}}_{C} 
(S \ot^{\mrm{L}}_{B} T) \cong 
T^{\vee} \ot^{\mrm{L}}_{B} (S^{\vee} \ot^{\mrm{L}}_{C} S )\ot^{\mrm{L}}_{B} T
\\ & \quad 
\cong T^{\vee} \ot^{\mrm{L}}_{B} B \ot^{\mrm{L}}_{B} T \cong A 
\end{aligned} \]
in $\dcat{D}(A^{\mrm{en}})$. Likewise 
\[ (S \ot^{\mrm{L}}_{B} T) \ot^{\mrm{L}}_{A} 
(T^{\vee} \ot^{\mrm{L}}_{B} S^{\vee}) \cong C \]
in $\dcat{D}(C^{\mrm{en}})$. 
\end{proof}

\begin{prop} \label{prop:3455}
Suppose $T$ is a tilting DG $B$-$A$-bimodule, and 
$S$ is a tilting DG $A$-$B$-bimodule. Then the triangulated functor 
\[ G_{T, S} : \dcat{D}(A^{\mrm{en}}) \to \dcat{D}(B^{\mrm{en}}) \]
with formula
$G_{T, S}(M) := T \ot^{\mrm{L}}_{A} M \ot^{\mrm{L}}_{A} S$ 
is an equivalence. 
\end{prop}

\begin{prop} \label{prop:3525}
Suppose $T$ is a tilting DG $B$-$A$-bimodule, with quasi-inverse 
$T^{\vee}$. Fix an isomorphism 
$T^{\vee} \ot^{\mrm{L}}_{B} T \iso A$
in $\dcat{D}(A^{\mrm{en}})$. 
Then for every $M_1, M_2 \in \dcat{D}(A^{\mrm{en}})$ there is an  
isomorphism 
\[ G_{T, T^{\vee}} (M_1 \ot^{\mrm{L}}_{A} M_2) \iso 
G_{T, T^{\vee}} (M_1) \ot^{\mrm{L}}_{B} G_{T, T^{\vee}} (M_2) \]
in $\dcat{D}(B^{\mrm{en}})$. This isomorphism is bifunctorial in
$(M_1, M_2)$. 
\end{prop}

\begin{exer} \label{exer:3590}
Prove Propositions \ref{prop:3455} and \ref{prop:3525}.
\end{exer}

\begin{dfn} \label{dfn:3457}
Let $A$ be a K-flat central DG $\K$-ring. The 
{\em noncommutative derived Picard group of $A$ relative to $\K$} 
\index{Derived Picard group}
\index{Picard group! noncommutative derived}
\index{1-DPic(A)@$\opn{DPic}_{\K}(A)$}
is the group $\opn{DPic}_{\K}(A)$, whose 
elements are the isomorphism classes in $\dcat{D}(A^{\mrm{en}})$ of the 
tilting DG $A$-bimodules. The multiplication in this group is induced from 
$(- \ot^{\mrm{L}}_{A} -)$, and the unit element is the class of the DG bimodule 
$A$. 
\end{dfn}

This definition is valid by Proposition \ref{prop:3490}. 
Propositions \ref{prop:3455} and \ref{prop:3525} say that: 

\begin{cor} \label{cor:3525}
A tilting DG $B$-$A$-bimodule $T$, with quasi-inverse $T^{\vee}$, 
and an isomorphism 
$T^{\vee} \ot^{\mrm{L}}_{B} T \cong A$ in $\dcat{D}(A^{\mrm{en}})$,
induce a group isomorphism 
\[ \opn{DPic}_{\K}(A) \iso \opn{DPic}_{\K}(B) . \]
\end{cor}

We don't know much about the structure of the derived Picard group 
$\opn{DPic}_{\K}(A)$ in general. However, when $A$ is either a ring
or a commutative DG ring, there are a few strong structural results -- see 
next subsection. 

\begin{rem} \label{rem:3595}
If $A$ is not K-flat over the base ring $\K$, then the derived Picard group 
should be defined using a K-flat resolution $\til{A} \to A$, as explained in 
Remark \ref{rem:1420}. Namely $\opn{DPic}_{\K}(A)$ is the group whose elements 
are the isomorphism classes in $\dcat{D}(\til{A}^{\mrm{en}})$ of the tilting DG 
$\til{A}$-bimodules, etc. This group is independent, up to a canonical 
group isomorphism, of the resolution $\til{A} \to A$. 
\end{rem}

The derived homothety morphism in the commutative setting,
and the related derived Morita property, were defined in subsection 
\ref{subsec:du-cplxs}. In the current noncommutative setting these notions 
become more involved, as we shall now see.

Given a DG bimodule $M \in \dcat{C}(B \ot_{\K} A^{\mrm{op}})$, 
there are DG ring homomorphisms 
\[  \opn{hm}_{M, A^{\mrm{op}}} : A^{\mrm{op}} \to \opn{End}_{B}(M) = 
\opn{Hom}_{B}(M, M) \]
and
\[  \opn{hm}_{M, B} : B \to 
\opn{End}_{A^{\mrm{op}}}(M) = \opn{Hom}_{A^{\mrm{op}}}(M, M) \]
that we call the {\em noncommutative homothety homomorphisms through 
$A^{\mrm{op}}$ and \lb through $B$} respectively. 
When we forget the ring structures, these become homomorphisms
\begin{equation} \label{eqn:4175}
\opn{hm}_{M, A^{\mrm{op}}} : A \to \opn{Hom}_{B}(M, M) 
\end{equation}
and
\begin{equation} \label{eqn:4176}
\opn{hm}_{M, B} : B \to \opn{Hom}_{A^{\mrm{op}}}(M, M)
\end{equation}
in $\dcat{C}_{\mrm{str}}(A^{\mrm{en}})$ and 
$\dcat{C}_{\mrm{str}}(B^{\mrm{en}})$ respectively. 

\begin{dfn} \label{dfn:3565}
Let $M \in \dcat{D}(B \ot_{} A^{\mrm{op}})$.
\index{Homothety morphism! noncommutative derived}
\begin{enumerate}
\item The {\em noncommutative derived homothety morphism of $M$ through
$A^{\mrm{op}}$} is the morphism 
\[ \opn{hm}^{\mrm{R}}_{M, A^{\mrm{op}}} := \eta^{\mrm{R}}_{M, M} \circ 
\opn{Q}(\opn{hm}_{M, A^{\mrm{op}}}) : A \to \opn{RHom}_{B}(M, M) \]
in $\dcat{D}(A^{\mrm{en}})$.

\item The {\em noncommutative derived homothety morphism of $M$ through $B$} is 
the morphism 
\[ \opn{hm}^{\mrm{R}}_{M, B} := \eta^{\mrm{R}}_{M, M} 
\circ \opn{Q}(\opn{hm}_{M, B}) : B \to \opn{RHom}_{A^{\mrm{op}}}(M, M)  \]
in $\dcat{D}(B^{\mrm{en}})$.
\end{enumerate}
\end{dfn}

Here is a commutative diagram in $\dcat{D}(A^{\mrm{en}})$ depicting the 
noncommutative derived homothety morphism through $A^{\mrm{op}}$.
\[ \UseTips \xymatrix @C=8ex @R=6ex {
A
\ar[d]_{\opn{Q}( \opn{hm}_{M, A^{\mrm{op}}} )}
\ar[dr]^{ \opn{hm}^{\mrm{R}}_{M, A^{\mrm{op}}} }
\\
\opn{Hom}_{B}(M, M)
\ar[r]^{\eta^{\mrm{R}}_{M, M}}
&
\opn{RHom}_{B}(M, M)
} \]
The morphism $\eta^{\mrm{R}}_{M, M}$ is part of the right derived bifunctor 
$\opn{RHom}_{B}(-, -)$; see Proposition \ref{prop:1031}. 

\begin{dfn} \label{3455}
Let $A$ and $B$ be DG rings, and let $M$ be an object of 
$\dcat{D}(B \ot_{} A^{\mrm{op}})$.
\index{Derived Morita property! noncommutative}
\begin{enumerate}
\item We say that $M$ has the {\em noncommutative derived Morita property on 
the $B$ side} if the derived homothety morphism
\[ \opn{hm}^{\mrm{R}}_{M, A^{\mrm{op}}} : A \to \opn{RHom}_{B}(M, M) \]
in $\dcat{D}(A^{\mrm{en}})$ is an isomorphism. 

\item We say that $M$ has the {\em noncommutative derived Morita property on 
the $A^{\mrm{op}}$ side} if the derived homothety morphism
\[ \opn{hm}^{\mrm{R}}_{M, B} : B \to 
\opn{RHom}_{A^{\mrm{op}}}(M, M)  \]
in $\dcat{D}(B^{\mrm{en}})$ is an isomorphism. 

\item We say that $M$ has the {\em noncommutative derived Morita property on 
both sides} if it has the noncommutative derived Morita property on 
the $B$ side and on the $A^{\mrm{op}}$ side.
\end{enumerate}
\end{dfn}

Recall the restriction functors from diagram (\ref{eqn:3536}). 
Compact generators were introduced in Definition \ref{dfn:3431}(3). 
The next definition resembles Definition \lb \ref{dfn:3465}.

\begin{dfn} \label{3505} 
Let $A$ and $B$ be DG rings, and let $T$ be an object of 
$\dcat{D}(B \ot_{} A^{\mrm{op}})$.
\index{Differential graded module! generator}
\index{Differential graded module! compact}
\begin{enumerate}
\item We say that $T$ is a {\em compact generator of $\dcat{D}(B)$},
or a  {\em compact generator on the $B$ side}, 
if $\opn{Rest}_B(T) \in \dcat{D}(B)$ is a compact generator. 

\item We say that $T$ is a {\em compact generator of 
$\dcat{D}(A^{\mrm{op}})$},
or a  {\em compact generator on the $A^{\mrm{op}}$ side}, if 
$\opn{Rest}_{A^{\mrm{op}}}(T) \in \dcat{D}(A^{\mrm{op}})$ is a compact 
generator. 

\item We say that $T$ is a {\em compact generator on both sides}
if it is a compact generator on the $B$ side and on the $A^{\mrm{op}}$ side.
\end{enumerate}
\end{dfn}

The next theorem is similar to results appearing in \cite{Ric2} and 
\cite{Kel2}. It is a derived version of the classical result for invertible 
bimodules over a ring (see Example \ref{exa:3435} or \cite[Section 4.1]{Row}). 

In Definition \ref{dfn:3430} we introduced the notion of pretilting DG 
bimodules, and in Definition \ref{dfn:3455} we introduced tilting DG bimodules. 

\begin{thm} \label{thm:3455}
Let $A$ and $B$ be K-flat DG central $\K$-rings. The following three 
conditions are equivalent for an object 
$T \in \dcat{D}(B \ot_{} A^{\mrm{op}})$. 
\begin{enumerate}
\rmitem{i} The DG $B$-$A$-bimodule $T$ is tilting.
\index{Differential graded bimodule! tilting}

\rmitem{ii} The DG $B$-$A$-bimodule $T$ is pretilting.
\index{Differential graded bimodule! pretilting}

\rmitem{iii} The DG $B$-$A$-bimodule $T$ is a compact generator on the 
$B$ side, and it has the noncommutative derived Morita property on the $B$ side.
\index{Differential graded module! compact}
\index{Differential graded module! generator}
\index{Derived Morita property! noncommutative}
\end{enumerate} 
\end{thm}

\begin{proof} \mbox{} 

\smallskip \noindent
(i) $\Rightarrow$ (ii): Let $S \in \dcat{D}(A \ot_{} B^{\mrm{op}})$ be the 
quasi-inverse of $T$. The functor 
\begin{equation} \label{eqn:3505}
G_T := T \ot^{\mrm{L}}_{A} (-) : \dcat{D}(A) \to \dcat{D}(B)
\end{equation}
has a quasi-inverse 
\begin{equation} \label{eqn:3513}
G_S := S \ot^{\mrm{L}}_{B} (-) : \dcat{D}(B) \to \dcat{D}(A) . 
\end{equation}
So $G_T$ is an equivalence, and the DG bimodule $T$ is pretilting. 

\medskip \noindent
(ii) $\Rightarrow$ (iii): By definition, the functor $G_T$ 
from formula (\ref{eqn:3505}) is an equivalence of triangulated categories. 
Also $G_T(A) \cong T$ in $\dcat{D}(B)$. 
Because $A$ is a compact generator of $\dcat{D}(A)$, Proposition 
\ref{prop:3545} says that $T$ is a compact generator of $\dcat{D}(B)$.

Next we shall prove that $T$ has the noncommutative derived Morita property 
on the $B$ side, namely that the derived homothety morphism 
\begin{equation} \label{eqn:3511}
\opn{hm}^{\mrm{R}}_{M, A^{\mrm{op}}} : A \to \opn{RHom}_B(T, T) 
\end{equation}
in $\dcat{D}(A^{\mrm{en}})$ is an isomorphism. 
According to Proposition \ref{prop:3445} the restriction functor
$\opn{Rest}_{A} : \dcat{D}(A^{\mrm{en}}) \to \dcat{D}(A)$ 
is conservative. Therefore, to prove that the morphism 
$\opn{hm}^{\mrm{R}}_{M, A^{\mrm{op}}}$ 
is an isomorphism in $\dcat{D}(A^{\mrm{en}})$, we can forget the 
$A^{\mrm{op}}$-module structures on the DG 
bimodules $A$ and $\opn{RHom}_B(T, T)$, and just prove that 
$\opn{Rest}_{A}(\opn{hm}^{\mrm{R}}_{M, A^{\mrm{op}}})$
is an isomorphism in $\dcat{D}(A)$. 

According to Proposition \ref{prop:3537} the functor 
\begin{equation} \label{eqn:3506}
F_T := \opn{RHom}_B(T, -) : \dcat{D}(B) \to \dcat{D}(A) , 
\end{equation}
which is the right adjoint of $G_T$, is an equivalence.
So there is an isomorphism of triangulated functors
$\ze : \opn{Id}_{\dcat{D}(A)} \iso F_T \circ G_T$
from $\dcat{D}(A)$ to itself. 
There is a diagram 
\[ \UseTips \xymatrix @C=6ex @R=6ex {
A
\ar[d]_{\opn{Rest}_{A}(\opn{hm}^{\mrm{R}}_{M, A^{\mrm{op}}}) }
\ar@(r,ul)[drr]^{\ze_A}_{\cong}
\\
\opn{RHom}_B(T, T)
\ar[r]^(0.45){\cong}
&
\opn{RHom}_B(T, T \ot^{\mrm{L}}_{A} A)
\ar[r]^(0.56){\cong}
&
(F_T \circ G_T)(A)
} \]
in $\dcat{D}(A)$, and a calculation (with elements, using a K-flat resolution 
$P \to T$ in $\dcat{C}_{\mrm{str}}(B \ot A^{\mrm{op}})$) shows that it is 
commutative. Because $\ze_A$ is an isomorphism, we conclude that 
$\opn{Rest}_{A}(\opn{hm}^{\mrm{R}}_{M, A^{\mrm{op}}})$ is an isomorphism. 

\medskip \noindent
(iii) $\Rightarrow$ (ii): 
Choose a resolution $T \to I$ in $\dcat{C}(B \ot_{} A^{\mrm{op}})$
such that $I$ is K-injective over $B$; this is possible by Lemma 
\ref{lem:3445}(2). (If $A$ is K-projective over $\K$ then we can also 
choose a resolution $P \to T$ in $\dcat{C}(B \ot_{} A^{\mrm{op}})$
such that $P$ is K-projective over $B$; and then the proof can proceed with $P$ 
instead of $I$.) Then the noncommutative derived homothety morphism 
$\opn{hm}^{\mrm{R}}_{T, A^{\mrm{op}}}$ in $\dcat{D}(A^{\mrm{en}})$ is 
represented by the canonical DG 
ring homomorphism $g : A \to A_I$, where 
$A_I := \opn{End}_B(I)^{\mrm{op}}$.
The noncommutative derived Morita property on the $B$ side says that 
$g$ is a quasi-isomorphism. Hence, by Theorem \ref{thm:2363}, 
the restriction functor 
$\opn{Rest}_{g} : \dcat{D}(A_I) \to \dcat{D}(A)$
is an equivalence. 

Because $I \cong T$ is a compact generator of $\dcat{D}(B)$, according to 
Theorem \ref{thm:3430} we know that the functor 
$F_I := \opn{RHom}_B(I, -) : \dcat{D}(B) \to \dcat{D}(A_I)$
is an equivalence. Taking $F_T$ to be the functor from (\ref{eqn:3506}),
we have a diagram of functors 
\[ \UseTips \xymatrix @C=8ex @R=6ex {
\dcat{D}(B)
\ar[d]_{F_I}
\ar[dr]^{F_T}
\\
\dcat{D}(A_I)
\ar[r]^{\opn{Rest}_{g}}
&
\dcat{D}(A)
} \]
that is commutative up to isomorphism (cf.\ Theorem \ref{thm:2363}(3)).
Therefore $F_T$ is an equivalence, and thus $T$ is pretilting. 

\medskip \noindent
(ii) $\Rightarrow$ (i): 
We assume that $T$ is a pretilting DG $B$-$A$-bimodule, i.e.\ 
the functors $G_T$ and $F_T$, from (\ref{eqn:3505}) and
(\ref{eqn:3506}) respectively, are equivalences.
Define the DG bimodule 
\begin{equation} \label{eqn:3507}
S := \opn{RHom}_B(T, B) \in \dcat{D}(A \ot_{} B^{\mrm{op}}) .
\end{equation}
Consider the tensor-evaluation morphism 
\begin{equation} \label{eqn:3508}
\opn{ev}^{\mrm{R, L}}_{T, B, T} : \opn{RHom}_B(T, B) 
\ot^{\mrm{L}}_{B} T \to 
\opn{RHom}_B(T, B \ot^{\mrm{L}}_{B} T)
\end{equation}
in $\dcat{D}(A^{\mrm{en}})$ from  Proposition \ref{prop:3565}.
Because the restriction functor
$\opn{Rest}_{\K} : \lb \dcat{D}(A^{\mrm{en}}) \to \dcat{D}(\K)$
is conservative, to prove that $\opn{ev}^{\mrm{R, L}}_{T, B, T}$ is an 
isomorphism in $\dcat{D}(A^{\mrm{en}})$, it suffices 
to prove it for 
$\opn{Rest}_{\K}(\opn{ev}^{\mrm{R, L}}_{T, B, T})$. 
But by Proposition \ref{prop:3431} the DG module $T$ is algebraically perfect 
on the $B$ side; so by Theorem \ref{thm:3400}, the morphism 
$\opn{Rest}_{\K}(\opn{ev}^{\mrm{R, L}}_{T, B, T})$ is an isomorphism.
Thus we get these isomorphisms 
\begin{equation} \label{eqn:3509}
\begin{aligned}
& S \ot^{\mrm{L}}_{B} T = \opn{RHom}_B(T, B) \ot^{\mrm{L}}_{B} T
\\
& \quad \xar{\opn{ev}^{\mrm{R, L}}_{T, B, T}}  
\opn{RHom}_B(T, B \ot^{\mrm{L}}_{B} T)
\cong \opn{RHom}_B(T, T) 
\end{aligned}
\end{equation}
in $\dcat{D}(A^{\mrm{en}})$. 
On the other hand, in the proof of ``(ii) $\Rightarrow$ (iii)'' above we 
already showed (formula (\ref{eqn:3511})) that 
$A \cong \opn{RHom}_B(T, T)$ in $\dcat{D}(A^{\mrm{en}})$. 
Combining this with (\ref{eqn:3509}) we deduce that 
\begin{equation} \label{eqn:3512}
S \ot^{\mrm{L}}_{B} T \cong A \ \ \tup{in} \ \ \dcat{D}(A^{\mrm{en}}) . 
\end{equation}

Next, consider the functor $G_S$ from formula (\ref{eqn:3513}), where now $S$ 
is the DG bimodule from (\ref{eqn:3507}). The same arguments as above -- those 
used for $\opn{ev}^{\mrm{R, L}}_{T, B, T}$ -- show that for every 
$N \in \dcat{D}(B)$ the morphism 
\[ \opn{ev}^{\mrm{R, L}}_{T, B, N} : \opn{RHom}_B(T, B) \ot^{\mrm{L}}_{B} N \to 
\opn{RHom}_B(T, B \ot^{\mrm{L}}_{B} N) \]
in $\dcat{D}(B)$ is an isomorphism. So, like (\ref{eqn:3509}), we get 
isomorphisms 
\begin{equation} \label{eqn:3514}
\begin{aligned}
& G_S(N) = S \ot^{\mrm{L}}_{B} N = \opn{RHom}_B(T, B) \ot^{\mrm{L}}_{B} N
\\
& \quad \xar{\opn{ev}^{\mrm{R, L}}_{T, B, N}}  
\opn{RHom}_B(T, B \ot^{\mrm{L}}_{B} N)
\cong \opn{RHom}_B(T, N) = F_T(N)
\end{aligned}
\end{equation}
in $\dcat{D}(A)$. Since these isomorphisms are functorial in $N$, we see that 
$G_S \cong F_T$ as functors. Therefore $G_S$ is an equivalence, and 
$S$ is a pretilting DG $A$-$B$-bimodule.

The same proof, but now for the pretilting DG $A$-$B$-bimodule $S$ instead of 
for $T$, shows that the DG $B$-$A$-bimodule 
\begin{equation} \label{eqn:3527}
T' := \opn{RHom}_A(S, A) \in \dcat{D}(B \ot_{} A^{\mrm{op}}) 
\end{equation}
satisfies the corresponding version of (\ref{eqn:3512}), i.e.\ 
$T' \ot^{\mrm{L}}_{A} S \cong B$
in $\dcat{D}(B^{\mrm{en}})$. 
By Lemma \ref{lem:3455} the DG $B$-$A$-bimodule $T \cong T'$ is tilting.
\end{proof}

\begin{cor} \label{cor:3455}
Let $T$ be a tilting DG $B$-$A$-bimodule. Then its quasi-inverse is the DG 
$A$-$B$-bimodule $S := \opn{RHom}_B(T, B)$.
\end{cor}

\begin{proof}
This was shown in the proof of the implication ``(ii) $\Rightarrow$ (i)''
above.
\end{proof}

\begin{cor} \label{cor:3505}
Let $T$ be a tilting DG $B$-$A$-bimodule. Then $T$ is a compact generator on 
both sides, and it has the noncommutative derived Morita property on both 
sides.
\end{cor}

\begin{proof}
By Theorem \ref{thm:3455}, $T$ is a compact generator on the $B$ side, 
and it has the noncommutative derived Morita property on the $B$ side. 

Using the DG ring isomorphism 
$A^{\mrm{op}} \ot B \iso B \ot A^{\mrm{op}}$
of (\ref{eqn:3552}), we can also view $T$ as an object of 
$\dcat{D}(A^{\mrm{op}} \ot B)$. And as such, $T$ is a tilting DG 
$A^{\mrm{op}}$-$B^{\mrm{op}}$-bimodule. Now Theorem \ref{thm:3455} tells us 
that $T$ is a compact generator on the $A^{\mrm{op}}$ side, 
and it has the noncommutative derived Morita property on the 
$A^{\mrm{op}}$ side. 
\end{proof}

The example of classical Morita Theory was already given; see Example 
\ref{exa:3435}. Next is an example that goes in another direction altogether.

\begin{exa} \label{eqn:3565}
Suppose $f : A \to B$ is a quasi-isomorphism in $\catt{DGRng} \fcentover \K$.
Then 
$T := B \in \dcat{D}(B \ot A^{\mrm{op}})$
is a tilting DG $B$-$A$-bimodule. 
The equivalence 
\[ F_T := \opn{RHom}_B(T, -) : \dcat{D}(B) \to \dcat{D}(A) \]
is just the restriction functor $\opn{Rest}_f$; and the equivalence 
\[ G_T := T \ot^{\mrm{L}}_{A} (-) :
\dcat{D}(A) \to \dcat{D}(B) \]
is just the derived induction functor $\opn{LInd}_f$. 
\end{exa}

\begin{thm}[Rickard-Keller] \label{thm:3565} 
Let $A$ and $B$ be K-flat DG central $\K$-rings. Assume that there exists a 
$\K$-linear equivalence of triangulated categories 
$F : \lb \dcat{D}(A) \to \dcat{D}(B)$,
and that $\opn{H}^i(A) = 0$ for all $i \neq 0$. 
Then there exists a tilting DG $B$-$A$-bimodule $T$. 
\end{thm}

This theorem is very similar to J. Rickard's \cite[Theorem 6.4]{Ric2}
and to B. Keller's \cite[Corollary 9.2]{Kel1}. See Remark \ref{rem:4189} for a 
brief discussion.

\begin{proof}
We begin with a warning: in the proof we will construct some new DG rings,
and they might fail to be K-flat over $\K$, thus violating Convention 
\ref{conv:3445}.

Consider the object $L := F(A) \in \dcat{D}(B)$.
It is a compact generator, by Propositions \ref{prop:3545} and \ref{prop:3547}.
Let us choose a K-projective resolution $P \to L$ in 
$\dcat{C}_{\mrm{str}}(B)$, and define the DG ring
$\til{A} := \opn{End}_B(P)^{\mrm{op}}$.
So $P \in \dcat{D}(B \ot_{\K} \til{A}^{\mrm{op}})$.

Because $F$ is an equivalence, we get ring isomorphisms 
\begin{equation} \label{eqn:3570}
\opn{H}^0(A) \cong \opn{End}_{\dcat{D}(A)}(A)^{\mrm{op}} \xar{F} 
\opn{End}_{\dcat{D}(B)}(L)^{\mrm{op}} \cong 
\opn{End}_{\dcat{D}(B)}(P)^{\mrm{op}} \cong \opn{H}^0(\til{A}) .
\end{equation}
For the same reason we have 
$\opn{H}^i(\til{A}) \cong \opn{H}^i(A)  = 0$ for all $i \neq 0$. 
Define the DG rings $A_1 := \opn{smt}^{\leq 0}(A)$ and 
$\til{A}_1 := \opn{smt}^{\leq 0}(\til{A})$,
and the rings $A_2 := \opn{H}^{0}(A)$ and
$\til{A}_2 := \opn{H}^{0}(\til{A})$. 
There are canonical quasi-isomorphisms 
$A_1 \to A$, $A_1 \to A_2$, 
$\til{A}_1 \to \til{A}$ and $\til{A}_1 \to \til{A}_2$ 
in $\catt{DGRng} \centover \K$. Equation (\ref{eqn:3570}) produces an 
isomorphism $A_2 \cong \til{A}_2$ in $\catt{Rng} \centover \K$.
Now the DG ring $B$ is K-flat over $\K$, so we have induced quasi-isomorphisms
\[ B \ot A^{\mrm{op}} \lar B \ot A_1^{\mrm{op}} \to B \ot A_2^{\mrm{op}}
\cong B \ot \til{A}_2^{\mrm{op}} \lar 
B \ot \til{A}_1^{\mrm{op}} \to B \ot \til{A}^{\mrm{op}} \]
in $\catt{DGRng} \centover \K$. According to Theorem \ref{thm:2363} we get 
$\K$-linear equivalences of  triangulated categories 
\[ \begin{aligned}
& \dcat{D}(B \ot A^{\mrm{op}}) \to \dcat{D}(B \ot A_1^{\mrm{op}}) \to 
\dcat{D}(B \ot A_2^{\mrm{op}}) 
\\
& \quad \to \dcat{D}(B \ot \til{A}_2^{\mrm{op}}) \to 
\dcat{D}(B \ot \til{A}_1^{\mrm{op}}) \to \dcat{D}(B \ot \til{A}^{\mrm{op}}) .
\end{aligned} \]
There is an object $T \in \dcat{D}(B \ot A^{\mrm{op}})$ that 
corresponds to $P \in \dcat{D}(B \ot \til{A}^{\mrm{op}})$
under this chain of equivalences. All these equivalence restrict to the 
identity functor of $\dcat{D}(B)$, and hence there is an isomorphism 
\begin{equation} \label{eqn:4181}
\opn{Rest}_B(T) \cong \opn{Rest}_B(P) \cong L 
\end{equation}
in $\dcat{D}(B)$. 

Since $L$ is a compact generator on the $B$ side, the same is true for $T$.
Finally, by equation (\ref{eqn:4181}), and using the equivalence $F$,  
we know that 
\begin{equation} \label{eqn:4182}
\opn{H}^i \bigl( \opn{RHom}_B(T, T) \bigr) \cong 
\opn{H}^i \bigl( \opn{RHom}_B(L, L) \bigr) \cong
\opn{H}^i \bigl( \opn{RHom}_A(A, A) \bigr) = 0
\end{equation}
for all $i \neq 0$. By equation (\ref{eqn:3570}) we know that 
\begin{equation} \label{eqn:4183}
\opn{H}^0 \bigl( \opn{RHom}_B(T, T) \bigr) \cong 
\opn{End}_{\dcat{D}(B)}(L) \cong \opn{End}_{\dcat{D}(A)}(A) \cong 
\opn{H}^0(A) 
\end{equation}
as $\opn{H}^0(A)$-bimodules. 
Once we trace the various morphisms, formulas (\ref{eqn:4182}) and 
(\ref{eqn:4183}) imply that the derived homothety morphism 
$\opn{hm}^{\mrm{R}}_{T, A^{\mrm{op}}} : A \to \lb \opn{RHom}_B(T, T)$
in $\cat{D}(A^{\mrm{en}})$ is an isomorphism. So $T$ has the derived Morita 
property on the $B$ side. 
By Theorem \ref{thm:3455}, $T$ is a tilting DG $B$-$A$-bimodule.
\end{proof}

\begin{rem} \label{rem:4189}
J. Rickard, in \cite[Theorem 6.4]{Ric1}, considers rings $A$ and $B$, and he 
does not assume that they are flat over a base ring $\K$. So instead of 
producing a tilting complex $T$, as we do in Theorem \ref{thm:3565} above, he 
actually produces a pretilting complex. Another difference is that Rickard 
considers an equivalence 
$F : \dcat{D}^{\mrm{b}}(A) \to \dcat{D}^{\mrm{b}}(B)$,
whereas we look at the unbounded derived categories. 

B. Keller, in \cite[Corollary 9.2]{Kel1}, works in much greater generality than 
we do: instead of DG rings $A$ and $B$, he looks at DG categories $\cat{A}$ and 
$\cat{B}$. Also he does not require $\cat{A}$ to be K-flat over $\K$. 

The apparent flatness limitation in our Theorem \ref{thm:3565} can be 
effectively overcome using K-flat resolutions 
$\til{A} \to A$ and $\til{B} \to B$ in $\catt{DGRng} \centover \K$,
as explained in Remark \ref{rem:1420}. Since the restriction functors 
$\dcat{D}(A) \to \dcat{D}(\til{A})$ and 
$\dcat{D}(B) \to \dcat{D}(\til{B})$ are equivalences,
we get an equivalence 
$\til{F} : \dcat{D}(\til{A}) \to \dcat{D}(\til{B})$.
Theorem \ref{thm:3565} produces a tilting DG bimodule 
$T \in \dcat{D}(\til{B} \ot \til{A}^{\mrm{op}})$. 
\end{rem}

\begin{rem} \label{rem:3565}
In the situation of Theorem \ref{thm:3565}, it is not known whether one can 
find a tilting DG $B$-$A$-bimodule $T$ such that 
$F \cong T \ot^{\mrm{L}}_{A} (-)$ as triangulated functors. 
This question was already raised by Rickard in \cite{Ric1}, and still remains 
open, except for a few special cases (see \cite{MiYe} for hereditary rings, and 
\cite{Ch} for triangular rings). 

If we drop the condition that the cohomology of $A$ is concentrated in  
degree $0$, then there is a counterexample to the assertion of Theorem 
\ref{thm:3565}, due to B. Shipley \cite[Section 5]{Shi}, based on her joint 
work with D. Dugger \cite{DuSh}. Here is its translation to our 
terminology. She considers the DG ring  
$A = \Z[\bar{x}] := \Z \bra{x} / (x^4)$,
where $x$ is a variable of degree $-1$, and $\d(\bar{x}) := 2$. 
(Note that $A$ is weakly commutative but not strongly commutative, since 
$\bar{x}^2 \neq 0$.) 
The second DG ring is $\opn{H}(A)$. A calculation shows that 
$\opn{H}(A) \cong  \mbb{F}_2[\bar{y}] := \mbb{F}_2 \bra{y} / (y^2)$,
where $y$ is a variable of degree $-2$, and $\d(\bar{y}) := 0$.
For the base ring we take $\K := \Z$. Because $\opn{H}(A)$ is not K-flat over 
$\K$, we choose a K-flat resolution $B \to \opn{H}(A)$ over $\K$. 

Shipley asserts that on the one hand there is an equivalence of triangulated 
categories $\dcat{D}(A) \to \dcat{D}(B)$, and on the other hand there does not 
exist a DG bimodule $T \in \dcat{D}(B \ot A^{\mrm{op}})$ 
that is a compact generator on the $B$ side, 
and has the noncommutative derived Morita property on the $B$ side
(condition (iii) of Theorem \ref{thm:3455}). Thus there does not exist a 
tilting DG $B$-$A$-bimodule. 
\end{rem}

\begin{rem} \label{rem:3596}
In algebraic geometry the role of tensoring with tilting bimodule complexes is 
played by {\em Fourier-Mukai transforms}. This theory is at the core of 
contemporary birational geometry. See the survey \cite{HiVdB} or the book 
\cite{Huy}. 
\end{rem}





\mysubsection{Tilting Bimodule Complexes over Rings} 
\label{subsec:tilting-rings}

We continue with Convention \ref{conv:3445}, but in this subsection 
we only look at rings. So our rings are flat central over the base ring $\K$, 
$\ot = \ot_{\K}$, and $A^{\mrm{en}} = A \ot A^{\mrm{op}}$ for a ring $A$. 
All bimodules are $\K$-central, and all additive functors are $\K$-linear.  
Rather than speaking about DG bimodules, here we speak about complexes of 
bimodules. By default our rings are nonzero. 

In classical Morita theory, an 
{\em invertible  $B$-$A$-bimodule}%
\index{Invertible! bimodule}
is a bimodule $P$ for which there exists an $A$-$B$-bimodule $Q$, such that  
$Q \ot_B P \cong A$ in $\dcat{M}(A^{\mrm{en}})$ 
and
$P \ot_A Q \cong B$ in $\dcat{M}(B^{\mrm{en}})$.
It is known that $P$ is a finitely generated projective module over $B$ and 
over $A^{\mrm{op}}$; and the reverse is true for $Q$. Furthermore, every 
$\K$-linear equivalence 
$F : \dcat{M}(A) \to \dcat{M}(B)$
is isomorphic, as a functor, to the functor $P \ot_A (-)$ for an invertible 
bimodule $P$, and this $P$ is unique up to a unique isomorphism. See 
\cite[Section 4.1]{Row} for proofs. 

We see that invertible bimodules are a very special kind of tilting bimodule 
complexes. Proposition \ref{prop:3575} clarifies matters. 
But first a lemma, which is a noncommutative variant of Lemma \ref{lem:2177}

\begin{lem}[K\"unneth Trick] \label{lem:3595} 
Let $M \in \dcat{D}^-(B \ot A^{\mrm{op}})$ and 
$N \in \dcat{D}^-(A \ot C^{\mrm{op}})$,  
and let $i_1, j_1 \in \Z$ be such that 
$\opn{sup}(\opn{H}(M)) \leq i_1$ and $\opn{sup}(\opn{H}(N)) \leq j_1$. 
Then 
\[ \opn{H}^{i_1 + j_1}(M \ot^{\mrm{L}}_A N) \cong 
\opn{H}^{i_1}(M) \ot_A \opn{H}^{j_1}(N)  \]
in $\dcat{M}(B \ot C^{\mrm{op}})$. 
\end{lem}

\begin{exer} \label{exer:3595}
Prove Lemma \ref{lem:3595}.
\end{exer}

\begin{prop} \label{prop:3575}
The following conditions are equivalent for a tilting bimodule complex 
$T \in \dcat{D}(B \ot A^{\mrm{op}})$, with quasi-inverse $S = T^{\vee}$. 
\begin{enumerate}
\rmitem{i} There is an isomorphism $T \cong P$ in 
$\dcat{D}(B \ot A^{\mrm{op}})$
for some invertible $B$-$A$-bimodule $P$. 

\rmitem{ii} $\opn{H}^0(T)$ is a projective $B$-module, and 
$\opn{H}^i(T) = 0$ for all $i \neq 0$. 

\rmitem{iii} $\opn{H}^i(T) = 0$ and $\opn{H}^i(S) = 0$ for all $i \neq 0$.  
\end{enumerate} 
\end{prop}

\begin{proof} \mbox{} 

\smallskip \noindent 
(i) $\Rightarrow$ (ii): this is trivial.

\medskip \noindent
(ii) $\Rightarrow$ (iii): 
Let $P := \opn{H}^0(T) \in \dcat{M}(B \ot A^{\mrm{op}})$,
so there is an isomorphism 
$T \cong P$ in $\dcat{D}(B \ot A^{\mrm{op}})$.
According to Corollary \ref{cor:3455} we have 
$S \cong \opn{RHom}_B(T, B) \cong \opn{Hom}_B(P, B)$.

\medskip \noindent
(iii) $\Rightarrow$ (i): 
Define $P := \opn{H}^{0}(T) \in \dcat{M}(B \ot A^{\mrm{op}})$
and $Q := \opn{H}^{0}(S) \in \dcat{M}(A \ot B^{\mrm{op}})$.
The K\"unneth Trick (Lemma \ref{lem:3595}) says that 
$P \ot_B Q \cong \opn{H}^{0}(T \ot^{\mrm{L}}_{B} S) \cong \opn{H}^{0}(A) = A$
and
$Q \ot_A P \cong \opn{H}^{0}(S \ot^{\mrm{L}}_{A} T) \cong \opn{H}^{0}(B) = B$
in $\dcat{M}(A^{\mrm{en}})$ and $\dcat{M}(B^{\mrm{en}})$ respectively. 
Thus $P$ and $Q$ are invertible bimodules. 
But $T \cong P$ in $\dcat{D}(B \ot A^{\mrm{op}})$. 
\end{proof}

The {\em Jacobson radical} of a ring $A$ is the intersection of all 
maximal left (or right) ideals of $A$. See \cite[Section 4.4.2]{Jac} or 
\cite[Section 2.5]{Row}. 

\begin{dfn}
A ring $A$, with Jacobson radical $\r$, is called {\em local} 
\index{Ring! local NC}
if $A / \r$ is a simple artinian ring. 
\end{dfn}

Note that this definition is wider than what is found in some books (e.g.\ 
\cite{Row}), where the condition is that $A / \r$ is a division ring. 

\begin{exa} \label{exa:3575}
If $C$ is a commutative local ring with maximal idea $\m$, and $n \geq 1$, 
then the matrix ring $A := \opn{Mat}_{n \times n}(C)$ is local, with Jacobson 
radical $\m \cd A = \opn{Mat}_{n \times n}(\m)$.  
\end{exa}

\begin{lem} \label{lem:3575}
Let $A$ be a local ring, let $M$ be a nonzero finitely generated right 
$A$-module, and let $N$ be a nonzero finitely generated left $A$-module. Then 
$M \ot_A N$ is nonzero. 
\end{lem}

\begin{proof}
Define $K := A / \r$, which is a simple artinian ring. By the noncommutative 
Nakayama Lemma (see \cite[Proposition 2.5.24]{Row}) the right $K$-module 
$M \ot_A K$ is nonzero, and also the left $K$-module $K \ot_A N$ is nonzero.
There is a canonical surjection 
\[ M \ot_A N \to (M \ot_A K) \ot_K (K \ot_A N) . \]
It suffices to prove that the target is nonzero. Thus we can assume that 
$A = K$ is a simple artinian ring. 

Every left module over $K$ is a  direct sum of simple ones; so we can assume 
that $N$ is simple. Likewise we can assume that $M$ is a simple right 
$K$-module. 

There is a ring isomorphism $K \cong \opn{Mat}_{n \times n}(D)$, 
where $D$ is a division ring and $n$ is a positive integer. The simple left 
$K$-module $N$ is isomorphic to $D^n$, seen as a column module; thus it is in 
fact a $K$-$D$-bimodule. Similarly $M \cong D^n$, seen as a row module; thus it 
is a $D$-$K$-bimodule. By an easy calculation (this is an elementary case of 
Morita equivalence) there is an isomorphism 
$M \ot_K N \cong D$ of $D$-$D$-bimodules. This is nonzero. 
\end{proof}

\begin{thm}[\cite{RoZi}, \cite{Ye4}] \label{thm:3575}
Let $A$ and $B$ be flat central $\K$-rings, with $A$ local, and let 
$T$ be a tilting $B$-$A$-bimodule complex%
\index{Tilting! complex of bimodules}%
\index{Invertible! bimodule}.
Then there is an iso\-morph\-ism $T \cong P[n]$
in $\dcat{D}(B \ot A^{\mrm{op}})$ for some 
invertible $B$-$A$-bimodule
$P$ and some integer $n$. 
\end{thm}

See Remark \ref{rem:3430} regarding the history of this theorem. 

\begin{proof}
Let $S \in \dcat{D}(A \ot B^{\mrm{op}})$
be the quasi-inverse of $T$. From Corollary \ref{cor:3505} and Theorem 
\ref{thm:3400} we know that $T$ and $S$ are algebraically perfect complexes on 
both sides. By Lemma \ref{lem:3415} we know that 
$\opn{con}(\opn{H}(T)) = [i_0, i_1]$ and 
$\opn{con}(\opn{H}(S)) = [j_0, j_1]$ for some integers $i_0 \leq i_1$ and 
$j_0 \leq j_1$; and also that $P := \opn{H}^{i_1}(T)$ and 
$Q := \opn{H}^{j_1}(S)$ are finitely presented modules on both 
sides. 

We now use the K\"unneth trick to obtain an isomorphism
\[ P \ot_A Q \cong 
\opn{H}^{i_1 + j_1}(T \ot^{\mrm{L}}_{A} S) \cong \opn{H}^{i_1 + j_1}(B) \]
in $\cat{M}(B^{\mrm{en}})$. By Lemma \ref{lem:3575} we know that 
$P \ot_A Q$ is nonzero. Therefore $i_1 + j_1 = 0$, and 
$P \ot_A Q \cong B$ in $\dcat{M}(B^{\mrm{en}})$. 
For similar reasons, on the reverse side we have 
$Q \ot_B P \cong \opn{H}^{0}(S \ot^{\mrm{L}}_{B} T) \cong A$
in $\dcat{M}(A^{\mrm{en}})$. Thus $P$ and $Q$ are invertible bimodules. 

We now restrict $T$ to $\dcat{D}(A^{\mrm{op}})$. Because its top cohomology 
module $\opn{H}^{i_1}(T) = P$ is projective, we can split it off (see Lemma 
\ref{lem:2189}, that holds also for a noncommutative ring), to get an 
isomorphism $T \cong T' \oplus P[-i_1]$ in $\dcat{D}(A^{\mrm{op}})$, 
with $\opn{sup}(\opn{H}(T')) \leq i_1 - 1$.
Similarly there's an isomorphism 
$S \cong S' \oplus Q[-j_1]$ in $\dcat{D}(A)$
with $\opn{sup}(\opn{H}(S')) \leq j_1 - 1$.
It follows that 
\[ \begin{aligned}
& B \cong T \ot^{\mrm{L}}_{A} S \cong 
\bigl( T' \oplus P[-i_1] \bigr) \ot^{\mrm{L}}_{A} 
\bigl( S' \oplus Q[-j_1] \bigr) 
\\
& \quad 
\cong (T' \ot^{\mrm{L}}_{A} S') \oplus
\bigl( P[-i_1] \ot_A S' \bigr) \oplus \bigl( T' \ot_A Q[-j_1] \bigr) 
\oplus B 
\end{aligned} \]
in $\dcat{D}(\K)$. The object $B \in \dcat{D}(\K)$ has cohomology concentrated 
in degree $0$. The direct summand $T' \ot_A Q[-j_1]$ has cohomology 
concentrated in degrees $\leq -1$, so its cohomology must be zero, and 
therefore $T' \ot_A Q[-j_1] = 0$ in $\dcat{D}(\K)$. 
But $Q$ is an invertible bimodule, and this forces $T' = 0$
in $\dcat{D}(A^{\mrm{op}})$. The conclusion is that $\opn{H}^{i}(T) = 0$ for 
all $i \neq i_1$.

Returning to the category $\dcat{D}(A \ot B^{\mrm{op}})$
we see that $T \cong P[-i_1]$. Finally we take $n := -i_1$. 
\end{proof}

\begin{cor} \label{cor:3580}
Let $A$ and $B$ be rings, with $A$ local.
If there is an equivalence of triangulated categories
$\dcat{D}(A) \to \dcat{D}(B)$, then there is an equivalence of 
abelian categories $\dcat{M}(A) \to \dcat{M}(B)$. In other words, $A$ and $B$ 
are Morita equivalent. 
\end{cor}

\begin{proof}
According to Theorem \ref{thm:3565} there exists a tilting $B$-$A$-bimodule 
complex $T$. Theorem \ref{thm:3575} says that $T \cong P[n]$ for an invertible 
$B$-$A$-bimodule $P$. Then 
$P \ot_A (-) : \dcat{M}(A) \to \dcat{M}(B)$
is an equivalence of abelian categories. 
\end{proof}

\begin{dfn} \label{dfn:3590}
For a central $\K$-ring $A$ we define the 
{\em noncommutative Picard group of $A$ relative to $\K$}
\index{Picard group! noncommutative}
to be the group $\opn{Pic}_{\K}(A)$, whose elements are 
the isomorphism classes in $\dcat{M}(A^{\mrm{en}})$ of the invertible 
bimodules. The operation is induced by $(- \ot_A -)$, and the unit element is 
the class of $A$. 
\end{dfn}

The derived Picard group was introduced in Definition \ref{dfn:3457}.

\begin{cor} \label{cor:3590}
If $A$ is a local ring, then 
\[ \opn{DPic}_{\K}(A) = \opn{Pic}_{\K}(A) \times \Z . \]
\end{cor}

\begin{proof}
By Theorem \ref{thm:3575} there is a surjective group homomorphism 
$\opn{Pic}_{\K}(A) \times \Z \to  \opn{DPic}_{\K}(A)$,
defined by $(P, n) \mapsto P[n]$.
It is injective because the functor 
$\dcat{M}(A^{\mrm{en}}) \to \dcat{D}(A^{\mrm{en}})$
is fully faithful. 
\end{proof}

The center of a ring $A$ is denoted by $\opn{Cent}(A)$. Of course 
$\opn{Cent}(A) = \opn{Cent}(A^{\mrm{op}})$. 
Given a complex $M \in  \dcat{D}(B \ot A^{\mrm{op}})$, 
there are ring homomorphisms 
\begin{equation} \label{eqn:3580}
\opn{chm}^{\dcat{D}}_{M, A^{\mrm{op}}} : \opn{Cent}(A) \to 
\opn{End}_{\dcat{D}(B \ot A^{\mrm{op}})}(M)
\end{equation}
and 
\begin{equation} \label{eqn:3581}
\opn{chm}^{\dcat{D}}_{M, B} : \opn{Cent}(B) \to 
\opn{End}_{\dcat{D}(B \ot A^{\mrm{op}})}(M) 
\end{equation}
that we call the {\em central homotheties} through $A^{\mrm{op}}$ and $B$
respectively. The formulas are the obvious ones: for an element 
$a \in \opn{Cent}(A)$ the action on $M$ is 
$\opn{chm}^{\dcat{D}}_{M, A^{\mrm{op}}}(a)(m) := m \cd a$ 
for $m \in M^i$. Likewise for $\opn{chm}^{\dcat{D}}_{M, B}$. 

\begin{lem} \label{lem:3580}
If $T$ is a tilting $B$-$A$-bimodule complex, then the ring homomorphisms 
$\opn{chm}^{\dcat{D}}_{M, A^{\mrm{op}}}$ and 
$\opn{chm}^{\dcat{D}}_{M, B}$ are both isomorphisms. 
\end{lem}

\begin{proof}
Let $S$ be the quasi-inverse of $T$. The functor 
\[ G := (-) \ot^{\mrm{L}}_{A} S : \dcat{D}(B \ot A^{\mrm{op}}) \to 
\dcat{D}(B \ot B^{\mrm{op}}) = \dcat{D}(B^{\mrm{en}}) \]
is an equivalence. It induces a ring isomorphism
\[ G : \opn{End}_{\dcat{D}(B \ot A^{\mrm{op}})}(T) \iso 
\opn{End}_{\dcat{D}(B^{\mrm{en}})}(B) , \]
and 
$G \circ \opn{chm}^{\dcat{D}}_{T, B} = \opn{chm}^{\dcat{D}}_{B, B}$.
But 
$\opn{End}_{\dcat{D}(B^{\mrm{en}})}(B) \cong 
\opn{End}_{\dcat{M}(B^{\mrm{en}})}(B)$,
and the ring homomorphism 
$\opn{Cent}(B) \to \opn{End}_{\dcat{M}(B^{\mrm{en}})}(B)$
is bijective. 

The proof for $\opn{chm}^{\dcat{D}}_{T, A^{\mrm{op}}}$ is similar.
\end{proof}

\begin{dfn} \label{dfn:5090}
Given a tilting $B$-$A$-bimodule complex $T$, we denote by 
$g_T : \opn{Cent}(A) \to \opn{Cent}(B)$
the ring isomorphism such that 
$\opn{chm}^{\dcat{D}}_{T, A^{\mrm{op}}}  = 
\opn{chm}^{\dcat{D}}_{T, B} \circ \, g_T$.
\end{dfn}

Even though the localization of a noncommutative ring $B$ with respect to a 
multiplicatively closed subset $Z \sub B$ is problematic in general, there is 
no difficulty at all if $Z \sub \opn{Cent}(B)$. In this case,
letting $C := \opn{Cent}(B)$, we get canonical $A$-ring isomorphisms
$B_{Z} \cong C_{Z} \ot_C B \cong B \ot_C C_Z$.

\begin{lem} \label{lem:3581}
Let $T$ be a tilting $B$-$A$-bimodule complex, and assume that the ring $A$ is 
commutative. Let $Z \sub A$ be a multiplicatively 
closed set, and define $A' := A_Z$ and $B' := B_{g_T(Z)}$. Then 
the complex 
$T' := B' \ot_B T \ot_A A'$ in 
$\dcat{D}(B' \ot A'^{\, \mrm{op}})$
is a tilting $B'$-$A'$-bimodule complex.
\end{lem}

\begin{proof}
The cohomology $\opn{H}(T)$ is a central graded $A$-bimodule, where the 
left action of $A$ on $\opn{H}(T)$ is via the isomorphism  $g_T$. (Warning: the 
complex of 
bimodules $T$ need not be central over $A$.) The flatness of $A \to A'$ and 
$B \to B'$ gives 
\[ \opn{H}(T') = \opn{H}(B' \ot_B T \ot_A A') \cong 
B' \ot_B \opn{H}(T) \ot_A A' \cong A' \ot_A \opn{H}(T) \ot_A A' . \]
Therefore 
$\opn{H}(T) \ot_A A' \to \opn{H}(T')$ 
is an isomorphism of graded modules, and this implies that 
\begin{equation} \label{eqn:3585}
T \ot_A A' \to T'
\end{equation}
is an isomorphism in 
$\dcat{D}(B \ot A'^{\, \mrm{op}})$. 

Let 
$S := \opn{RHom}_B(T, B)$ in $\dcat{D}(A \ot B^{\mrm{op}})$
be the quasi-inverse of $T$. 
In analogy to Definition \ref{dfn:5090}, there is a ring isomorphism 
$h_S : A \to \opn{Cent}(B)$
such that 
$\opn{chm}^{\dcat{D}}_{S, A}  = 
\opn{chm}^{\dcat{D}}_{S, B^{\mrm{op}}} \circ \, h_S$.
The cohomology $\opn{H}(S)$ is a central graded $A$-bimodule, 
where the right action of $A$ on $\opn{H}(S)$ is via $h_S$.

Let us consider the derived tensor product 
$S \ot^{\mrm{L}}_{B} T \in \dcat{D}(A^{\mrm{en}})$,
and the four ways that an element $a \in A$ acts on it as an 
object of $\dcat{D}(A^{\mrm{en}})$. 
The action of $a$ from the right side of $T$ the same as the action of 
$g_T(a) \in \opn{Cent}(B)$ from the left side of $T$. This is also the action 
of $g_T(a)$ from the right side of $S$, and also the action of 
$h_S^{-1}(g_T(a))$ from the left on $S$. 
But there are ring isomorphisms 
\[ \opn{End}_{\dcat{D}(A^{\mrm{en}})}(S \ot^{\mrm{L}}_{B} T) \cong 
\opn{End}_{\dcat{D}(A^{\mrm{en}})}(A) \cong A , \]
and $a$ acts on $A$ the same from both sides. The conclusion is that 
$h_S^{-1}(g_T(a)) = a$. 
Because $a$ is an arbitrary element, we see that $g_T = h_S$ as isomorphisms 
$A \iso \opn{Cent}(B)$. So we have a single ring isomorphism 
$A \cong \opn{Cent}(B)$, for which both $\opn{H}(T)$ and $\opn{H}(S)$ are 
central graded $A$-bimodules. 

Define 
$S' := A' \ot_A S \ot_B B' \in \dcat{D}(A' \ot B'^{\, \mrm{op}})$. 
The same arguments as for $T$ show that 
$A' \ot_A S \to S'$ is an isomorphism in 
$\dcat{D}(A' \ot B^{\mrm{op}})$.

We now calculate:
\[ \begin{aligned}
& S' \ot^{\mrm{L}}_{B'} T' = (A' \ot_A S \ot_B B') \ot^{\mrm{L}}_{B'}
(B' \ot_B T \ot_A A') 
\\
& \quad \cong (A' \ot_A S) \ot^{\mrm{L}}_{B} (B' \ot_B T \ot_A A') 
\cong^{\dag} (A' \ot_A S) \ot^{\mrm{L}}_{B} (T \ot_A A') 
\\
& \quad \cong A' \ot_A (S \ot^{\mrm{L}}_{B} T) \ot_A A') 
\cong A' \ot_A A \ot_A A' \cong A'
\end{aligned} \]
in $\dcat{D}(A'^{\, \mrm{en}})$. We used the associativity of 
$(- \ot^{\mrm{L}}_{(-)} -)$ several times, and also the ring isomorphism 
$A' \ot_A A' \cong A'$. The isomorphism $\cong^{\dag}$ is by formula 
(\ref{eqn:3585}). Similarly we show that 
$T' \ot^{\mrm{L}}_{A'} S' \cong B'$
in $\dcat{D}(B'^{\, \mrm{en}})$. 
\end{proof}

\begin{thm}[\cite{RoZi}, \cite{Ye4}] \label{thm:3585}
Let $A$ and $B$ be flat central $\K$-rings, and assume $A$ is commutative with 
connected spectrum. Let $T$ be a tilting 
$B$-$A$-bimodule complex% 
\index{Differential graded bimodule! tilting}.
Then there is an isomorphism $T \cong P[n]$  
in $\dcat{D}(B \ot A^{\mrm{op}})$, for some 
invertible $B$-$A$-bimodule $P$ and some integer $n$. 
\end{thm}

See Remark \ref{rem:3430} regarding the history of this theorem. 

\begin{proof}
Because $A$ is commutative we have $A = A^{\mrm{op}}$.
The ring homomorphism $g_T$ makes $B$ into a central $A$-ring. 
We may assumed that $A \neq 0$, and hence $T \neq 0$.
The complex $T$ is algebraically perfect over $A$, 
so it has bounded cohomology, say with 
$\opn{sup}(\opn{H}(T)) = i_1 \in \Z$.
By Lemma \ref{lem:3415} the $A$-module $P := \opn{H}^{i_1}(T)$ is finitely 
presented.

For a prime $\p \in \opn{Spec}(A)$, with corresponding local ring
$A_{\p}$, we write $P_{\p} := P \otimes_A A_{\p}$.
Define $Y \sub \opn{Spec}(A)$ to be the support of $P$, i.e.\ 
\[ Y := \{ \p \in \opn{Spec}(A) \mid P_{\p} \neq 0 \} . \]
This is a nonempty set. 
Since $P$ is finitely generated it follows that $Y$ is a closed subset of 
$\opn{Spec}(A)$; see \cite[Proposition II.4.17]{Bou}.

Take any prime $\p \in Y$, and let $B_{\p} :=  B \otimes_A A_{\p}$.
Then, by Lemma \ref{lem:3581}, the complex
\[ T_{\p} := B_{\p} \otimes_B T \otimes_A A_{\p} \in 
\dcat{D}(B_{\p} \otimes A_{\p}^{\mrm{op}}) \]
is a tilting $B_{\p}$-$A_{\p}$-bimodule complex. Since
$\mrm{H}^{i_1}(T_{\p}) \cong P_{\p} \neq 0$,
Theorem \ref{thm:3575} implies that 
\begin{equation} \label{eqn:3587}
T_{\p} \cong P_{\p}[-i_1] \in \dcat{D}(B_{\p} \otimes A_{\p}^{\mrm{op}}) ,
\end{equation}
and that $P_{\p}$ is an invertible $B_{\p}$-$A_{\p}$-bimodule. 
In particular, $P_{\p}$ is a nonzero finitely generated projective 
$A_{\p}$-module. Thus $P_{\p}$ is a free $A_{\p}$-module, of rank 
$r_{\p} > 0$. Recall that $P$ is a finitely presented $A$-module. 
According to \cite[Section II.5.1,  Corollary]{Bou} there is an 
open neighborhood $U$ of $\p$ in $\opn{Spec}(A)$ on which $P$ is free of rank 
$r_{\p}$. In particular $P_{\q} \neq 0$ for all $\q \in U$.
Therefore $U \sub Y$. 

The conclusion is that $Y$ is also open in $\opn{Spec}(A)$.
Since $\opn{Spec}(A)$ is connected, it follows that 
$Y = \opn{Spec}(A)$. Another conclusion is that $P$ is projective as an 
$A$-module -- see \cite[Section II.5.2, Theorem 1]{Bou}.

Going back to equation (\ref{eqn:3587}) we see that 
$\mrm{H}^i(T)_{\p} \cong \mrm{H}^i(T_{\p}) = 0$
for all $i \neq i_1$. Therefore 
$\mrm{H}^i(T) = 0$ for $i \neq i_1$. By truncation we get an 
isomorphism $T \cong P[n]$ in 
$\dcat{D}(B \otimes A^{\mrm{op}})$,
where $n := -i_1$.
Finally, by Proposition \ref{prop:3575} the $B$-$A$-bimodule $P$ is invertible.
\end{proof}

Let $A$ be a commutative ring. An $A$-module $P$ can 
be viewed as a central $A$-bimodule. If $P$ is a rank $1$ projective 
$A$-module, then as a bimodule it is invertible. The usual Picard group of $A$
(see \cite[Section II.6]{Har})
is then the subgroup $\opn{Pic}_{A}(A)$ of $\opn{Pic}_{\K}(A)$, whose elements 
are the isomorphism classes of the central $A$-bimodules; and we refer to it 
here as the {\em commutative Picard group of $A$}. 

\begin{exer} \label{exer:3592}
Let $A$ be a commutative ring. We denote by $\opn{Aut}_{\K}(A)$ the group 
of ring automorphisms of $A$ (i.e.\ the Galois group). Show that 
\[ \opn{Pic}_{\K}(A) \cong \opn{Aut}_{\K}(A) \ltimes \opn{Pic}_{A}(A) . \]
(Hints: (1) Every invertible $A$-bimodule $P$ is isomorphic to 
$P' \ot_A P''$, where $P'$ is free of rank $1$ as a left $A$-module, and 
$P''$ is a central invertible bimodule.
This is proved like Theorem \ref{thm:3585}. 
(2) If $P'$ is an invertible $A$-bimodule that's free of rank $1$ as a left 
$A$-module, then it is also free as a right $A$-module, with the same basis 
element $e \in P'$. The proof is like those of Lemma \ref{lem:3575} and Theorem 
\ref{thm:3585}. Conclude that $P' \cong A(\ga)$, where 
$\ga \in \opn{Aut}_{\K}(A)$, and $A(\ga)$ is the twisted bimodule that is the
left $A$-module $A$, with right action 
$p \cd a := \ga(a) \cd p$ for $p \in A(\ga)$ and $a \in A$.)
\end{exer}

\begin{cor} \label{cor:3591} 
\index{1-DPic(A)@$\opn{DPic}_{\K}(A)$}
Let $A$ be a commutative ring with connected spectrum. Then 
\[ \opn{DPic}_{\K}(A) = \opn{Pic}_{\K}(A) \times \Z . \]
\end{cor}

\begin{proof}
The same as that of Corollary \ref{cor:3590}. 
\end{proof}

\begin{cor} \label{cor:3597} 
Let $A$ and $B$ be rings, and assume that $A$ is commutative ring with 
connected spectrum. 
If there is an equivalence of triangulated categories
$\dcat{D}(A) \to \dcat{D}(B)$, then there is an equivalence of 
abelian categories $\dcat{M}(A) \to \dcat{M}(B)$. In other words, $A$ and $B$ 
are Morita equivalent.
\end{cor}

\begin{proof}
According to Theorem \ref{thm:3565} there exists a tilting $B$-$A$-bimodule 
complex $T$. Theorem \ref{thm:3585} says that $T \cong P[n]$ for an 
invertible $B$-$A$-bimodule $P$. Then 
$P \ot_A (-) : \dcat{M}(A) \to \dcat{M}(B)$ 
is $\K$-linear equivalence of abelian categories. 
Moreover, $P$ is a projective $A$-module of rank $r$ for some positive integer 
$r$, and 
$B \cong \opn{End}_A(P)^{\mrm{op}}$. 
\end{proof}

\begin{rem} \label{rem:3610}
In the situation of Corollary \ref{cor:3597}, we can view $B$ as a central 
$A$-ring, using the homomorphism $g_T$ from Definition \ref{dfn:5090}. Then $B$ 
is an Azumaya $A$-ring; see \cite[Section 5.3]{Row}. Moreover, letting 
$X := \opn{Spec}(A)$, and letting $\BB$ be the sheafification of $B$ to $X$, 
then $\BB$ is a trivial Azumaya $\OO_X$-ring, in the sense of 
\cite[Section IV.2]{Mil}.
\end{rem}

Let $A$ be a commutative ring. If $\opn{Spec}(A)$ has a finite connected 
component decomposition, in the sense of Definition \ref{dfn:3210}(1),
then the connected component decomposition of $A$ is 
$A = \prod_{i = 1}^m \, A_i$, see Definition \ref{dfn:3210}(2).

\begin{cor} \label{cor:3585}
Let $A$ and $B$ be rings. Assume $A$ is commutative, and it 
has a finite connected component decomposition 
$A = \prod_{i = 1}^m \, A_i$.
Let $T$ be a tilting $B$-$A$-bimodule complex. We consider $B$ as a central 
$A$-ring via the isomorphism $g_T$ from Definition \tup{\ref{dfn:5090}}. Then 
$T \cong \bigoplus\nolimits_{i = 1}^m P_i[n_i]$
in $\dcat{D}(B \ot A^{\mrm{op}})$,
where $B_i := A_i \ot_A B$, $P_i$ is an invertible $B_i$-$A_i$-bimodule, and 
$n_i \in \Z$. 
\end{cor}

\begin{cor} \label{cor:3595} 
Let $A$ be a commutative ring, and assume it has 
a finite connected component decomposition $A = \prod_{i = 1}^m \, A_i$.
Then 
\[ \opn{DPic}_{\K}(A) = \opn{Pic}_{\K}(A) \times \Z^m . \] 
\end{cor}

\begin{cor} \label{cor:3598} 
Let $A$ and $B$ be rings. Assume that $A$ is commutative and it has a  
finite connected component decomposition. 
If there is an equivalence of triangulated categories
$\dcat{D}(A) \to \dcat{D}(B)$, then there is an equivalence of 
abelian categories $\dcat{M}(A) \to \dcat{M}(B)$. 
\end{cor}

\begin{exer} \label{exer:3585}
Prove Corollaries \ref{cor:3585}, \ref{cor:3595} and \ref{cor:3598}. Cf.\ 
Exercise \ref{exer:2176}. 
\end{exer}

The general case of these last three corollaries, namely when $A$ 
does not have a finite connected component decomposition, is dealt with in 
\cite[Theorem 6.13]{Ye8}. 

\begin{rem} \label{rem:4800}
The results in this subsection have variants in which the flatness condition
on the $\K$-rings $A$ and $B$ is dropped. 

For instance, Theorem \ref{thm:3575} can be modified as follows: If $A$ and $B$ 
are rings, with $A$ local, and if $T$ is a pretilting $B$-$A$-bimodule complex, 
then there is an isomorphism 
$T \cong P[n]$ in $\dcat{D}(A \ot B^{\mrm{op}})$ 
for some invertible $B$-$A$-bimodule $P$ and some integer $n$. 
Here is an outline of the proof. One chooses nonpositive K-flat resolutions
$\til{A} \to A$ and $\til{B} \to B$
in $\catt{DGRng} \centover \K$. 
Next a version of the K\"unneth trick (Lemma \ref{lem:3595}) for nonpositive DG 
rings has to be proved: For 
$M \in \dcat{D}^-(\til{B} \ot \til{A}^{\mrm{op}})$
and 
$N \in \dcat{D}^-(\til{A} \ot \til{B}^{\mrm{op}})$
satisfying 
$\opn{sup}(\opn{H}(M)) \leq i_1$ and $\opn{sup}(\opn{H}(N)) \leq j_1$,
there is an isomorphism
\[ \opn{H}^{i_1 + j_1}(M \ot^{\mrm{L}}_{\til{A}} N) \cong 
\opn{H}^{i_1}(M) \ot_A \opn{H}^{j_1}(N)  \]
in $\dcat{M}(B^{\mrm{en}})$. 
Then the proof of Theorem \ref{thm:3575} is 
repeated, for the tilting DG bimodule 
$T \in \dcat{D}(\til{A} \ot \til{B}^{\mrm{op}})$. 

The proofs of the other results are along the same lines. 
\end{rem}

\begin{rem} \label{rem:3590}
Assume $A$ is a commutative DG ring (Definitions \ref{dfn:3091} and 
\ref{dfn:3090}), with reduction $\bar{A} = \opn{H}^0(A)$. Here we do not need 
to assume flatness over the base ring $\K$. A DG $A$-module can 
be viewed as a central DG $A$-bimodule, so we have a derived tensor product 
\[ (- \ot^{\mrm{L}}_{A} -) : \dcat{D}(A) \times \dcat{D}(A) \to \dcat{D}(A) . 
\]
This operation is symmetric, i.e.\ 
$S \ot^{\mrm{L}}_{A} T \cong T \ot^{\mrm{L}}_{A} S$. 

A DG $A$-module $T$ is called a {\em
tilting DG $A$-module} if there is some $S \in \dcat{D}(A)$ such that 
$S \ot^{\mrm{L}}_{A} T \cong A$.
The {\em commutative derived Picard group} of $A$ is the group 
$\opn{DPic}_{A}(A)$, whose elements are the isomorphism classes in $\dcat{D}(A)$
of the tilting DG modules, the multiplication is induced by 
$(- \ot^{\mrm{L}}_{A} -)$, and the unit is the class of $A$. It is an abelian 
group. 

In \cite[Theorem 6.14]{Ye8} it was proved that the group homomorphism 
$\opn{DPic}_{A}(A) \to \opn{DPic}_{\bar{A}}(\bar{A})$,
which induced by the derived reduction functor 
$\bar{A} \ot^{\mrm{L}}_{A} (-)$,
is bijective. If the commutative ring $\bar{A}$
has a finite connected component decomposition 
$A = \prod_{i = 1}^m A_i$, then there is an isomorphism
\[ \opn{DPic}_{\bar{A}}(\bar{A}) \cong 
\opn{Pic}_{\bar{A}}(\bar{A}) \times \Z^m  . \]
This is according to Corollary \ref{cor:3595} with 
$\K = \bar{A}$. 
\end{rem}

As Corollaries \ref{cor:3590} and \ref{cor:3591} show, if the ring $A$ is 
either local, or commutative with connected spectrum, then the group 
$\opn{DPic}_{\K}(A)$ is not very interesting: it is 
\begin{equation} \label{4180}
\opn{DPic}_{\K}(A) = \opn{Pic}_{\K}(A) \times \bra{\si} , 
\end{equation}
where $\opn{Pic}_{\K}(A)$ is the classical contribution, and 
$\bra{\si} \cong \Z$ is the subgroup generated by the element $\si$, which is 
the class of the tilting complex $A[1]$.

The next example shows that matters are very different when $A$ is neither 
commutative nor local.

\begin{exa} \label{exa:4180}
Let $\K$ be an algebraically closed field and let $n$ be an integer $\geq 2$. 
Consider the ring $A$ of upper-triangular $n \times n$ matrices. For $n = 2$ it 
is $A = \sbmat{ \K & \K \\[0.1em] 0 & \K}$.
The group $\opn{DPic}_{\K}(A)$ contains a new element in this case. 
The classical NC Picard group $\opn{Pic}_{\K}(A)$ is trivial here. But the 
bimodule $A^* := \opn{Hom}_{\K}(A, \K)$
is tilting, and its class $\mu \in \opn{DPic}_{\K}(A)$ satisfies 
\begin{equation} \label{eqn:4185}
\mu^{n + 1} = \si^{n - 1} . 
\end{equation}
This was calculated by A. Yekutieli (for $n = 2$) and E. Kreines (for 
$n \geq 3$) in \cite{Ye4}. When $n = 2$ this says that 
$\si = \mu^3$, so the group $\opn{DPic}_{\K}(A)$ is larger than the classical 
part 
$\opn{Pic}_{\K}(A) \times \bra{\si} = \bra{\si} \cong \Z$.

Later, in \cite{MiYe}, J.-I. Miyachi and Yekutieli showed that the group 
$\opn{DPic}_{\K}(A)$ is abelian, it is generated by $\si$ and $\mu$, and the 
only relation is (\ref{eqn:4185}). Note that the ring $A$ is isomorphic to the 
path ring $\K[Q]$ of the Dynkin quiver $Q$ of type $\mrm{A}_n$ with all
arrows going in the same direction. 
The paper \cite{MiYe} contains calculations 
of the groups $\opn{DPic}_{\K}(A)$ for several other types of path rings of 
quivers. 

M. Kontsevich interprets the relation (\ref{eqn:4185}) as follows: 
the {\em fractional Calabi-Yau dimension} of the ``noncommutative smooth proper 
space'' $\dcat{D}^{\mrm{b}}_{\mrm{f}}(A)$ is $\frac{n - 1}{n + 1}$. 
More on this in Remark \ref{rem:4505} and Example \ref{exa:4870}. 
\end{exa}

\begin{rem} \label{rem:3430}
Here are a few historical notes on Theorems \ref{thm:3575} and  \ref{thm:3585}.
The concept of derived Picard group was discovered independently, around 
1997, by R. Rouquier and A. Zimmermann \cite{RoZi} (who had used the notation
$\opn{TrPic}(A)$ for this group) and Yekutieli \cite{Ye4}. The motivations of 
the two teams of authors were very different: Rouquier and Zimmermann were 
interested in invariants of finite dimensional algebras, whereas Yekutieli was 
trying to classify noncommutative dualizing complexes (See Section
\ref{sec:rigid-DC-NC}).
\end{rem}

\begin{rem} \label{rem:4650}
Suppose $\K$ is an algebraically closed field and $A$ is a finite $\K$-ring. 
It turns out that in this case the group $\opn{DPic}_{\K}(A)$ is a {\em locally 
algebraic group}. Namely there is a connected algebraic group 
$\opn{DPic}^0_{\K}(A)$, and $\opn{DPic}_{\K}(A)$ it is a (usually infinite) 
disjoint union of connected algebraic varieties, that are left and right cosets 
of  $\opn{DPic}^0_{\K}(A)$. This same sort of geometric structure can be found 
in the {\em Picard scheme} of a smooth projective algebraic variety $X$, in 
which the identity component is the {\em Jacobian variety} of $X$. But whereas 
the Picard scheme is an abelian group scheme, the group $\opn{DPic}_{\K}(A)$ is 
not 
abelian; and this means that there is a nontrivial geometric action of 
$\opn{DPic}^0_{\K}(A)$, by conjugation, on the connected components. 

The result above is in the paper \cite{Ye4.6}, and it is based on the papers 
\cite{HuSa} by B. Huisgen-Zimmermann and M. Saorin, and \cite{Rou} by Rouquier. 

Later B. Keller \cite{Kel3} used an abstraction of this idea -- essentially  
viewing the derived Picard group as a derived group stack (evaluating it on 
commutative artinian DG rings) -- to prove that the {\em Hochschild cohomology} 
of $A$ is the {\em Lie algebra}, in the derived sense, of $\opn{DPic}_{\K}(A)$. 
This allowed him to prove that the Hochschild cohomology of $A$, with its {\em 
Gerstenhaber Lie bracket}, is a derived invariant of $A$. 
\end{rem}






